\chapter{Modelos de Atención}\label{ch:attention}

Los mecanismos de atención constituyen el complemento natural de las convoluciones geométricas desarrolladas en el capítulo anterior: donde las convoluciones imponen estructura rígida mediante kernels compartidos con pesos independientes del contenido, la atención introduce flexibilidad adaptativa, ponderando dinámicamente las contribuciones de cada elemento según la información que transporta. Este capítulo desarrolla los fundamentos matemáticos de los mecanismos de atención, desde su formulación algebraica clásica hasta las extensiones geométricas que permiten operar sobre dominios no euclidianos preservando simetrías.

Partimos de los mecanismos de atención tradicionales, formalizados como operadores integrales sobre espacios de Hilbert, para después incorporar progresivamente estructura geométrica: primero la equivarianza respecto a grupos de simetrías, después la operación sobre grafos mediante redes de atención y análisis espectral, la extensión a variedades riemannianas con transporte paralelo en espacios hiperbólicos y esféricos, la formulación en álgebras de Clifford que unifica escalares, vectores y multivectores en un marco algebraico coherente, y finalmente la atención sobre complejos simpliciales y estructuras cohomológicas. Cada extensión incorpora estructura geométrica más rica mientras preserva la capacidad de adaptación dinámica que distingue a la atención de la convolución.

Esta progresión de lo algebraico a lo geométrico conecta directamente con los capítulos precedentes: el \autoref{ch:cnn} estableció que la atención generaliza la convolución como caso límite cuando el kernel gaussiano relaja su soporte local, y el \autoref{ch:gdl} identificó los \glspl{Transformer} como instancias del marco \gls{MPNN} sobre grafos completos adaptativos (\autoref{prop:attention_mpnn}). Este capítulo desarrolla en profundidad ambas conexiones.

\section{Mecanismos de Atención Tradicionales}\label{sec:atencion_tradicional}

\subsection{Introducción e Historia}\label{subsec:historia_atencion}

Los mecanismos de atención emergen históricamente de la necesidad de superar las limitaciones inherentes de los modelos secuenciales en el procesamiento de lenguaje natural. Los primeros trabajos de \citet{bahdanau2014neural} y \citet{luong2015effective} introdujeron mecanismos de atención en sistemas de traducción automática basados en redes neuronales recurrentes, permitiendo al modelo ``atender'' a diferentes partes de la secuencia de entrada durante la generación.

El punto de inflexión llegó con el trabajo seminal \textit{"Attention is All You Need"} de \citet{vaswani2017attentionneed}, que demostró que los mecanismos de atención por sí solos, sin recurrencia ni convolución, podían lograr rendimiento de estado del arte en tareas de traducción automática. Esta arquitectura \gls{Transformer} revolucionó no solo el procesamiento de lenguaje natural, sino que se convirtió en el fundamento de modelos de visión por computador, procesamiento de grafos y, eventualmente, geometric deep learning.

La evolución histórica puede resumirse en tres fases principales:

\begin{enumerate}
    \item \textbf{Atención como complemento} (2014-2016): Mecanismos de atención añadidos a \glspl{RNN} para mejorar el manejo de dependencias a largo alcance
    \item \textbf{Atención como arquitectura principal} (2017-2019): \glspl{Transformer} puros reemplazando arquitecturas recurrentes
    \item \textbf{Atención geométrica} (2020-presente): Extensión a datos no euclidianos con preservación de simetrías
\end{enumerate}

\subsection{Fundamentos Algebraicos}\label{subsec:fundamentos_algebraicos_atencion}

Los mecanismos de atención se fundamentan en conceptos de álgebra lineal y geometría de espacios vectoriales. En su esencia, la atención opera sobre tres espacios vectoriales interrelacionados que codifican diferentes aspectos de la información.

\subsubsection{Espacios de Características}\label{subsubsec:espacios_caracteristicas}

Consideremos un espacio de características $\mathcal{H} = \mathbb{R}^d$ donde residen las representaciones de entrada. Para un conjunto de $n$ elementos de entrada, tenemos:

\begin{equation}
    \mathbf{X} = [\mathbf{x}_1, \mathbf{x}_2, \ldots, \mathbf{x}_n]^T \in \mathbb{R}^{n \times d}
    \label{eq:input_matrix}
\end{equation}

donde cada $\mathbf{x}_i \in \mathcal{H}$ representa un elemento de la secuencia de entrada.

\subsubsection{Transformaciones Lineales Fundamentales}

El mecanismo de atención define tres transformaciones lineales fundamentales:

\begin{align}
    \mathbf{Q} &= \mathbf{X}\mathbf{W}_Q \in \mathbb{R}^{n \times d_k} \quad \text{(Consultas)} \\
    \mathbf{K} &= \mathbf{X}\mathbf{W}_K \in \mathbb{R}^{n \times d_k} \quad \text{(Claves)} \\
    \mathbf{V} &= \mathbf{X}\mathbf{W}_V \in \mathbb{R}^{n \times d_v} \quad \text{(Valores)}
    \label{eq:linear_projections}
\end{align}

donde las matrices de peso $\mathbf{W}_Q, \mathbf{W}_K \in \mathbb{R}^{d \times d_k}$ y $\mathbf{W}_V \in \mathbb{R}^{d \times d_v}$ son parámetros entrenables que proyectan el espacio de entrada a subespacios especializados. 

Las dimensiones se relacionan como sigue: si $\mathbf{X} \in \mathbb{R}^{n \times d}$ representa $n$ elementos con $d$ características cada uno, entonces las proyecciones resultan en consultas y claves de dimensión $d_k$ y valores de dimensión $d_v$. En la implementación multi-cabeza estándar, se utiliza $d_k = d_v = d/h$ donde $h$ es el número de cabezas de atención.

\begin{definition}[Producto Punto como Medida de Similaridad]\label{def:dot_product_similarity}
    La función de compatibilidad fundamental en atención se define mediante el producto punto:
    \begin{equation}
        \text{sim}(\mathbf{q}_i, \mathbf{k}_j) = \mathbf{q}_i^T \mathbf{k}_j = \|\mathbf{q}_i\| \|\mathbf{k}_j\| \cos(\theta_{ij})
        \label{eq:dot_product_similarity}
    \end{equation}
    donde $\theta_{ij}$ es el ángulo entre los vectores $\mathbf{q}_i$ y $\mathbf{k}_j$.
\end{definition}

Esta interpretación geométrica revela que la atención mide la similaridad coseno normalizada entre representaciones en el espacio de características.

\subsubsection{Geometría de los Espacios Q, K, V}

La estructura geométrica subyacente de los espacios de consultas, claves y valores proporciona perspectivas profundas sobre el funcionamiento de la atención.

\begin{definition}[Matriz de Similaridad en el Espacio de Atención]\label{def:gram_metric_attention}
    La \textit{matriz de similaridad} (o matriz de productos internos cruzados) asociada a las consultas y claves se define como:
    \begin{equation}
        \mathbf{S} = \mathbf{Q}\mathbf{K}^T, \quad S_{ik} = \langle \mathbf{q}_i, \mathbf{k}_k \rangle \in \mathbb{R}^{n \times n}
        \label{eq:gram_matrix_attention}
    \end{equation}
    Esta matriz codifica todas las relaciones geométricas entre consultas y claves. Cuando $\mathbf{Q} = \mathbf{K}$ (auto-atención pura), $\mathbf{S}$ coincide con la matriz de Gram de las claves.
\end{definition}

\begin{proposition}[Propiedades Geométricas de la Matriz de Similaridad]\label{prop:geometric_properties_gram}
    La matriz $\mathbf{S} = \mathbf{Q}\mathbf{K}^T$ satisface:
    \begin{enumerate}
        \item \textbf{Descomposición en valores singulares}: $\mathbf{S} = \mathbf{U}\mathbf{\Sigma}\mathbf{V}^T$ (SVD), donde $\mathbf{U}, \mathbf{V}$ capturan las direcciones principales de consultas y claves respectivamente y $\mathbf{\Sigma}$ contiene los valores singulares.
        \item \textbf{Norma de Frobenius}: $\|\mathbf{S}\|_F^2 = \text{tr}(\mathbf{Q}\mathbf{K}^T\mathbf{K}\mathbf{Q}^T) = \|\mathbf{K}\mathbf{Q}^T\|_F^2$
        \item \textbf{Interpretación angular}: $S_{ij} = \|\mathbf{q}_i\|\|\mathbf{k}_j\|\cos(\theta_{ij})$, donde $\theta_{ij}$ es el ángulo entre $\mathbf{q}_i$ y $\mathbf{k}_j$.
    \end{enumerate}
\end{proposition}

\begin{proof}
    La propiedad (1) es la descomposición en valores singulares estándar, aplicable a toda matriz real. La propiedad (2) se sigue de $\|\mathbf{S}\|_F^2 = \text{tr}(\mathbf{S}\mathbf{S}^T) = \text{tr}(\mathbf{Q}\mathbf{K}^T(\mathbf{Q}\mathbf{K}^T)^T) = \text{tr}(\mathbf{Q}\mathbf{K}^T\mathbf{K}\mathbf{Q}^T)$, y por la propiedad cíclica de la traza, $\text{tr}(\mathbf{Q}\mathbf{K}^T\mathbf{K}\mathbf{Q}^T) = \text{tr}(\mathbf{K}\mathbf{Q}^T\mathbf{Q}\mathbf{K}^T) = \|\mathbf{K}\mathbf{Q}^T\|_F^2$. La propiedad (3) es consecuencia inmediata de la definición del producto interno en espacios euclidianos.
\end{proof}

\begin{theorem}[Teorema de Proyección Geométrica en Atención]\label{theorem:geometric_projection_attention}
    El mecanismo de atención puede interpretarse como una proyección geométrica del espacio de valores $\mathcal{V}$ sobre el subespacio generado por las consultas $\text{span}(\mathbf{Q})$, mediada por la métrica inducida por las claves.
    
    Formalmente, para cada consulta $\mathbf{q}_i$:
    \begin{equation}
        \mathbf{y}_i = \sum_{j=1}^n \alpha_{ij} \mathbf{v}_j = \mathbf{V}^T \boldsymbol{\alpha}_i
        \label{eq:geometric_projection_formula}
    \end{equation}
    donde $\boldsymbol{\alpha}_i = \text{softmax}(\mathbf{G}_{i,:}/\sqrt{d_k})$ define una distribución de probabilidad sobre el índice de valores.
\end{theorem}

\begin{proof}
    La demostración procede reconociendo que los pesos de atención $\alpha_{ij}$ forman una combinación convexa (probabilística) sobre los valores. Esta combinación está determinada por la proximidad geométrica entre $\mathbf{q}_i$ y cada clave $\mathbf{k}_j$, medida por el producto interno escalado y normalizado por softmax.
    
    La proyección resultante preserva la estructura del espacio de valores mientras enfatiza direcciones alineadas con la consulta según la métrica definida por las claves.
\end{proof}

\begin{figure}[ht]
    \centering
    \includegraphics[width=0.95\textwidth]{figures/qkv_geometric_projection.pdf}
    \caption{Interpretación geométrica del mecanismo de atención. \textbf{Izquierda}: en el espacio de queries y keys, la similitud entre $\mathbf{q}_i$ y cada $\mathbf{k}_j$ se mide por el ángulo $\theta_{ij}$ (producto interno escalado). \textbf{Derecha}: los valores $\mathbf{v}_j$ se ponderan por los pesos de atención $\alpha_{ij}$, y la salida $\mathbf{y}_i$ es su combinación convexa.}
    \label{fig:qkv_geometric_projection}
\end{figure}

\subsubsection{Análisis Multiresolución de Atención}

El análisis de la atención desde una perspectiva de análisis armónico revela estructura jerárquica en múltiples escalas.

\begin{definition}[Espectro de Atención]\label{def:attention_spectrum}
    Para una matriz de atención $\mathbf{A} = \text{softmax}(\mathbf{Q}\mathbf{K}^T/\sqrt{d_k})$, definimos su \textit{espectro} como el conjunto de eigenvalores:
    \begin{equation}
        \text{Spec}(\mathbf{A}) = \{\lambda_1, \lambda_2, \ldots, \lambda_n : \mathbf{A}\mathbf{v}_i = \lambda_i \mathbf{v}_i\}
        \label{eq:attention_spectrum}
    \end{equation}
    Los eigenvectores $\{\mathbf{v}_i\}$ revelan los patrones principales de atención.
\end{definition}

\begin{example}[Cálculo Explícito: Atención $3 \times 3$]\label{ex:explicit_3x3_attention}
    Consideremos un ejemplo concreto con $n=3$ elementos y $d_k=2$:
    
    \begin{equation}
        \mathbf{Q} = \begin{bmatrix} 1 & 0 \\ 0 & 1 \\ 1 & 1 \end{bmatrix}, \quad
        \mathbf{K} = \begin{bmatrix} 1 & 0 \\ 1 & 1 \\ 0 & 1 \end{bmatrix}, \quad
        \mathbf{V} = \begin{bmatrix} 2 & 0 \\ 0 & 2 \\ 1 & 1 \end{bmatrix}
    \end{equation}
    
    \textbf{Paso 1: Matriz de compatibilidad}
    \begin{equation}
        \mathbf{S} = \mathbf{Q}\mathbf{K}^T = \begin{bmatrix} 1 & 1 & 0 \\ 0 & 1 & 1 \\ 1 & 2 & 1 \end{bmatrix}
    \end{equation}

    \textbf{Paso 2: Matriz de atención (sin escalado para simplificar)}
    \begin{equation}
        \mathbf{A} = \text{softmax}(\mathbf{S}) \approx \begin{bmatrix} 0.42 & 0.42 & 0.16 \\ 0.16 & 0.42 & 0.42 \\ 0.21 & 0.58 & 0.21 \end{bmatrix}
    \end{equation}

    \textbf{Paso 3: Salida de atención}
    \begin{equation}
        \mathbf{Y} = \mathbf{A}\mathbf{V} \approx \begin{bmatrix} 1.00 & 1.00 \\ 0.74 & 1.26 \\ 0.63 & 1.37 \end{bmatrix}
    \end{equation}
    
    Este ejemplo ilustra cómo diferentes patrones de consulta-clave producen diferentes combinaciones convexas de valores.
\end{example}

\subsection{Mecanismo de Atención Básico}\label{subsec:atencion_basico}

\subsubsection{Ecuación Fundamental de Atención}

El mecanismo de atención de producto punto escalado se define formalmente como:

\begin{equation}
    \text{Attention}(\mathbf{Q}, \mathbf{K}, \mathbf{V}) = \text{softmax}\left(\frac{\mathbf{Q}\mathbf{K}^T}{\sqrt{d_k}}\right)\mathbf{V}
    \label{eq:scaled_dot_product_attention}
\end{equation}

La matriz de scores de atención $\mathbf{S} = \mathbf{Q}\mathbf{K}^T \in \mathbb{R}^{n \times n}$ codifica la compatibilidad entre cada par consulta-clave, mientras que el factor de escala $\frac{1}{\sqrt{d_k}}$ previene que los valores del producto punto crezcan excesivamente con la dimensión.

\subsubsection{Interpretación Geométrica}

La operación de atención puede interpretarse como una interpolación convexa en el espacio de valores:

\begin{equation}
    \mathbf{y}_i = \sum_{j=1}^{n} \alpha_{ij} \mathbf{v}_j, \quad \sum_{j=1}^{n} \alpha_{ij} = 1, \quad \alpha_{ij} \geq 0
    \label{eq:convex_combination}
\end{equation}

donde los pesos de atención $\alpha_{ij}$ forman una distribución de probabilidad sobre los valores.

\begin{definition}[Función de Atención Generalizada]\label{def:generalized_attention}
    Una \textit{función de atención generalizada} es una función $f: \mathcal{X}^n \times \mathcal{Y}^n \times \mathcal{Z}^n \to \mathcal{W}^n$ donde:
    \begin{align}
        f(\mathbf{Q}, \mathbf{K}, \mathbf{V}) &= \sum_{j=1}^{n} \alpha_{ij} \mathbf{v}_j \\
        \alpha_{ij} &= \frac{\exp(e_{ij})}{\sum_{k=1}^{n} \exp(e_{ik})} \\
        e_{ij} &= a(\mathbf{q}_i, \mathbf{k}_j)
    \end{align}
    donde $a: \mathcal{X} \times \mathcal{Y} \to \mathbb{R}$ es una función de compatibilidad.
\end{definition}

\subsubsection{Auto-atención y Dependencias Globales}

La auto-atención constituye el caso especial donde las consultas, claves y valores provienen de la misma entrada:

\begin{equation}
    \mathbf{Q} = \mathbf{K} = \mathbf{V} = \mathbf{X}
    \label{eq:self_attention_case}
\end{equation}

Esta formulación permite que cada elemento de la secuencia atienda a todos los demás elementos, capturando dependencias globales en una sola operación.

\begin{theorem}[Universalidad de la Auto-atención]\label{theorem:self_attention_universality}
    Sea $K \subseteq \mathbb{R}^{n \times d}$ un conjunto compacto. Para cualquier función continua 
    $f: K \to \mathbb{R}^{n \times d}$ que sea permutación-equivariante (i.e., $f(P\mathbf{X}) = Pf(\mathbf{X})$ 
    para cualquier matriz de permutación $P$), existe una red de auto-atención que puede aproximar 
    $f$ uniformemente en $K$ con precisión arbitraria.
    
    Este resultado extiende los teoremas de aproximación universal para funciones simétricas
    \citep{zaheer2017deepsets} al contexto de mecanismos de atención. La universalidad plena requiere la arquitectura Transformer completa (auto-atención combinada con capas feed-forward y conexiones residuales), no únicamente capas de auto-atención aisladas; véase la demostración rigurosa en \citet{yun2020transformers}.
\end{theorem}

\begin{proof}
    La demostración procede en tres pasos principales:
    
    \textbf{Paso 1: Caracterización de funciones permutación-equivariantes.}
    Sea $f: K \to \mathbb{R}^{n \times d}$ una función continua y permutación-equivariante, donde $K \subseteq \mathbb{R}^{n \times d}$ es compacto. Por el teorema de representación para funciones simétricas \citep{zaheer2017deepsets}, cualquier función permutación-equivariante puede expresarse como:
    
    \begin{equation}
        f(\mathbf{X})_i = \phi\left(\mathbf{x}_i, \sum_{j=1}^n \rho(\mathbf{x}_j)\right)
        \label{eq:symmetric_representation}
    \end{equation}
    
    donde $\phi: \mathbb{R}^d \times \mathbb{R}^d \to \mathbb{R}^d$ y $\rho: \mathbb{R}^d \to \mathbb{R}^d$ son funciones continuas apropiadas.
    
    \textbf{Paso 2: Aproximación mediante interacciones por pares.}
    Por el teorema de aproximación universal para redes neuronales, para cualquier $\epsilon > 0$, existen funciones $g_1, g_2, g_3: \mathbb{R}^d \to \mathbb{R}^{d'}$ y $h: \mathbb{R}^{d'} \to \mathbb{R}^d$ tales que:
    
    \begin{equation}
        \left\|f(\mathbf{X})_i - \sum_{j=1}^n \text{softmax}(g_1(\mathbf{x}_i)^T g_2(\mathbf{x}_j)) \cdot h(g_3(\mathbf{x}_j))\right\| < \epsilon
        \label{eq:pairwise_approximation}
    \end{equation}
    
    uniformemente para $\mathbf{X} \in K$.
    
    \textbf{Paso 3: Implementación exacta mediante auto-atención.}
    Definimos las matrices de proyección como:
    \begin{align}
        \mathbf{W}_Q &= \begin{bmatrix} \mathbf{W}_{g_1} \\ \mathbf{0} \end{bmatrix} \in \mathbb{R}^{d \times d_k} \\
        \mathbf{W}_K &= \begin{bmatrix} \mathbf{W}_{g_2} \\ \mathbf{0} \end{bmatrix} \in \mathbb{R}^{d \times d_k} \\
        \mathbf{W}_V &= \begin{bmatrix} \mathbf{W}_{g_3} \\ \mathbf{0} \end{bmatrix} \in \mathbb{R}^{d \times d_v}
    \end{align}
    
    donde $\mathbf{W}_{g_1}, \mathbf{W}_{g_2}, \mathbf{W}_{g_3}$ implementan las funciones $g_1, g_2, g_3$ respectivamente mediante capas lineales.
    
    Con esta construcción, el mecanismo de auto-atención produce exactamente:
    \begin{equation}
        \text{SelfAttention}(\mathbf{X})_i = \sum_{j=1}^n \text{softmax}\left(\frac{g_1(\mathbf{x}_i)^T g_2(\mathbf{x}_j)}{\sqrt{d_k}}\right) h(g_3(\mathbf{x}_j))
        \label{eq:self_attention_implementation}
    \end{equation}
    
    El factor de escala $\sqrt{d_k}$ puede absorberse en las funciones $g_1$ y $g_2$ sin pérdida de generalidad.
    
    Por composición de funciones universalmente aproximables y la densidad de funciones continuas en espacios compactos, concluimos que las redes de auto-atención pueden aproximar uniformemente cualquier función continua permutación-equivariante en conjuntos compactos.
\end{proof}

\subsection{Multi-Head Attention}

\subsubsection{Descomposición en Subespacios Paralelos}

El mecanismo Multi-Head Attention extiende la atención básica mediante el procesamiento paralelo en múltiples subespacios de representación:

\begin{equation}
    \text{MultiHead}(\mathbf{Q}, \mathbf{K}, \mathbf{V}) = \text{Concat}(\text{head}_1, \ldots, \text{head}_h)\mathbf{W}_O
    \label{eq:multihead_attention}
\end{equation}

donde cada cabeza de atención se define como:

\begin{equation}
    \text{head}_i = \text{Attention}(\mathbf{Q}\mathbf{W}_Q^i, \mathbf{K}\mathbf{W}_K^i, \mathbf{V}\mathbf{W}_V^i)
    \label{eq:attention_head}
\end{equation}

con matrices de proyección $\mathbf{W}_Q^i, \mathbf{W}_K^i \in \mathbb{R}^{d \times d_k}$, $\mathbf{W}_V^i \in \mathbb{R}^{d \times d_v}$ para $i = 1, \ldots, h$.

\subsubsection{Ventajas de la Descomposición}

La arquitectura multi-cabeza proporciona varios beneficios fundamentales:

\begin{enumerate}
    \item \textbf{Especialización de representaciones}: Cada cabeza puede especializarse en capturar diferentes tipos de relaciones
    \item \textbf{Incremento de capacidad}: El modelo puede atender simultáneamente a múltiples posiciones y subespacios
    \item \textbf{Robustez}: La redundancia parcial entre cabezas mejora la estabilidad del entrenamiento
\end{enumerate}

\begin{figure}[ht]
\centering
\includegraphics[width=0.85\textwidth]{figures/multi_head_attention.pdf}
\caption{Diagrama de Multi-Head Attention. La entrada se proyecta en $h$ subespacios paralelos, cada uno con sus propias matrices $\mathbf{W}^Q_i$, $\mathbf{W}^K_i$, $\mathbf{W}^V_i$. Cada cabeza computa atención independientemente y los resultados se concatenan y proyectan mediante $\mathbf{W}^O$.}
\label{fig:multi_head_attention}
\end{figure}

\begin{proposition}[Capacidad Representacional Multi-Cabeza para Scores Pre-Softmax]\label{theorem:multihead_representational_capacity}
    El mecanismo Multi-Head Attention con $h$ cabezas puede representar cualquier patrón de scores pre-softmax que admita descomposición como suma de $h$ funciones de rango bajo.

    Formalmente, si $\mathbf{S}^* = \sum_{i=1}^h \mathbf{U}_i \mathbf{\Sigma}_i \mathbf{V}_i^T$ donde cada $\mathbf{U}_i \mathbf{\Sigma}_i \mathbf{V}_i^T$ tiene rango $r_i \leq d_k$, entonces existe una configuración Multi-Head con $h$ cabezas cuya matriz de scores pre-softmax es exactamente $\mathbf{S}^*$.
\end{proposition}

\begin{proof}
    La demostración procede por construcción directa. Para cada término $\mathbf{U}_i \mathbf{\Sigma}_i \mathbf{V}_i^T$ en la descomposición, definimos la cabeza $i$-ésima como:
    \begin{align}
        \mathbf{W}_Q^i &= \mathbf{U}_i \mathbf{\Sigma}_i^{1/2} \\
        \mathbf{W}_K^i &= \mathbf{V}_i \mathbf{\Sigma}_i^{1/2} \\
        \mathbf{W}_V^i &= \mathbf{I}
    \end{align}

    Con esta elección, la matriz de scores pre-softmax de la cabeza $i$ es $\mathbf{Q}_i \mathbf{K}_i^T = \mathbf{X}\mathbf{U}_i \mathbf{\Sigma}_i \mathbf{V}_i^T \mathbf{X}^T$, y la suma sobre cabezas reproduce la estructura deseada $\mathbf{S}^*$.
\end{proof}

\begin{remark}
La construcción anterior opera a nivel de scores pre-softmax. En la práctica, cada cabeza aplica softmax a sus scores antes de la agregación con los valores, lo que introduce la restricción de que los pesos de atención efectivos son no negativos y suman 1 por fila. Esta no linealidad impide que el mecanismo multi-cabeza represente de forma exacta funciones de atención arbitrarias a nivel de la salida.
\end{remark}

\subsubsection{Análisis de Ortogonalidad entre Cabezas}

\begin{definition}[Ortogonalidad Funcional entre Cabezas]\label{def:functional_orthogonality_heads}
    Dos cabezas de atención $\text{head}_i$ y $\text{head}_j$ son \textit{funcionalmente ortogonales} si sus matrices de proyección satisfacen:
    \begin{equation}
        \langle \mathbf{W}_Q^i, \mathbf{W}_Q^j \rangle_F = \text{tr}((\mathbf{W}_Q^i)^T \mathbf{W}_Q^j) = 0
        \label{eq:head_orthogonality}
    \end{equation}
    donde $\langle \cdot, \cdot \rangle_F$ denota el producto interno de Frobenius.
\end{definition}

\begin{proposition}[Beneficios de Ortogonalidad]\label{prop:orthogonality_benefits}
    Si las cabezas de atención son mutuamente ortogonales en el sentido de la \autoref{def:functional_orthogonality_heads}, entonces:
    \begin{enumerate}
        \item La información capturada por cada cabeza es complementaria (no redundante)
        \item La norma de Frobenius del Multi-Head es la suma de las normas individuales: $\|\text{MultiHead}\|_F^2 = \sum_{i=1}^h \|\text{head}_i\|_F^2$
        \item El gradiente durante backpropagation se descompone limpiamente entre cabezas
    \end{enumerate}
\end{proposition}

\begin{proof}
\leavevmode
\begin{enumerate}
    \item \textbf{Complementariedad.} Si $\langle \mathbf{W}_Q^i, \mathbf{W}_Q^j \rangle_F = 0$ para $i \neq j$, los subespacios de queries son ortogonales. Luego cada cabeza proyecta sobre un subespacio distinto de $\mathbb{R}^d$, capturando información linealmente independiente.

    \item \textbf{Descomposición de norma.} La salida del Multi-Head es $\text{MultiHead} = \sum_{i=1}^h \text{head}_i \mathbf{W}_O^i$. Por la ortogonalidad de las cabezas y la identidad de Parseval para subespacios ortogonales:
    \[
    \|\text{MultiHead}\|_F^2 = \left\|\sum_{i=1}^h \text{head}_i \mathbf{W}_O^i\right\|_F^2 = \sum_{i=1}^h \|\text{head}_i \mathbf{W}_O^i\|_F^2 = \sum_{i=1}^h \|\text{head}_i\|_F^2
    \]
    donde la última igualdad asume $\mathbf{W}_O^i$ ortonormales.

    \item \textbf{Gradientes.} Por la ortogonalidad, $\frac{\partial \mathcal{L}}{\partial \mathbf{W}_Q^i}$ no tiene componentes en la dirección de los subespacios de $\text{head}_j$ ($j \neq i$), evitando interferencia entre las actualizaciones de distintas cabezas. \qedhere
\end{enumerate}
\end{proof}

\subsubsection{Análisis de Complejidad}

La complejidad computacional del Multi-Head Attention es:

\begin{align}
    \text{Tiempo} &: O(n^2 d + ndhd_k) \\
    \text{Espacio} &: O(n^2 + ndh)
\end{align}

donde el término $n^2$ domina para secuencias largas, motivando posteriores optimizaciones como la atención dispersa y lineal.

\subsection{Formalización como Operador Integral}\label{subsec:attention_integral_operator}

El mecanismo de auto-atención puede formalizarse rigurosamente como un operador integral actuando sobre funciones definidas en variedades, proporcionando una perspectiva matemática profunda que conecta con la teoría clásica de operadores.

\subsubsection{Formulación Continua}

Sea $\mathcal{M}$ una variedad riemanniana compacta equipada con una medida $\mu$. Para una función $f: \mathcal{M} \to \mathbb{R}^d$, el operador de atención se define como:

\begin{equation}
    \mathcal{T}f(x) = \int_{\mathcal{M}} K(x, y) f(y) \, d\mu(y)
    \label{eq:attention_integral_operator}
\end{equation}

donde el \textbf{núcleo de atención} $K: \mathcal{M} \times \mathcal{M} \to \mathbb{R}^{d \times d}$ satisface:

\begin{equation}
    K(x, y) = \sigma\left(\frac{\langle Q(x), K(y) \rangle}{\sqrt{d_k}}\right) V(y)^T
    \label{eq:attention_kernel}
\end{equation}

con $Q, K, V: \mathcal{M} \to \mathbb{R}^{d_k}, \mathbb{R}^{d_k}, \mathbb{R}^{d_v}$ respectivamente, y $\sigma$ la función softmax generalizada.

\subsubsection{Propiedades del Operador}

\begin{proposition}[Compacidad del Operador de Atención]\label{prop:attention_compactness}
    Si el núcleo $K(x,y)$ es continuo en $\mathcal{M} \times \mathcal{M}$ y $\mathcal{M}$ es compacta, entonces el operador $\mathcal{T}: L^2(\mathcal{M}) \to L^2(\mathcal{M})$ es compacto.
\end{proposition}

\begin{proof}
    Demostraremos que $\mathcal{T}: L^2(\mathcal{M}) \to L^2(\mathcal{M})$ es compacto aplicando el criterio de Arzelà-Ascoli. Para esto, mostraremos que la imagen de la bola unitaria $B = \{f \in L^2(\mathcal{M}) : \|f\|_{L^2} \leq 1\}$ bajo $\mathcal{T}$ es relativamente compacta en $C(\mathcal{M})$ (y por tanto en $L^2(\mathcal{M})$).
    
    \textbf{Paso 1: Acotamiento uniforme.}
    Para $f \in B$ y $x \in \mathcal{M}$, tenemos:
    \begin{align}
        \|\mathcal{T}f(x)\| &= \left\|\int_{\mathcal{M}} K(x, y) f(y) \, d\mu(y)\right\| \\
        &\leq \int_{\mathcal{M}} \|K(x, y)\| \|f(y)\| \, d\mu(y) \\
        &\leq \sup_{x,y \in \mathcal{M}} \|K(x, y)\| \cdot \|f\|_{L^2} \cdot \sqrt{\mu(\mathcal{M})}
    \end{align}
    
    Como $K$ es continuo en el conjunto compacto $\mathcal{M} \times \mathcal{M}$, la función $\|K(x,y)\|$ alcanza su máximo. Por tanto, $\{\mathcal{T}f : f \in B\}$ está uniformemente acotada.
    
    \textbf{Paso 2: Equicontinuidad.}
    Sean $x_1, x_2 \in \mathcal{M}$ y $f \in B$. Entonces:
    \begin{align}
        \|\mathcal{T}f(x_1) - \mathcal{T}f(x_2)\| &= \left\|\int_{\mathcal{M}} [K(x_1, y) - K(x_2, y)] f(y) \, d\mu(y)\right\| \\
        &\leq \int_{\mathcal{M}} \|K(x_1, y) - K(x_2, y)\| \|f(y)\| \, d\mu(y) \\
        &\leq \|f\|_{L^2} \sqrt{\mu(\mathcal{M})} \sup_{y \in \mathcal{M}} \|K(x_1, y) - K(x_2, y)\|
    \end{align}
    
    Como $K$ es continuo en $\mathcal{M} \times \mathcal{M}$ compacto, es uniformemente continuo. Por tanto, para cualquier $\epsilon > 0$, existe $\delta > 0$ tal que si $d_{\mathcal{M}}(x_1, x_2) < \delta$, entonces:
    \[\sup_{y \in \mathcal{M}} \|K(x_1, y) - K(x_2, y)\| < \frac{\epsilon}{\sqrt{\mu(\mathcal{M})}}\]
    
    Esto implica $\|\mathcal{T}f(x_1) - \mathcal{T}f(x_2)\| < \epsilon$ para $f \in B$, estableciendo la equicontinuidad.
    
    \textbf{Paso 3: Aplicación del teorema de Arzelà-Ascoli.}
    Por los pasos anteriores, $\{\mathcal{T}f : f \in B\}$ es una familia equicontinua y uniformemente acotada en $C(\mathcal{M})$. Como $\mathcal{M}$ es compacta, el teorema de Arzelà-Ascoli garantiza que esta familia es relativamente compacta en $C(\mathcal{M})$.
    
    Como $C(\mathcal{M})$ se embebe densamente en $L^2(\mathcal{M})$ y la convergencia uniforme implica convergencia en $L^2$, concluimos que $\mathcal{T}$ es un operador compacto.
\end{proof}

\subsubsection{Conexión con Análisis Espectral}

La formulación integral permite analizar la atención mediante teoría espectral, siguiendo el marco de análisis desarrollado por \citet{sander2023attention}. Los \textbf{eigenfunciones} del operador $\mathcal{T}$ capturan los patrones fundamentales de atención:

\begin{equation}
    \mathcal{T}\phi_i = \lambda_i \phi_i, \quad \lambda_1 \geq \lambda_2 \geq \cdots \geq 0
    \label{eq:attention_eigendecomposition}
\end{equation}

Los eigenvalores $\lambda_i$ indican la importancia relativa de cada patrón de atención, mientras que las eigenfunciones $\phi_i$ definen las direcciones principales en el espacio de características.

\begin{figure}[ht]
    \centering
    \includegraphics[width=0.95\textwidth]{figures/attention_spectrum.pdf}
    \caption{Descomposición espectral del operador de atención. \textbf{Izquierda}: los eigenvalores $\lambda_i$ del núcleo de atención exhiben decaimiento rápido ($\lambda_i \sim i^{-\alpha}$), lo que indica un rango efectivo finito $r$ del operador. La región sombreada corresponde a los modos dominantes. \textbf{Derecha}: las eigenfunciones $\phi_i$ asociadas muestran frecuencia creciente, análogamente a la base de Fourier; las barras púrpura indican la magnitud del eigenvalor correspondiente.}
    \label{fig:attention_spectrum}
\end{figure}

\subsubsection{Núcleos de Mercer y Teoría Espectral Avanzada}

El mecanismo de atención $\text{softmax}(QK^T/\sqrt{d})V$ puede interpretarse como un operador integral con un núcleo dependiente de los datos. Cuando este núcleo satisface las condiciones de Mercer (simetría y definición positiva), la potente maquinaria del análisis funcional (en particular el teorema espectral para operadores compactos) proporciona garantías sobre las propiedades del operador de atención, como su convergencia y capacidad de aproximación.

\begin{definition}[Núcleo de Atención Gaussiano]\label{def:gaussian_attention_kernel}
    Una variante común del núcleo de atención es el \textit{núcleo gaussiano}:
    \begin{equation}
        K_{\text{gauss}}(x, y) = \exp\left(-\frac{\|Q(x) - K(y)\|^2}{2\sigma^2}\right)
        \label{eq:gaussian_attention_kernel}
    \end{equation}
    donde $\sigma > 0$ es un parámetro de escala que controla la localidad de la atención.
\end{definition}

\begin{theorem}[Descomposición Espectral del Núcleo de Atención]\label{theorem:spectral_decomposition_attention}
    Para un núcleo de atención $K: \mathcal{X} \times \mathcal{X} \to \mathbb{R}$ que es núcleo de Mercer, existe una descomposición espectral:
    \begin{equation}
        K(x, y) = \sum_{i=1}^{\infty} \lambda_i \phi_i(x) \phi_i(y)
        \label{eq:mercer_decomposition_attention}
    \end{equation}
    donde $\{\lambda_i\}$ son eigenvalores no negativos en orden decreciente y $\{\phi_i\}$ son eigenfunciones ortonormales.
    
    El operador de atención correspondiente admite la representación:
    \begin{equation}
        \mathcal{T}f = \sum_{i=1}^{\infty} \lambda_i \langle f, \phi_i \rangle_{L^2} \phi_i
        \label{eq:spectral_representation_attention}
    \end{equation}
\end{theorem}

\begin{proof}
    La demostración sigue del teorema espectral para operadores compactos auto-adjuntos. Como $K$ es núcleo de Mercer, el operador integral asociado $\mathcal{T}$ es compacto y positivo. El teorema espectral garantiza la existencia de la descomposición espectral, donde los eigenvalores son reales y no negativos.
    
    La convergencia de la serie es en la norma $L^2$, y la representación espectral proporciona una caracterización completa del comportamiento del operador.
\end{proof}

\begin{corollary}[Filtrado Espectral Automático]\label{cor:spectral_filtering}
    El operador de atención actúa como un filtro pasa-bajas natural, donde:
    \begin{enumerate}
        \item Componentes de alta frecuencia (eigenvalores pequeños) son atenuadas
        \item La tasa de decaimiento de eigenvalores determina la suavidad de la función de atención
        \item El rango efectivo del operador está determinado por el número de eigenvalores significativos
    \end{enumerate}
\end{corollary}

\begin{proof}
Consecuencia directa del \autoref{theorem:spectral_decomposition_attention}. Por la representación espectral $\mathcal{T}f = \sum_{i=1}^\infty \lambda_i \langle f, \phi_i \rangle \phi_i$ con eigenvalores en orden decreciente $\lambda_1 \geq \lambda_2 \geq \cdots \geq 0$: (1)~las componentes de alta frecuencia (índice $i$ grande) se multiplican por eigenvalores $\lambda_i$ pequeños, atenuándose naturalmente; (2)~la tasa de decaimiento $\lambda_i \to 0$ controla la suavidad de $\mathcal{T}f$ (decaimiento rápido implica mayor regularidad); (3)~el rango efectivo es $r_\epsilon = |\{i : \lambda_i > \epsilon\}|$, que determina la dimensionalidad intrínseca de la representación.
\end{proof}

\subsubsection{Análisis de Convergencia Espectral}

\begin{theorem}[Convergencia de Aproximaciones de Rango Finito]\label{theorem:finite_rank_convergence}
    Sea $\mathcal{T}_r$ la aproximación de rango $r$ del operador de atención:
    \begin{equation}
        \mathcal{T}_r f = \sum_{i=1}^{r} \lambda_i \langle f, \phi_i \rangle \phi_i
        \label{eq:finite_rank_approximation}
    \end{equation}
    
    Entonces, para cualquier $f \in L^2(\mathcal{M})$:
    \begin{equation}
        \|\mathcal{T}f - \mathcal{T}_r f\|_{L^2} \leq \lambda_{r+1} \|f\|_{L^2}
        \label{eq:approximation_error_bound}
    \end{equation}
\end{theorem}

\begin{proof}
    Por la representación espectral:
    \begin{align}
        \mathcal{T}f - \mathcal{T}_r f &= \sum_{i=r+1}^{\infty} \lambda_i \langle f, \phi_i \rangle \phi_i \\
        \|\mathcal{T}f - \mathcal{T}_r f\|_{L^2}^2 &= \sum_{i=r+1}^{\infty} \lambda_i^2 |\langle f, \phi_i \rangle|^2 \\
        &\leq \lambda_{r+1}^2 \sum_{i=r+1}^{\infty} |\langle f, \phi_i \rangle|^2 \\
        &\leq \lambda_{r+1}^2 \|f\|_{L^2}^2
    \end{align}
    donde la última desigualdad usa la identidad de Parseval.
\end{proof}

\subsubsection{Interpretación Física y Geométrica}

El análisis espectral revela que los mecanismos de atención implementan naturalmente principios de análisis armónico:

\begin{enumerate}
    \item \textbf{Descomposición en modos}: Las eigenfunciones $\phi_i$ representan "modos de atención" fundamentales
    \item \textbf{Jerarquía de importancia}: Los eigenvalores $\lambda_i$ ordenan estos modos por importancia
    \item \textbf{Regularización implícita}: El decaimiento de eigenvalores proporciona regularización automática
    \item \textbf{Interpretación frecuencial}: Modos con eigenvalores grandes corresponden a patrones de "baja frecuencia" en el espacio de características
\end{enumerate}

\textbf{Conexión con Análisis de Fourier.}\label{rem:fourier_connection} Cuando el núcleo de atención es invariante a traslaciones, la descomposición espectral se reduce al análisis clásico de Fourier, donde las eigenfunciones son exponenciales complejas y los eigenvalores están relacionados con la transformada de Fourier del núcleo.

\subsubsection{Aplicaciones Computacionales}

El análisis espectral tiene implicaciones prácticas importantes:

\begin{itemize}
    \item \textbf{Compresión de modelos}: Truncar la expansión espectral reduce el costo computacional
    \item \textbf{Análisis de capacidad}: El rango efectivo indica la capacidad representacional
    \item \textbf{Inicialización}: Los modos espectrales pueden guiar la inicialización de pesos
    \item \textbf{Regularización}: Penalizar eigenvalores grandes previene overfitting
\end{itemize}

\begin{theorem}[Aproximación por Atención Finita]\label{theorem:finite_attention_approximation}
    Sea $\{x_i\}_{i=1}^n$ una discretización $\epsilon$-densa de $\mathcal{M}$. El operador de atención discreto:
    \begin{equation}
        \mathcal{T}_n f(x_i) = \sum_{j=1}^n \alpha_{ij} f(x_j), \quad \alpha_{ij} = \frac{\exp(K(x_i, x_j))}{\sum_{k=1}^n \exp(K(x_i, x_k))}
    \end{equation}
    converge al operador continuo $\mathcal{T}$ en norma de operadores cuando $n \to \infty$ y $\epsilon \to 0$.
\end{theorem}

\begin{proof}[Esquema]
Sea $f \in L^2(\mathcal{M})$ y $\mathcal{T}f(x) = \int_{\mathcal{M}} \alpha(x,y) f(y) \, d\mu(y)$ el operador continuo. La aproximación discreta es $\mathcal{T}_n f(x_i) = \sum_{j=1}^n \alpha_{ij} f(x_j)$, que puede verse como una regla de cuadratura. Para una discretización $\epsilon$-densa, la densidad de los puntos asegura que $|\mu(B_\epsilon(x_i)) - 1/n| \to 0$. El error se descompone en: (1) error de cuadratura $\|\mathcal{T}f - \mathcal{T}_n f\| = O(\epsilon^s)$ para núcleos $C^s$, por convergencia de sumas de Riemann sobre variedades compactas; y (2) error de normalización $|\sum_k \exp(K(x_i,x_k)) - \int \exp(K(x_i,y)) d\mu(y)| \to 0$ por el mismo argumento. Bajo condiciones de regularidad del kernel (Lipschitz en ambos argumentos), ambos errores convergen a cero uniformemente.
\end{proof}

Esta formulación revela que los \glspl{Transformer} estándar pueden verse como aproximaciones discretas de operadores integrales sobre variedades, proporcionando una base teórica sólida para extensiones geométricas.

\subsection{Análisis Funcional de Atención}\label{subsec:functional_analysis_attention}

La teoría de espacios de Hilbert con kernel reproductor (RKHS) proporciona un marco riguroso para entender las propiedades funcionales de los mecanismos de atención.

\subsubsection{Atención como Operador en RKHS}

\begin{definition}[Núcleo de Atención como Núcleo de Mercer]\label{def:attention_mercer_kernel}
    El núcleo de atención $K_{\text{att}}: \mathcal{X} \times \mathcal{X} \to \mathbb{R}$ definido por:
    \begin{equation}
        K_{\text{att}}(x, y) = \exp\left(\frac{\langle Q(x), K(y) \rangle}{\sqrt{d_k}}\right)
        \label{eq:attention_mercer_kernel}
    \end{equation}
    es un \textit{núcleo de Mercer} si es simétrico y positivo definido.
\end{definition}

\begin{theorem}[RKHS Inducido por Atención]\label{theorem:attention_rkhs}
    Si $K_{\text{att}}$ es un núcleo de Mercer, entonces existe un espacio de Hilbert $\mathcal{H}_{\text{att}}$ con kernel reproductor tal que:
    \begin{enumerate}
        \item $K_{\text{att}}(x, \cdot) \in \mathcal{H}_{\text{att}}$ para todo $x \in \mathcal{X}$
        \item $\langle f, K_{\text{att}}(x, \cdot) \rangle_{\mathcal{H}_{\text{att}}} = f(x)$ para todo $f \in \mathcal{H}_{\text{att}}$ (propiedad reproductora)
        \item La atención puede expresarse como: $\text{Attention}(x) = \sum_{i} \alpha_i K_{\text{att}}(x, x_i) V_i$
    \end{enumerate}
\end{theorem}

\begin{proof}
    La demostración sigue la construcción estándar de Moore-Aronszajn:
    
    \textbf{Paso 1}: Definimos el pre-Hilbert $\mathcal{H}_0$ como el span lineal de $\{K_{\text{att}}(x, \cdot) : x \in \mathcal{X}\}$.
    
    \textbf{Paso 2}: El producto interno se define por bilinealidad:
    \[\left\langle \sum_i \alpha_i K_{\text{att}}(x_i, \cdot), \sum_j \beta_j K_{\text{att}}(y_j, \cdot) \right\rangle = \sum_{i,j} \alpha_i \beta_j K_{\text{att}}(x_i, y_j)\]
    
    \textbf{Paso 3}: $\mathcal{H}_{\text{att}}$ es la completación de $\mathcal{H}_0$ bajo esta norma.
    
    La propiedad reproductora se verifica directamente por construcción.
\end{proof}

\subsubsection{Regularización y Estabilidad}

\begin{proposition}[Estabilidad de Lipschitz para Atención]\label{prop:lipschitz_stability_attention}
    Si las funciones $Q, K, V$ son $L$-Lipschitz y los datos satisfacen $\|x\|, \|y\| \leq R$ (hipótesis de acotación esencial, ya que la auto-atención estándar no es Lipschitz en dominios no acotados), entonces el operador de atención satisface:
    \begin{equation}
        \|\text{Attention}(x) - \text{Attention}(y)\| \leq C_L \|x - y\|
        \label{eq:lipschitz_bound_attention}
    \end{equation}
    donde $C_L$ depende de $L$, $R$, $d_k$ y el número de elementos de entrada $n$.
\end{proposition}

\begin{proof}
Sea $\alpha_j(x) = \frac{\exp(\langle Q(x), K(x_j)\rangle / \sqrt{d_k})}{\sum_l \exp(\langle Q(x), K(x_l)\rangle / \sqrt{d_k})}$ los pesos de atención. La función softmax es Lipschitz con constante 1 en la norma $\ell^\infty$. Como $Q$ y $K$ son $L$-Lipschitz y $\|x\| \leq R$, la función $x \mapsto \langle Q(x), K(x_j)\rangle / \sqrt{d_k}$ es Lipschitz con constante $O(LR/\sqrt{d_k})$. Componiendo con softmax y la combinación lineal con $V(x_j)$ (también $L$-Lipschitz y acotada), se obtiene $\|\text{Attention}(x) - \text{Attention}(y)\| \leq C_L \|x - y\|$ donde $C_L = O(nL^2R/\sqrt{d_k})$. Para una derivación detallada, véase \citep{kim2021lipschitz}.
\end{proof}

\begin{theorem}[Regularización Espectral Automática]\label{theorem:spectral_regularization}
    El operador de atención actúa como un regularizador espectral natural, filtrando componentes de alta frecuencia según:
    \begin{equation}
        \|\mathcal{T}f\|_{\dot{H}^s} \leq C_s \|f\|_{\dot{H}^{s-\sigma}}
        \label{eq:spectral_regularization}
    \end{equation}
    donde $\dot{H}^s$ denota el espacio de Sobolev homogéneo y $\sigma > 0$ es el orden de regularización determinado por la suavidad del núcleo de atención.
\end{theorem}

\begin{proof}
Por la representación espectral (\autoref{theorem:spectral_decomposition_attention}), $\mathcal{T}f = \sum_i \lambda_i \langle f, \phi_i\rangle \phi_i$. En la norma de Sobolev $\dot{H}^s$, la componente $i$-ésima tiene peso $\mu_i^s$ donde $\mu_i$ son los eigenvalores del Laplaciano de $\mathcal{M}$. Entonces:
\[
\|\mathcal{T}f\|_{\dot{H}^s}^2 = \sum_i \lambda_i^2 \mu_i^{2s} |\langle f, \phi_i\rangle|^2
\]
Si el núcleo de atención $K$ tiene regularidad de Sobolev $\sigma$ (es decir, $\lambda_i = O(\mu_i^{-\sigma})$ cuando $i \to \infty$), entonces $\lambda_i^2 \mu_i^{2s} \leq C_s^2 \mu_i^{2(s-\sigma)}$, de donde $\|\mathcal{T}f\|_{\dot{H}^s} \leq C_s \|f\|_{\dot{H}^{s-\sigma}}$.
\end{proof}

\begin{remark}
La demostración anterior supone implícitamente que las eigenfunciones $\phi_i$ del operador de atención coinciden (o están estrechamente relacionadas) con las eigenfunciones del Laplaciano de $\mathcal{M}$. Esta suposición se cumple para núcleos isótropos que conmutan con el Laplaciano (por ejemplo, núcleos de Matérn en variedades compactas), pero no es válida para núcleos de atención generales.
\end{remark}

\subsubsection{Convergencia y Aproximación}

\begin{theorem}[Convergencia de Monte Carlo para Atención]\label{theorem:monte_carlo_convergence}
    Sea $\mathcal{T}_n$ el operador de atención discreto con $n$ puntos muestreados i.i.d. según una medida $\mu$. Entonces:
    \begin{equation}
        \mathbb{E}[\|\mathcal{T}_n - \mathcal{T}\|_{L^2 \to L^2}^2] \leq \frac{C}{n}
        \label{eq:monte_carlo_rate}
    \end{equation}
    donde $C$ es una constante que depende de las propiedades del núcleo y la medida.
\end{theorem}

\begin{proof}
Consideramos la versión con normalización $1/n$ (sin softmax) para facilitar el análisis. El operador continuo es $\mathcal{T}f(x) = \int_\mathcal{M} K(x,y) f(y)\, d\mu(y)$ y su aproximación Monte Carlo es $\mathcal{T}_n f(x) = \frac{1}{n}\sum_{j=1}^n K(x, y_j) f(y_j)$ con $y_j \sim \mu$ i.i.d. Para $f$ fijo, $\mathcal{T}_n f(x)$ es la media empírica de $n$ variables aleatorias i.i.d. $Z_j = K(x, y_j)f(y_j)$ con esperanza $\mathcal{T}f(x)$. Para la tasa:
\[
\mathbb{E}[\|\mathcal{T}_n f - \mathcal{T}f\|_{L^2}^2] = \frac{1}{n}\int_\mathcal{M} \text{Var}(K(x,\cdot)f(\cdot))\, d\mu(x) \leq \frac{C}{n}
\]
donde $C = \|K\|_\infty^2 \|f\|_{L^2}^2 \mu(\mathcal{M})$. La extensión a norma de operadores y al caso con normalización softmax (como en el enunciado) requiere argumentos adicionales de cobertura y control de la convergencia del denominador, respectivamente.
\end{proof}

\section{Teoría de Grupos y Equivarianza}

La incorporación de simetrías geométricas en mecanismos de atención constituye uno de los desarrollos más significativos en geometric deep learning, permitiendo que los modelos preserven estructuras fundamentales de los datos.

\subsection{Fundamentos de Teoría de Grupos}

Los fundamentos de teoría de grupos (grupos, acciones de grupo, equivarianza) necesarios para mecanismos de atención equivariante se desarrollan completamente en el \autoref{ch:gdl}, Secciones~\ref{subsec:teoria_grupos_representaciones} y~\ref{subsec:equivarianza_invarianza}. 

Los mecanismos de atención extienden estos principios mediante operaciones adaptativas que preservan simetrías de grupos arbitrarios, superando las limitaciones de equivarianza fija en \glspl{CNN}.

\subsubsection{Grupos Relevantes para Atención Geométrica}

En el contexto de mecanismos de atención, los grupos más importantes incluyen:

\begin{enumerate}
    \item \textbf{$SO(3)$}: Rotaciones 3D, esencial para Point \glspl{Transformer} y datos moleculares
    \item \textbf{$SE(3)$}: Grupo euclidiano especial, fundamental en SE(3)-\glspl{Transformer} 
    \item \textbf{$E(3)$}: Grupo euclidiano completo, incluyendo reflexiones
    \item \textbf{Grupos de Lie generales}: Marco unificado para Lie\glspl{Transformer}
    \item \textbf{Grupos de permutaciones}: Para invarianza a reordenamientos en conjuntos
\end{enumerate}

El \autoref{theorem:equivariant_composition} del \autoref{ch:gdl} asegura que las capas de atención equivariante pueden componerse preservando las propiedades de simetría deseadas.

\subsection{Atención Equivariante}\label{subsec:equivariant_attention}\label{sec:equivariant_attention}

\subsubsection{Definición y Construcción}

Aplicando la definición de equivarianza del \autoref{ch:gdl} (\autoref{def:equivariance_gdl}) a mecanismos de atención, un modelo de atención es $G$-equivariante si preserva las transformaciones de grupo en consultas, claves y valores.

\begin{theorem}[Construcción de Atención Equivariante]
    Para un grupo $G$ actuando en el espacio de características, un mecanismo de atención es $G$-equivariante si y solo si:
    \begin{enumerate}
        \item Las funciones de compatibilidad son $G$-invariantes
        \item Las transformaciones de consulta, clave y valor conmutan con la acción del grupo
    \end{enumerate}
\end{theorem}

\begin{proof}[Esbozo de la demostración]
\textbf{Dirección directa} ($\Leftarrow$): Si las funciones de compatibilidad $a(\mathbf{q}_i, \mathbf{k}_j)$ son $G$-invariantes, entonces $a(\rho(g)\mathbf{q}_i, \rho(g)\mathbf{k}_j) = a(\mathbf{q}_i, \mathbf{k}_j)$ para todo $g \in G$, por lo que los pesos de atención $\alpha_{ij}$ son invariantes. Si además las transformaciones de valor conmutan con $G$ ($W_V \rho(g) = \rho(g) W_V$), la salida $\sum_j \alpha_{ij} W_V \rho(g)\mathbf{x}_j = \rho(g) \sum_j \alpha_{ij} W_V \mathbf{x}_j$, lo que da equivarianza.

\textbf{Dirección conversa} ($\Rightarrow$): Si la atención es $G$-equivariante, los pesos $\alpha_{ij}$ deben ser invariantes bajo la acción simultánea de $G$ en Q y K (pues determinan la combinación lineal), lo que exige invarianza de las funciones de compatibilidad. La equivarianza de la salida con pesos fijos implica que $W_V$ conmuta con $\rho(g)$.
\end{proof}

\begin{lemma}[Condición Suficiente para Equivarianza en Auto-atención]\label{lemma:self_attention_equivariance}
Sea $G$ un grupo actuando en el espacio de entrada $\mathbb{R}^{n \times d}$ mediante permutaciones. 
Si las matrices de proyección satisfacen $\mathbf{W}_Q = \mathbf{W}_K = \mathbf{W}_V = \mathbf{W}$ 
y $\mathbf{W}$ conmuta con la acción del grupo, entonces el mecanismo de auto-atención 
es $G$-equivariante.
\end{lemma}

\begin{proof}
Para cualquier permutación $P \in G$ y entrada $\mathbf{X} \in \mathbb{R}^{n \times d}$, tenemos:
\begin{align}
\text{SelfAttention}(P\mathbf{X}) &= \text{softmax}\left(\frac{(P\mathbf{X}\mathbf{W})(P\mathbf{X}\mathbf{W})^T}{\sqrt{d_k}}\right)(P\mathbf{X}\mathbf{W}) \\
&= \text{softmax}\left(\frac{P(\mathbf{X}\mathbf{W})(\mathbf{X}\mathbf{W})^T P^T}{\sqrt{d_k}}\right)P(\mathbf{X}\mathbf{W}) \\
&= P \cdot \text{SelfAttention}(\mathbf{X})
\end{align}
donde el segundo paso usa la conmutatividad y el tercero la equivarianza de softmax a permutaciones.
\end{proof}

\subsubsection{Representaciones Irreducibles}

La teoría de representaciones irreducibles (\autoref{def:representacion_irreducible}, \autoref{ch:gdl}) proporciona el marco matemático fundamental para diseñar mecanismos de atención equivariantes bajo grupos de Lie compactos. Esta teoría, desarrollada inicialmente por \citet{Peter1927} y extendida al contexto del aprendizaje profundo por \citet{kondor2018generalizationconvolutionalneural}, establece que cualquier representación de un grupo compacto se descompone en bloques fundamentales irreducibles.

Por el teorema de Peter-Weyl (\autoref{thm:peter_weyl}, \autoref{ch:gdl}), toda representación se descompone en suma directa de representaciones irreducibles:
\begin{equation}
V = \bigoplus_{\lambda \in \widehat{G}} n_\lambda \mathcal{H}_\lambda
\label{eq:peter_weyl_decomposition_attention}
\end{equation}
donde $\widehat{G}$ indexa las representaciones irreducibles inequivalentes, $\mathcal{H}_\lambda$ son los espacios de representación, y $n_\lambda$ son las multiplicidades.

\paragraph{Representaciones Irreducibles de SO(3)}

Para el grupo de rotaciones tridimensionales $SO(3)$, las representaciones irreducibles se clasifican por un número entero no negativo $\ell \in \mathbb{N}_0$ llamado \textit{grado} o \textit{momento angular}. Cada representación irreducible tiene dimensión $2\ell + 1$.

\begin{theorem}[Clasificación de Irreps de SO(3)]\label{thm:so3_irreps}
Las representaciones irreducibles de $SO(3)$ están indexadas por $\ell \in \mathbb{N}_0$ y tienen dimensión $d_\ell = 2\ell + 1$. El espacio de representación $\mathcal{H}_\ell$ está generado por los armónicos esféricos $Y_\ell^m$ con $m \in \{-\ell, -\ell+1, \ldots, \ell-1, \ell\}$.
\end{theorem}

\begin{proof}[Esquema]
La clasificación se establece mediante la relación entre $SO(3)$ y $SU(2)$. El álgebra de Lie $\mathfrak{so}(3) \cong \mathfrak{su}(2)$ tiene generadores $\{J_x, J_y, J_z\}$ con relaciones de conmutación $[J_i, J_j] = i\epsilon_{ijk}J_k$. El operador de Casimir $\mathbf{J}^2 = J_x^2 + J_y^2 + J_z^2$ conmuta con todo el álgebra, y por el lema de Schur actúa como escalar en cada irreducible. La diagonalización simultánea de $\mathbf{J}^2$ y $J_z$ produce valores propios $\ell(\ell+1)$ y $m$ respectivamente, con $-\ell \leq m \leq \ell$, dando dimensión $2\ell+1$. Los armónicos esféricos $Y_\ell^m$ realizan estas representaciones como funciones en $S^2$. Véase \citep{brocker1985representations} para una demostración completa.
\end{proof}

\begin{definition}[Armónicos Esféricos]\label{def:spherical_harmonics_attention}
Los \textit{armónicos esféricos} $Y_\ell^m: S^2 \to \mathbb{C}$ de grado $\ell$ y orden $m$ son funciones propias del operador laplaciano en la esfera:
\begin{equation}
Y_\ell^m(\theta, \phi) = \sqrt{\frac{2\ell+1}{4\pi}\frac{(\ell-m)!}{(\ell+m)!}} P_\ell^m(\cos\theta) e^{im\phi}
\label{eq:spherical_harmonics_definition}
\end{equation}
donde $P_\ell^m$ son los polinomios asociados de Legendre. Bajo una rotación $R \in SO(3)$, los armónicos esféricos transforman según:
\begin{equation}
(R \cdot Y_\ell^m)(\mathbf{x}) = Y_\ell^m(R^{-1}\mathbf{x}) = \sum_{m'=-\ell}^{\ell} D_\ell^{m'm}(R) Y_\ell^{m'}(\mathbf{x})
\label{eq:spherical_harmonics_transformation}
\end{equation}
donde $D_\ell^{m'm}(R)$ son las matrices de Wigner.
\end{definition}

\begin{definition}[Matrices de Wigner]\label{def:wigner_matrices}
Las \textit{matrices de Wigner} $D_\ell: SO(3) \to GL(2\ell+1, \mathbb{C})$ proporcionan la representación matricial explícita de las rotaciones en el espacio de armónicos esféricos de grado $\ell$. Para rotaciones parametrizadas por ángulos de Euler $(\alpha, \beta, \gamma)$, tienen la forma:
\begin{equation}
D_\ell^{m'm}(\alpha, \beta, \gamma) = e^{-im'\alpha} d_\ell^{m'm}(\beta) e^{-im\gamma}
\label{eq:wigner_d_matrices}
\end{equation}
donde $d_\ell^{m'm}(\beta)$ son las matrices reducidas de Wigner, que dependen solo del ángulo de rotación $\beta$.
\end{definition}

\begin{figure}[ht]
    \centering
    \includegraphics[width=0.85\textwidth]{figures/spherical_harmonics.pdf}
    \caption{Armónicos esféricos $Y_\ell^m(\theta, \varphi)$ para grados $\ell = 0, 1, 2$. Las regiones azules y rojas indican valores positivos y negativos, respectivamente. El grado $\ell$ controla la frecuencia angular (número de cambios de signo) y el orden $m$ determina la orientación del patrón nodal. Estas funciones forman la base de las representaciones irreducibles de $SO(3)$ de dimensión $2\ell+1$ utilizadas en la descomposición de características equivariantes~\eqref{eq:irrep_decomposition}.}
    \label{fig:spherical_harmonics_attention}
\end{figure}

En el contexto de mecanismos de atención equivariantes, las características en cada capa se descomponen según tipos de representación:

\begin{equation}
    \mathbf{f} = \bigoplus_{\ell=0}^{L_{\max}} \mathbf{f}^{(\ell)}, \quad \mathbf{f}^{(\ell)} \in \mathbb{R}^{n_\ell \times (2\ell+1)}
    \label{eq:irrep_decomposition}
\end{equation}

donde $\mathbf{f}^{(\ell)}$ contiene $n_\ell$ canales de tipo $\ell$, cada uno transformando como un armónico esférico de grado $\ell$ bajo rotaciones.

\paragraph{Coeficientes de Clebsch-Gordan y Productos Tensoriales}

La composición de características de diferentes tipos requiere el producto tensorial de representaciones irreducibles. Los coeficientes de Clebsch-Gordan describen cómo se descompone este producto.

\begin{theorem}[Descomposición del Producto Tensorial]\label{thm:tensor_product_so3}
El producto tensorial de dos representaciones irreducibles de $SO(3)$ de grados $\ell_1$ y $\ell_2$ se descompone como:
\begin{equation}
D^{(\ell_1)} \otimes D^{(\ell_2)} = \bigoplus_{\ell=|\ell_1-\ell_2|}^{\ell_1+\ell_2} D^{(\ell)}
\label{eq:tensor_product_decomposition_attention}
\end{equation}
Esta descomposición es fundamental para construir operaciones equivariantes que combinen características de diferentes tipos. La demostración se obtiene aplicando la regla de adición de momentos angulares, véase \citep{brocker1985representations}.
\end{theorem}

\begin{definition}[Coeficientes de Clebsch-Gordan]\label{def:clebsch_gordan_attention}
Los \textit{coeficientes de Clebsch-Gordan} $C_{l_1 m_1, l_2 m_2}^{l m}$ (véase notación en el archivo \texttt{notation.tex}) proporcionan los coeficientes de expansión del producto de armónicos esféricos:
\begin{equation}
Y_{\ell_1}^{m_1} \otimes Y_{\ell_2}^{m_2} = \sum_{\ell=|\ell_1-\ell_2|}^{\ell_1+\ell_2} \sum_{m=-\ell}^{\ell} C_{l_1 m_1, l_2 m_2}^{l m} Y_{\ell}^{m}
\label{eq:clebsch_gordan_expansion}
\end{equation}
Estos coeficientes satisfacen relaciones de ortogonalidad y pueden calcularse mediante fórmulas recursivas o tablas precalculadas.
\end{definition}

\begin{example}[Casos de Bajo Grado]\label{ex:low_degree_clebsch_gordan}
Para ilustrar concretamente, consideremos algunos casos de bajo grado:

\textbf{Producto $\ell_1 = 0, \ell_2 = 1$} (escalar $\times$ vector):
\begin{equation}
D^{(0)} \otimes D^{(1)} = D^{(1)}
\end{equation}
Este producto preserva el tipo vectorial, simplemente escalando el vector.

\textbf{Producto $\ell_1 = 1, \ell_2 = 1$} (vector $\times$ vector):
\begin{equation}
D^{(1)} \otimes D^{(1)} = D^{(0)} \oplus D^{(1)} \oplus D^{(2)}
\end{equation}
Este producto genera un escalar ($\ell=0$, el producto punto), un pseudovector ($\ell=1$, el producto cruz), y un tensor simétrico traceless ($\ell=2$).
\end{example}

\textbf{Aplicación en Atención Equivariante.} En mecanismos de atención equivariantes como SE(3)-Transformers \citep{fuchs2020se3transformers} y Equiformer \citep{liao2022equiformer}, los productos tensoriales mediante coeficientes de Clebsch-Gordan permiten:
\begin{itemize}
    \item \textbf{Combinar consultas y claves}: $\mathbf{q}_i^{(\ell_1)} \otimes \mathbf{k}_j^{(\ell_2)}$ para calcular scores de atención equivariantes
    \item \textbf{Mezclar información geométrica}: Incorporar direcciones relativas $\hat{\mathbf{r}}_{ij}$ (tipo $\ell=1$) con características de nodos
    \item \textbf{Mantener equivarianza exacta}: Cada operación respeta la estructura de grupo de $SO(3)$ por construcción
\end{itemize}

\begin{example}[Descomposición de Espacio de Características]\label{ex:feature_space_decomposition}
Consideremos un espacio de características con $L_{\max} = 2$:
\begin{align}
\mathbf{f} &= \mathbf{f}^{(0)} \oplus \mathbf{f}^{(1)} \oplus \mathbf{f}^{(2)} \\
&= \underbrace{\mathbb{R}^{n_0 \times 1}}_{\text{escalares}} \oplus \underbrace{\mathbb{R}^{n_1 \times 3}}_{\text{vectores}} \oplus \underbrace{\mathbb{R}^{n_2 \times 5}}_{\text{tensores grado 2}}
\label{eq:feature_space_example}
\end{align}

Una operación de atención equivariante calcula:
\begin{equation}
e_{ij} = \sum_{\ell_q, \ell_k, \ell_{\text{out}}} \text{MLP}_{\ell_q, \ell_k, \ell_{\text{out}}} \left( \mathbf{q}_i^{(\ell_q)} \otimes \mathbf{k}_j^{(\ell_k)} \right)^{(\ell_{\text{out}})}
\label{eq:equivariant_attention_score_irreps}
\end{equation}

donde el producto tensorial se descompone según la Ecuación~\eqref{eq:tensor_product_decomposition_attention} y las MLPs operan independientemente en cada tipo de representación.
\end{example}

\subsection{Arquitecturas Equivariantes}\label{sec:equivariant_transformers}

\subsubsection[SE(3)-Transformers]{SE(3)-\glspl{Transformer}}

Los SE(3)-\glspl{Transformer} de \citet{fuchs2020se3transformers} extienden la atención para ser equivariante bajo rotaciones y traslaciones en 3D.

Para características escalares $f_0$, vectoriales $\mathbf{f}_1$ y tensoriales de orden superior, la atención se construye como:

\begin{align}
    f_0^{\text{out}} &= \sum_j \alpha_{ij} f_{0,j} \\
    \mathbf{f}_1^{\text{out}} &= \sum_j \alpha_{ij} \mathbf{f}_{1,j} \\
    \alpha_{ij} &= \text{softmax}(\text{\gls{MLP}}(f_{0,i}, f_{0,j}, \|\mathbf{r}_{ij}\|, \mathbf{f}_{1,i} \cdot \mathbf{f}_{1,j}))
    \label{eq:se3_transformer_attention}
\end{align}

donde $\mathbf{r}_{ij} = \mathbf{r}_i - \mathbf{r}_j$ es el vector de posición relativa.

\begin{theorem}[SE(3)-Equivarianza de Atención]
    El mecanismo de atención definido en la ecuación~\eqref{eq:se3_transformer_attention} es equivariante bajo el grupo $SE(3)$ de rotaciones y traslaciones euclidianas \citep{fuchs2020se3transformers}.
\end{theorem}

\begin{proof}
Sea $(R, \mathbf{t}) \in SE(3)$ una roto-traslación. Los pesos de atención $\alpha_{ij}$ dependen de: (i) escalares $f_{0,i}, f_{0,j}$ (invariantes por definición), (ii) la norma $\|\mathbf{r}_{ij}\| = \|\mathbf{r}_i - \mathbf{r}_j\|$ que es invariante bajo $SE(3)$: $\|R\mathbf{r}_i + \mathbf{t} - R\mathbf{r}_j - \mathbf{t}\| = \|R(\mathbf{r}_i - \mathbf{r}_j)\| = \|\mathbf{r}_i - \mathbf{r}_j\|$, y (iii) $\mathbf{f}_{1,i} \cdot \mathbf{f}_{1,j}$ que es invariante ya que $\mathbf{f}_1$ transforma como vector tipo-1 bajo $SO(3)$: $(R\mathbf{f}_{1,i}) \cdot (R\mathbf{f}_{1,j}) = \mathbf{f}_{1,i} \cdot \mathbf{f}_{1,j}$. Por tanto, los scores $\alpha_{ij}$ son \emph{invariantes} bajo $SE(3)$.

Para la salida vectorial: $\mathbf{f}_1^{\text{out}} = \sum_j \alpha_{ij} \mathbf{f}_{1,j}$. Bajo $(R, \mathbf{t})$: $\sum_j \alpha_{ij} R\mathbf{f}_{1,j} = R \sum_j \alpha_{ij} \mathbf{f}_{1,j} = R \mathbf{f}_1^{\text{out}}$. Así, la salida tipo-1 es \emph{equivariante}. Las salidas escalares (tipo-0) son trivialmente invariantes.
\end{proof}

\subsubsection{Equiformer: Atención \gls{MLP} Equivariante}

El Equiformer de \citet{liao2022equiformer} supera el producto punto tradicional mediante atención basada en \gls{MLP} equivariante.

En lugar del producto punto $\mathbf{q}_i^T \mathbf{k}_j$, utiliza:

\begin{equation}
    e_{ij} = \text{\gls{MLP}}_{\text{attn}}(\mathbf{q}_i \otimes \mathbf{k}_j, \|\mathbf{r}_{ij}\|)
    \label{eq:equiformer_attention_scores}
\end{equation}

donde $\otimes$ denota el producto tensorial de características irreducibles.

El producto tensorial de representaciones irreducibles se descompone como:

\begin{equation}
    D^{(\ell_1)} \otimes D^{(\ell_2)} = \bigoplus_{\ell=|\ell_1-\ell_2|}^{\ell_1+\ell_2} D^{(\ell)}
    \label{eq:tensor_product_decomposition}
\end{equation}

Esta construcción preserva la equivarianza SO(3) mientras permite interacciones más ricas que el simple producto punto.

\subsubsection{Lie\gls{Transformer}: Framework General}

El Lie\gls{Transformer} de \citet{hutchinson2020lietransformer} proporciona un framework unificado para atención equivariante a grupos de Lie arbitrarios.

El enfoque se basa en ``lifting'' las representaciones al grupo:

\begin{enumerate}
    \item \textbf{Lift}: $\mathbf{x} \mapsto \{(g, \mathbf{x})\}_{g \in G}$
    \item \textbf{Procesar}: Aplicar atención en el espacio del grupo
    \item \textbf{Project}: Reducir de vuelta al espacio original
\end{enumerate}

\begin{equation}
    \text{LieAttn}(\mathbf{x}) = \text{Project}(\text{Attention}(\text{Lift}(\mathbf{x})))
    \label{eq:lie_attention}
\end{equation}

Para un grupo $G$, la representación regular izquierda se define como:

\begin{equation}
    (L_g f)(h) = f(g^{-1}h), \quad g, h \in G, f \in L^2(G)
    \label{eq:left_regular_representation}
\end{equation}

Esta construcción garantiza equivarianza para cualquier grupo de Lie.

\subsubsection{EPT: Modelos Fundacionales Equivariantes}

El Equivariant Pretrained Transformer (EPT) de \citet{jiao2024ept} representa el primer modelo fundacional que unifica el aprendizaje geométrico de moléculas pequeñas y proteínas bajo un marco equivariante.

\begin{definition}[Preentrenamiento Multi-Dominio Equivariante]
    Un modelo fundacional $F_{\text{EPT}}$ es multi-dominio equivariante si para dominios $\mathcal{D}_1, \mathcal{D}_2$ (e.g., moléculas, proteínas) y grupo $G$:
    \begin{equation}
        F_{\text{EPT}}(g \cdot \mathbf{x}) = g \cdot F_{\text{EPT}}(\mathbf{x}), \quad \forall \mathbf{x} \in \mathcal{D}_1 \cup \mathcal{D}_2, g \in G
        \label{eq:multi_domain_equivariance}
    \end{equation}
\end{definition}

La arquitectura EPT incorpora bloques de representación mejorados:

\begin{equation}
    \mathbf{h}_i^{\text{EPT}} = \text{E(3)-TransformerBlock}(\mathbf{h}_i, \{\mathbf{h}_j\}_{j \in \mathcal{N}_{\text{ext}}(i)})
    \label{eq:ept_enhanced_block}
\end{equation}

EPT utiliza múltiples objetivos de preentrenamiento que preservan la geometría:

\begin{align}
    \mathcal{L}_{\text{EPT}} &= \mathcal{L}_{\text{masked}} + \mathcal{L}_{\text{geom}} + \mathcal{L}_{\text{cross}} \\
    \mathcal{L}_{\text{geom}} &= \|\mathbf{f}(g \cdot \mathbf{x}) - g \cdot \mathbf{f}(\mathbf{x})\|^2 \\
    \mathcal{L}_{\text{cross}} &= \text{KL}(\mathbf{p}_{\text{mol}}, \mathbf{p}_{\text{prot}})
    \label{eq:ept_loss}
\end{align}

\begin{proposition}[Transferencia de Conocimiento Geométrico]
    Si un modelo $F$ es preentrenado en un dominio $\mathcal{D}_1$ con equivarianza $G$, entonces la transferencia a un dominio $\mathcal{D}_2$ con la misma simetría $G$ preserva las propiedades de equivarianza \citep{jiao2024ept}.
\end{proposition}

\begin{proof}
Si $F$ es $G$-equivariante, esto significa que sus capas están construidas con operaciones equivariantes (productos tensoriales, convoluciones de grupo, capas lineales equivariantes). La equivarianza es una propiedad \emph{arquitectónica}, no depende de los datos de entrenamiento ni de los valores de los parámetros aprendidos. Al transferir $F$ a $\mathcal{D}_2$ (con o sin fine-tuning), la estructura de las capas no cambia, por lo que $F(g \cdot \mathbf{x}) = g \cdot F(\mathbf{x})$ se mantiene para todo $g \in G$ y $\mathbf{x} \in \mathcal{D}_2$. La condición necesaria es que ambos dominios admitan la misma acción del grupo $G$.
\end{proof}

\subsection{Atención Equivariante en Nubes de Puntos}

Las nubes de puntos 3D presentan desafíos únicos para la atención debido a su naturaleza no estructurada y la necesidad de preservar equivarianza bajo transformaciones rígidas.

\subsubsection{Point \gls{Transformer}: Atención Vectorial}

El Point \gls{Transformer} de \citet{zhao2021point} introduce atención vectorial específicamente diseñada para nubes de puntos:

\begin{equation}
    \mathbf{y}_i = \sum_{j \in \mathcal{N}(i)} \rho(\gamma(\phi(\mathbf{x}_i) - \psi(\mathbf{x}_j) + \delta))(\alpha(\mathbf{x}_j) + \delta)
    \label{eq:point_transformer_attention}
\end{equation}

donde:
\begin{itemize}
    \item $\phi, \psi, \alpha$ son transformaciones lineales
    \item $\delta = \text{\gls{MLP}}(\mathbf{p}_i - \mathbf{p}_j)$ codifica información posicional
    \item $\gamma, \rho$ son \glspl{MLP}
\end{itemize}

La función $\delta$ captura la geometría local:

\begin{equation}
    \delta_{ij} = \text{\gls{MLP}}(\mathbf{p}_i - \mathbf{p}_j, \|\mathbf{p}_i - \mathbf{p}_j\|)
    \label{eq:positional_encoding_3d}
\end{equation}

La formulación del Point \gls{Transformer} puede entenderse como una instancia del marco de atención equivariante desarrollado en la \autoref{subsec:equivariant_attention}. La codificación posicional $\delta_{ij}$ depende exclusivamente de la posición relativa $\mathbf{p}_i - \mathbf{p}_j$, que es invariante bajo traslaciones. Asimismo, si la \gls{MLP} que computa $\delta$ recibe como entrada únicamente la diferencia relativa y su norma, la arquitectura es invariante bajo rotaciones en el cómputo de los pesos de atención, conectando directamente con la condición de invarianza de las funciones de compatibilidad del \autoref{subsec:equivariant_attention}. La atención vectorial (en contraste con la atención escalar estándar) permite que diferentes dimensiones del vector de atención modulen independientemente las características, proporcionando mayor capacidad expresiva que un único peso escalar $\alpha_{ij}$ por par de puntos.

Una ventaja computacional del Point \gls{Transformer} frente a arquitecturas como SE(3)-\glspl{Transformer} reside en evitar el producto tensorial de representaciones irreducibles (Ecuación~\eqref{eq:tensor_product_decomposition_attention}), que tiene coste $O(\ell_{\max}^3)$ para grado máximo $\ell_{\max}$. Al operar directamente sobre diferencias de posiciones y características, el Point \gls{Transformer} logra complejidad $O(|\mathcal{N}(i)| \cdot d)$ por punto, a costa de renunciar a la equivarianza exacta bajo $SO(3)$ para las características de salida (solo las posiciones relativas son invariantes, no las características transformadas).

\section{Teoría de Grafos y Análisis Espectral}

Los mecanismos de atención desarrollados hasta aquí operan sobre datos con estructura regular (secuencias, imágenes, nubes de puntos). Sin embargo, muchos problemas de interés (redes sociales, moléculas, grafos de conocimiento) involucran datos con conectividad irregular. La extensión natural de la atención a estos dominios requiere combinar la teoría de grafos con el análisis espectral, proporcionando fundamentos matemáticos sólidos para los pesos adaptativos en datos estructurados irregulares.

\subsection{Fundamentos de Grafos}

Para los fundamentos de teoría de grafos (definición de grafo, matriz de adyacencia, Laplaciano) necesarios en mecanismos de atención, véase el \autoref{ch:gdl}, \autoref{sec:grafos_discretizacion}. Los conceptos de análisis espectral se desarrollan en el \autoref{ch:cnn}.

\subsubsection{Conexión con Message Passing}

Los mecanismos de atención en grafos generalizan el paradigma de message passing:

\begin{equation}
    \text{Atención} \equiv \text{MessagePassing con pesos adaptativos}
    \label{eq:attention_message_passing_equivalence}
\end{equation}

\subsection{Graph Attention Networks}\label{subsec:graph_attention_networks}\label{sec:graph_attention}

Los \glspl{GAT} representan la intersección natural entre la teoría espectral de grafos desarrollada en la \autoref{sec:analisis_armonico} y los mecanismos de atención adaptativos. Mientras que los métodos espectrales tradicionales requieren descomposición propia costosa, los \glspl{GAT} logran poder expresivo similar mediante coeficientes de atención aprendibles que capturan implícitamente la estructura del grafo y dependencias de largo alcance.

\subsubsection{Fundamentos Matemáticos: Atención sobre Grafos}\label{subsubsec:fundamentos_atencion_grafos}

\begin{definition}[Grafo Etiquetado para Atención]\label{def:labeled_graph_attention}
Un \textit{grafo etiquetado para atención} es una tupla $\mathcal{G} = (V, E, \mathbf{H})$ donde:
\begin{itemize}
    \item $V = \{v_1, v_2, \ldots, v_n\}$ es el conjunto de nodos
    \item $E \subseteq V \times V$ es el conjunto de aristas
    \item $\mathbf{H} = [\mathbf{h}_1, \mathbf{h}_2, \ldots, \mathbf{h}_n]^T \in \mathbb{R}^{n \times d}$ son las características de los nodos
\end{itemize}
La estructura del grafo define las restricciones de atención: típicamente, el nodo $v_i$ puede atender únicamente a su vecindario $\mathcal{N}(i) = \{j : (i,j) \in E\}$.
\end{definition}

\begin{definition}[Función de Atención en Grafos]\label{def:graph_attention_function}
Una \textit{función de atención en grafos} $f_{\text{GAT}}: \mathcal{G} \to \mathbb{R}^{n \times d'}$ satisface:
\begin{enumerate}
    \item \textbf{Equivarianza a permutaciones}: Para cualquier permutación $\pi$ de los nodos:
    \begin{equation}
        f_{\text{GAT}}(\pi(\mathcal{G})) = \pi(f_{\text{GAT}}(\mathcal{G}))
        \label{eq:gat_permutation_equivariance}
    \end{equation}

    \item \textbf{Respeto de la estructura}: Solo considera conexiones definidas por $E$:
    \begin{equation}
        \frac{\partial f_{\text{GAT}}(\mathbf{h}_i)}{\partial \mathbf{h}_j} \neq \mathbf{0} \Rightarrow (i,j) \in E \text{ o } i = j
        \label{eq:gat_structural_constraint}
    \end{equation}
\end{enumerate}
\end{definition}

\subsubsection{Graph Attention Networks Clásicos}\label{subsubsec:gat_clasicos}

El trabajo pionero de \citet{velickovic2017graph} introduce la primera formulación exitosa de atención específica para grafos.

\begin{definition}[Mecanismo de Atención GAT]\label{def:gat_mechanism}
Para un nodo $i$ con vecindario $\mathcal{N}(i)$, el mecanismo \gls{GAT} computa:

\textbf{Paso 1: Proyección lineal}
\begin{equation}
    \mathbf{z}_i = \mathbf{W}\mathbf{h}_i, \quad \mathbf{z}_j = \mathbf{W}\mathbf{h}_j
    \label{eq:gat_linear_projection}
\end{equation}
donde $\mathbf{W} \in \mathbb{R}^{d' \times d}$ es una matriz de peso compartida.

\textbf{Paso 2: Coeficientes de atención}
\begin{equation}
    e_{ij} = a^T[\mathbf{z}_i \| \mathbf{z}_j] = a^T[\mathbf{W}\mathbf{h}_i \| \mathbf{W}\mathbf{h}_j]
    \label{eq:gat_attention_coefficients}
\end{equation}
donde $a \in \mathbb{R}^{2d'}$ es un vector de atención aprendible y $\|$ denota concatenación.

\textbf{Paso 3: Normalización softmax enmascarada}
\begin{equation}
    \alpha_{ij} = \frac{\exp(\text{LeakyReLU}(e_{ij}))}{\sum_{k \in \mathcal{N}(i) \cup \{i\}} \exp(\text{LeakyReLU}(e_{ik}))}
    \label{eq:gat_softmax_normalization}
\end{equation}

\textbf{Paso 4: Agregación ponderada}
\begin{equation}
    \mathbf{h}_i' = \sigma\left(\sum_{j \in \mathcal{N}(i) \cup \{i\}} \alpha_{ij} \mathbf{z}_j\right)
    \label{eq:gat_weighted_aggregation}
\end{equation}
\end{definition}

\begin{figure}[ht]
\centering
\includegraphics[width=0.75\textwidth]{figures/gat_mechanism.pdf}
\caption{Mecanismo de Graph Attention Network (GAT). El nodo central (rojo) computa coeficientes de atención $\alpha_{ij}$ con cada vecino (azul). El grosor de las flechas refleja la magnitud de los pesos de atención: la red aprende a asignar diferentes importancias a cada vecino según las características de los nodos.}
\label{fig:gat_mechanism}
\end{figure}

\begin{proposition}[Poder Expresivo de GAT]\label{prop:gat_expressive_power}
Los Graph Attention Networks estándar están acotados en poder expresivo por el test 1-Weisfeiler-Lehman (1-WL) \citep{xu2019how}. Esto implica:
\begin{enumerate}
    \item \textbf{Cota superior}: GAT no puede distinguir pares de grafos no isomorfos que el test 1-WL no distingue.
    \item \textbf{Agregación anisotrópica}: A diferencia de \gls{GCN} con pesos fijos, GAT puede aprender funciones de agregación adaptativas mediante coeficientes de atención, lo que permite mayor flexibilidad dentro de la clase de funciones accesibles por arquitecturas de message passing.
    \item \textbf{Equivarianza a permutaciones}: GAT es permutación-equivariante por construcción, ya que los coeficientes de atención dependen únicamente de las características de los nodos involucrados y no de un ordenamiento.
\end{enumerate}
\end{proposition}

\begin{proof}
GAT es una instancia del framework MPNN (message passing neural network) con agregación ponderada por atención sobre el vecindario de 1-salto. \citet{xu2019how} demuestran que toda arquitectura MPNN está acotada por el test 1-WL en su capacidad de distinguir grafos. La equivarianza a permutaciones se sigue de que la operación de atención es invariante al reordenamiento de los nodos vecinos: los coeficientes $\alpha_{ij}$ dependen de $(\mathbf{h}_i, \mathbf{h}_j)$ y la agregación $\sum_{j \in \mathcal{N}(i)} \alpha_{ij} \mathbf{z}_j$ es simétrica respecto al orden de los vecinos.
\end{proof}

\subsubsection{Multi-Head Graph Attention}\label{subsubsec:multihead_graph_attention}

Para mejorar la estabilidad y capacidad expresiva, los \glspl{GAT} emplean atención multi-cabeza.

\begin{definition}[Multi-Head GAT]\label{def:multihead_gat}
Un Multi-Head GAT con $h$ cabezas computa:

\begin{equation}
    \mathbf{h}_i' = \sigma\left(\frac{1}{h}\sum_{k=1}^{h} \sum_{j \in \mathcal{N}(i)} \alpha_{ij}^{(k)} \mathbf{W}^{(k)}\mathbf{h}_j\right)
    \label{eq:multihead_gat_average}
\end{equation}

o mediante concatenación:

\begin{equation}
    \mathbf{h}_i' = \sigma\left(\|_{k=1}^{h} \sum_{j \in \mathcal{N}(i)} \alpha_{ij}^{(k)} \mathbf{W}^{(k)}\mathbf{h}_j\right)
    \label{eq:multihead_gat_concat}
\end{equation}

donde cada cabeza $k$ tiene sus propios parámetros $\mathbf{W}^{(k)}$ y $\mathbf{a}^{(k)}$.
\end{definition}

\begin{proposition}[Ventajas del Multi-Head GAT]\label{prop:multihead_gat_advantages}
El Multi-Head GAT proporciona:
\begin{enumerate}
    \item \textbf{Estabilidad de entrenamiento}: Reduce la varianza de los gradientes
    \item \textbf{Especialización de características}: Diferentes cabezas capturan diferentes aspectos relacionales
    \item \textbf{Mayor capacidad expresiva}: Aproxima funciones más complejas con el mismo número total de parámetros
\end{enumerate}
\end{proposition}

\begin{proof}[Justificación]
\textbf{(1)} Sea $\hat{\mathbf{h}}_i^{(k)} = \frac{1}{H}\sum_{h=1}^H \mathbf{h}_i^{(k,h)}$ la salida promediada. La varianza del gradiente $\text{Var}[\nabla_\theta \hat{\mathbf{h}}_i]$ se reduce por un factor $1/H$ respecto a una única cabeza, por la independencia de los mecanismos de atención en cada cabeza (análogo al efecto de los \textit{ensemble methods}).

\textbf{(2)} Empíricamente, las cabezas tienden a especializarse: con $H$ cabezas y dimensión $d/H$ cada una, cada cabeza opera en un subespacio diferente de $\mathbb{R}^d$. Los pesos de atención $\alpha_{ij}^{(h)}$ varían significativamente entre cabezas, indicando que capturan diferentes patrones relacionales.

\textbf{(3)} El espacio de funciones representables por $H$ cabezas de dimensión $d/H$ incluye estrictamente al de una cabeza de dimensión $d$: la concatenación $\|_{h=1}^H \text{Attn}_h$ puede representar funciones que requieren atender simultáneamente a distintas combinaciones de vecinos, lo cual una única cabeza no puede lograr.
\end{proof}

\subsubsection{Análisis Espectral de Graph Attention}\label{subsubsec:analisis_espectral_gat}

La conexión entre GATs y análisis espectral revela perspectivas profundas sobre su funcionamiento.

\begin{definition}[Interpretación Espectral de Atención]\label{def:spectral_interpretation_attention}
Sea $\mathbf{A}^{(GAT)} \in \mathbb{R}^{n \times n}$ la matriz de atención con entradas $A^{(GAT)}_{ij} = \alpha_{ij}$. Esta matriz puede interpretarse como una \textbf{versión aprendible del operador de difusión}:
\begin{equation}
    \mathbf{A}^{(GAT)} \approx f(\mathbf{L}^{-1})
    \label{eq:gat_spectral_approximation}
\end{equation}
donde $f$ es una función aprendible y $\mathbf{L}$ es el \gls{laplaciano_grafo}.
\end{definition}

\begin{theorem}[Conexión GAT-Espectral]\label{thm:gat_spectral_connection}
En el límite de muchas capas, un \gls{GAT} con pesos de atención estables converge a realizar filtrado espectral adaptativo:
\begin{equation}
    \lim_{L \to \infty} \text{GAT}^{(L)}(\mathbf{H}) = \mathbf{U} \text{diag}(\hat{g}(\lambda_1), \ldots, \hat{g}(\lambda_n)) \mathbf{U}^T \mathbf{H}
    \label{eq:gat_spectral_limit}
\end{equation}
donde $\mathbf{U}$ contiene los eigenvectores del Laplaciano y $\hat{g}$ es la función de filtro aprendida implícitamente.
\end{theorem}

\begin{proof}
\textbf{Paso 1}: Estabilidad de los pesos de atención.
En el límite, los pesos de atención convergen: $\lim_{L \to \infty} \alpha_{ij}^{(L)} = \alpha_{ij}^*$.

\textbf{Paso 2}: Matriz de transición límite.
La operación \gls{GAT} se convierte en: $\mathbf{H}^{(L+1)} = \mathbf{A}^* \mathbf{H}^{(L)}$ donde $\mathbf{A}^*$ es la matriz de atención estable.

\textbf{Paso 3}: Descomposición espectral.
Bajo la hipótesis (idealizada) de que $\mathbf{A}^*$ comparte eigenvectores con el Laplaciano del grafo: $\mathbf{A}^* = \mathbf{U} \mathbf{\Lambda}^* \mathbf{U}^T$.

\textbf{Paso 4}: Filtrado implícito.
Los eigenvalores $\mathbf{\Lambda}^*$ definen implícitamente la función de filtro $\hat{g}$.
\end{proof}

\begin{remark}
El Paso~3 de la demostración anterior supone que la matriz de atención aprendida $\mathbf{A}^*$ comparte eigenvectores con el Laplaciano del grafo $\mathbf{L}$, lo cual es una simplificación. En la práctica, $\mathbf{A}^*$ es generalmente asimétrica y sus eigenvectores no coinciden con los de $\mathbf{L}$. Esta conexión espectral debe entenderse como una analogía que se cumple exactamente solo cuando los coeficientes de atención son funciones de la estructura del grafo compatibles con su descomposición espectral. Además, la aplicación repetida de una matriz estocástica fija $\mathbf{A}^*$ converge a una distribución estacionaria (fenómeno de over-smoothing), no al filtrado espectral rico descrito en la ecuación~\eqref{eq:gat_spectral_limit}.
\end{remark}

\subsubsection{Codificación Posicional Espectral}\label{subsubsec:codificacion_posicional_espectral}

Una limitación fundamental de los \glspl{GAT} es su incapacidad para distinguir nodos estructuralmente equivalentes. La codificación posicional espectral resuelve este problema.

\begin{definition}[Codificación Posicional Espectral]\label{def:spectral_positional_encoding}
Sea $\mathbf{L} = \mathbf{U}\mathbf{\Lambda}\mathbf{U}^T$ la descomposición espectral del \gls{laplaciano_grafo}. La codificación posicional espectral del nodo $i$ se define como:
\begin{equation}
    \text{PE}_i = [\mathbf{u}_{i,1}, \mathbf{u}_{i,2}, \ldots, \mathbf{u}_{i,k}] \in \mathbb{R}^k
    \label{eq:spectral_positional_encoding}
\end{equation}
donde $\mathbf{u}_{i,j}$ es el $i$-ésimo componente del $j$-ésimo eigenvector más pequeño (excluyendo el trivial).
\end{definition}

\begin{definition}[Distancias Espectrales para Atención]\label{def:spectral_distances_attention}
Para mejorar la atención entre nodos distantes, se definen métricas espectrales:

\textbf{Distancia de Difusión}:
\begin{equation}
    d_D^2(i,j) = \sum_{k=2}^{n} e^{-2t\lambda_k} (\mathbf{u}_{i,k} - \mathbf{u}_{j,k})^2
    \label{eq:diffusion_distance}
\end{equation}

\textbf{Función de Green}:
\begin{equation}
    G(i,j) = \sum_{k=2}^{n} \frac{\mathbf{u}_{i,k} \mathbf{u}_{j,k}}{\lambda_k}
    \label{eq:green_function_distance}
\end{equation}

\textbf{Distancia Biarmónica}:
\begin{equation}
    d_B^2(i,j) = \sum_{k=2}^{n} \frac{(\mathbf{u}_{i,k} - \mathbf{u}_{j,k})^2}{\lambda_k^2}
    \label{eq:biharmonic_distance}
\end{equation}
\end{definition}

\subsubsection{Spectral Attention Networks (SAN)}\label{subsubsec:spectral_attention_networks}

Los \citet{kreuzer2021spectral} introducen el primer Transformer completamente conectado para grafos que compite con GNNs especializados.

\begin{definition}[Arquitectura SAN]\label{def:san_architecture}
Un Spectral Attention Network procesa grafos mediante:

\textbf{Paso 1: Codificación posicional aprendible}
\begin{equation}
    \text{LPE}_i = \text{MLP}_{\text{PE}}([\lambda_1, \ldots, \lambda_k, \mathbf{u}_{i,1}, \ldots, \mathbf{u}_{i,k}])
    \label{eq:learned_positional_encoding}
\end{equation}

\textbf{Paso 2}: Atención dual (real + virtual)
\begin{align}
    \mathbf{h}_i^{\text{real}} &= \text{Attention}_{\text{sparse}}(\mathbf{h}_i, \{\mathbf{h}_j : j \in \mathcal{N}(i)\}) \\
    \mathbf{h}_i^{\text{virtual}} &= \text{Attention}_{\text{full}}(\mathbf{h}_i, \{\mathbf{h}_j : j \in V\})
    \label{eq:dual_attention_san}
\end{align}

\textbf{Paso 3: Balanceado adaptativo}
\begin{equation}
    \mathbf{h}_i' = \gamma \cdot \mathbf{h}_i^{\text{real}} + (1-\gamma) \cdot \mathbf{h}_i^{\text{virtual}}
    \label{eq:adaptive_balancing_san}
\end{equation}
donde $\gamma \in [0,1]$ es un parámetro aprendible.
\end{definition}

\begin{theorem}[Superación del Over-Squashing]\label{thm:over_squashing_resolution}
Los Spectral Attention Networks resuelven el problema del over-squashing al proporcionar conexiones directas entre nodos arbitrariamente distantes:
\begin{equation}
    \text{Distancia de información} = \mathcal{O}(1) \quad \forall i,j \in V
    \label{eq:information_distance_constant}
\end{equation}
independientemente de la distancia geodésica en el grafo original.
\end{theorem}

\begin{proof}
En una \gls{GNN} con message passing local, la información del nodo $j$ tarda al menos $d_G(i,j)$ capas en alcanzar el nodo $i$, donde $d_G$ es la distancia geodésica. Tras $L$ capas, la información se comprime en representaciones de dimensión fija $d$, causando over-squashing cuando $|\{j: d_G(i,j) \leq L\}|$ crece exponencialmente \citep{topping2022understanding}. En SAN, la atención global permite que el nodo $i$ atienda directamente al nodo $j$ en una sola capa, con score $\alpha_{ij} = \text{softmax}(\mathbf{q}_i^T \mathbf{k}_j / \sqrt{d})$. La distancia de información $D_{\text{info}}(i,j) = -\log \alpha_{ij}$ es $O(1)$ (depende de la similitud de representaciones, no de la topología del grafo), eliminando el cuello de botella.
\end{proof}

\subsubsection{Extensiones Escalables: GraphGPS}\label{subsubsec:graphgps}

\begin{definition}[Arquitectura GraphGPS]\label{def:graphgps_architecture}
GraphGPS \citep{rampasek2022graphgps} combina message passing local con atención global mediante:

\begin{equation}
    \mathbf{H}^{(\ell+1)} = \text{MLP}^{(\ell)}\left(\mathbf{H}_{\text{MPNN}}^{(\ell+1)} + \mathbf{H}_{\text{Global}}^{(\ell+1)}\right)
    \label{eq:graphgps_architecture}
\end{equation}

donde:
\begin{align}
    \mathbf{H}_{\text{MPNN}}^{(\ell+1)} &= \text{MPNN}^{(\ell)}(\mathbf{H}^{(\ell)}) \\
    \mathbf{H}_{\text{Global}}^{(\ell+1)} &= \text{GlobalTransformer}^{(\ell)}(\mathbf{H}^{(\ell)} + \text{PE})
    \label{eq:graphgps_components}
\end{align}
\end{definition}

\begin{proposition}[Escalabilidad Lineal de GraphGPS]\label{prop:graphgps_scalability}
Cuando se utiliza un mecanismo de atención de complejidad lineal (Performer \citep{choromanski2021performer} o BigBird) para la componente global, y excluyendo el coste de precomputar codificaciones posicionales espectrales, GraphGPS logra complejidad por capa $\mathcal{O}(N + E)$ mediante:
\begin{enumerate}
    \item \textbf{Atención sub-cuadrática}: Performer o BigBird con complejidad $\mathcal{O}(N)$
    \item \textbf{Message passing local}: Complejidad $\mathcal{O}(E)$ para la componente MPNN
    \item \textbf{Separación modular}: Desacoplamiento de información local y global
\end{enumerate}
Con atención completa estándar, la componente global tendría complejidad $\mathcal{O}(N^2)$. Para el desglose preciso en función de los parámetros del modelo, véase la \autoref{prop:graphgps_efficient_complexity}.
\end{proposition}

\begin{proof}
La complejidad se obtiene sumando las contribuciones de cada componente. La componente MPNN procesa cada arista una vez, resultando en $\mathcal{O}(E)$. Para la componente global, mecanismos como Performer aproximan la atención softmax mediante $\text{softmax}(\mathbf{Q}\mathbf{K}^\top) \approx \hat{\phi}(\mathbf{Q})\hat{\phi}(\mathbf{K})^\top$, donde $\hat{\phi}: \mathbb{R}^d \to \mathbb{R}^k$ con $k \ll N$. Esto reduce el cómputo de $\mathcal{O}(N^2)$ a $\mathcal{O}(Nk) = \mathcal{O}(N)$ para $k$ fijo. La separación modular permite sumar ambas contribuciones independientemente, obteniendo $\mathcal{O}(N + E)$.
\end{proof}

\subsubsection{Limitaciones y Direcciones Futuras}\label{subsubsec:limitaciones_gat}

\paragraph{Limitaciones Actuales}

\begin{enumerate}
    \item \textbf{Limitación 1-WL}: Los GATs estándar no pueden distinguir grafos no-isomorfos que el test 1-Weisfeiler-Lehman no puede distinguir

    \item \textbf{Over-smoothing}: En grafos profundos, las representaciones nodales convergen a valores similares:
    \begin{equation}
        \lim_{L \to \infty} \|\mathbf{h}_i^{(L)} - \mathbf{h}_j^{(L)}\| \to 0 \quad \forall i,j
        \label{eq:over_smoothing_limit}
    \end{equation}

    \item \textbf{Ambigüedad de signo}: Los eigenvectores están definidos salvo multiplicación por $\pm 1$, requiriendo técnicas como SignNet

    \item \textbf{Grafos isoespectrales}: Grafos con espectros idénticos pero estructuras diferentes no pueden distinguirse espectralmente
\end{enumerate}

\paragraph{Investigación Futura}

\begin{definition}[k-GAT de Alto Orden]\label{def:higher_order_gat}
Los GATs de orden superior consideran $k$-tuplas de nodos:
\begin{equation}
    \alpha_{i,\mathcal{S}} = \text{softmax}\left(a^T\left[\|_{j \in \mathcal{S}} \mathbf{W}\mathbf{h}_j\right]\right)
    \label{eq:higher_order_gat}
\end{equation}
donde $\mathcal{S} \subseteq \mathcal{N}^{(k)}(i)$ es un $k$-subconjunto del $k$-vecindario.
\end{definition}

\begin{definition}[Atención Geométrica Continua]\label{def:continuous_geometric_attention}
Para grafos embebidos en espacios continuos:
\begin{equation}
    e_{ij} = a(\mathbf{h}_i, \mathbf{h}_j, \|\mathbf{x}_i - \mathbf{x}_j\|, \angle(\mathbf{x}_i, \mathbf{x}_j))
    \label{eq:continuous_geometric_attention}
\end{equation}
incorporando información métrica y angular.
\end{definition}

\subsubsection{Conexiones con Chapters Anteriores}

Los desarrollos en GATs conectan directamente con conceptos establecidos:

\begin{itemize}
    \item \textbf{Análisis Espectral (\autoref{sec:analisis_armonico})}: Los GATs implementan implícitamente filtrado espectral adaptativo
    \item \textbf{Message Passing (Sección~\ref{subsec:message_passing_unificador})}: Los GATs son una instanciación específica del marco unificador con agregación por atención
    \item \textbf{Equivarianza (Sección~\ref{subsec:equivarianza_invarianza})}: Los GATs preservan equivarianza a permutaciones mientras ganan flexibilidad adaptativa
\end{itemize}

\begin{proposition}[Síntesis GAT en Marco Unificador]\label{prop:gat_unified_framework}
Los Graph Attention Networks se expresan naturalmente en el marco de Message Passing Geométrico:
\begin{align}
    \mathbf{h}_i^{(k+1)} &= \phi^{(k)}\left(\mathbf{h}_i^{(k)}, \bigoplus_{j \in \mathcal{N}(i)} \alpha_{ij} \cdot \psi^{(k)}(\mathbf{h}_j^{(k)})\right) \\
    \alpha_{ij} &= \text{Attention}(\mathbf{h}_i^{(k)}, \mathbf{h}_j^{(k)}, e_{ij})
    \label{eq:gat_in_unified_framework}
\end{align}
donde los pesos de atención $\alpha_{ij}$ reemplazan la agregación fija con agregación adaptativa aprendible.
\end{proposition}

\begin{proof}
Sea $G = (V, E)$ un grafo y $\mathbf{h}_i^{(k)}$ la representación del nodo $i$ en la capa $k$. En el marco MPNN estándar, la actualización toma la forma $\mathbf{h}_i^{(k+1)} = \phi(\mathbf{h}_i^{(k)}, \bigoplus_{j \in \mathcal{N}(i)} \psi(\mathbf{h}_j^{(k)}))$. En un GAT, la función de mensaje es $\psi^{(k)}(\mathbf{h}_j) = \mathbf{W}^{(k)} \mathbf{h}_j$ y la agregación es $\bigoplus = \sum_{j} \alpha_{ij} (\cdot)$, donde $\alpha_{ij} = \operatorname{softmax}_j(e_{ij})$ con $e_{ij} = \text{LeakyReLU}(\mathbf{a}^\top [\mathbf{W}\mathbf{h}_i \| \mathbf{W}\mathbf{h}_j])$. Sustituyendo en la ecuación general del MPNN se obtiene directamente la expresión de la proposición.
\end{proof}

Esta formulación demuestra que los GATs representan la evolución natural de message passing de agregación rígida hacia agregación flexible, manteniendo garantías de equivarianza mientras ganando expresividad adaptativa.

\subsection{Análisis Espectral en Grafos}

Para los fundamentos de análisis espectral del Laplaciano, incluyendo eigenvalores, eigenvectores y descomposición espectral, véase el \autoref{ch:cnn}, Sección de Análisis Espectral Detallado. La codificación posicional espectral (\autoref{def:spectral_positional_encoding}) y las distancias espectrales (\autoref{def:spectral_distances_attention}) introducidas anteriormente constituyen las herramientas fundamentales para incorporar información estructural del grafo en mecanismos de atención.

\subsection{Arquitecturas Avanzadas de Grafos}\label{subsec:arquitecturas_avanzadas_grafos}

Las arquitecturas avanzadas de grafos representan la confluencia de mecanismos de atención con los principios fundamentales del \gls{GDL} desarrollados en el \autoref{ch:gdl}. Esta sección examina cómo diferentes diseños arquitectónicos integran información espectral, estructural y geométrica para superar las limitaciones de las arquitecturas básicas como GAT. Cada arquitectura aborda desafíos específicos: los Spectral Attention Networks resuelven el problema de la ambigüedad posicional, GraphGPS combina procesamiento local y global eficientemente, los Graph \glspl{Transformer} incorporan estructura topológica en la atención, y los Structure-Aware \glspl{Transformer} explotan patrones de subgrafos. Unificaremos estas arquitecturas mostrando cómo instancian el marco de Message Passing (Sección~\ref{subsec:message_passing_unificador}) con diferentes mecanismos de agregación y sesgos inductivos.

\subsubsection{Spectral Attention Networks}\label{subsubsec:san_advanced}

Los Spectral Attention Networks de \citet{kreuzer2021spectral} constituyen el primer \gls{Transformer} completamente conectado para grafos que compite con \glspl{GNN} especializadas. Su contribución fundamental radica en el uso de codificaciones posicionales aprendibles basadas en el espectro del Laplaciano del grafo.

\paragraph{Fundamento teórico: el problema de la isotropía}

Los \glspl{Transformer} estándar operan sobre secuencias con orden natural implícito, lo que proporciona información posicional intrínseca. Los grafos, sin embargo, carecen de ordenamiento canónico: son invariantes bajo permutaciones de nodos. Esta invarianza, aunque deseable desde la perspectiva de equivarianza (véase Sección~\ref{subsec:equivarianza_invarianza}), elimina toda información posicional estructural. El desafío central es: \textit{¿cómo dotar a un \gls{Transformer} de información posicional en un dominio sin orden natural?}

La respuesta de SAN es utilizar el \textbf{espectro del Laplaciano del grafo}. El Teorema Espectral (\autoref{ch:cnn}) garantiza que para el Laplaciano simétrico normalizado $\tilde{L} = I - D^{-1/2}AD^{-1/2}$ existe una descomposición:
\begin{equation}
    \tilde{L} = \sum_{k=1}^{n} \lambda_k \mathbf{v}_k \mathbf{v}_k^T
    \label{eq:spectral_decomposition_san}
\end{equation}
donde $0 = \lambda_1 \leq \lambda_2 \leq \cdots \leq \lambda_n \leq 2$ son los eigenvalores y $\{\mathbf{v}_k\}_{k=1}^{n}$ los eigenvectores ortonormales correspondientes.

\begin{definition}[Codificación Posicional Espectral Aprendible]\label{def:learned_spectral_pe}
Para un grafo $G$ con Laplaciano $\tilde{L}$, la codificación posicional espectral aprendible para el nodo $i$ se define como:
\begin{equation}
    \mathbf{PE}_i = \text{MLP}_{\text{enc}}(\lambda_1, \ldots, \lambda_k, \mathbf{v}_{1,i}, \ldots, \mathbf{v}_{k,i}) \in \mathbb{R}^{d_{\text{model}}}
    \label{eq:san_learned_positional_encoding}
\end{equation}
donde:
\begin{itemize}
    \item $k \ll n$ es el número de componentes espectrales utilizadas (típicamente $k \in [8, 20]$)
    \item $\text{MLP}_{\text{enc}}$ es una red multicapa con función de activación no lineal
    \item $\lambda_1, \ldots, \lambda_k$ son los primeros $k$ eigenvalores (información global del grafo)
    \item $\mathbf{v}_{k,i}$ es la $i$-ésima componente del eigenvector $\mathbf{v}_k$ (información posicional local del nodo)
\end{itemize}
\end{definition}

\textbf{Justificación espectral.}\label{rem:justification_spectral_pe} Los primeros eigenvectores del Laplaciano capturan estructura topológica fundamental:
\begin{itemize}
    \item $\mathbf{v}_1 = \frac{1}{\sqrt{n}}\mathbf{1}$ (vector constante, componente conectada)
    \item $\mathbf{v}_2$ (Fiedler vector) identifica comunidades y cortes mínimos
    \item Eigenvectores subsecuentes capturan patrones geométricos de frecuencia creciente
\end{itemize}
Este enfoque conecta con análisis armónico en grafos (\autoref{ch:cnn}), donde los eigenvectores forman una base de Fourier para señales sobre el grafo.

\paragraph{Arquitectura SAN: atención dual con balanceo adaptativo}

La arquitectura completa de SAN implementa un esquema de \textbf{atención dual} que combina información local (vecindad del grafo) con información global (grafo completo):

\begin{definition}[Capa SAN]\label{def:san_layer}
Una capa SAN procesa características de nodos $\mathbf{H}^{(\ell)} \in \mathbb{R}^{n \times d}$ mediante tres etapas:

\textbf{Etapa 1: Inyección de codificación posicional}
\begin{equation}
    \tilde{\mathbf{H}}^{(\ell)} = \mathbf{H}^{(\ell)} + \mathbf{PE}
    \label{eq:san_pe_injection}
\end{equation}

\textbf{Etapa 2: Atención dual}
\begin{align}
    \mathbf{H}_{\text{local}}^{(\ell+1)} &= \text{MultiHeadAttn}_{\text{sparse}}\left(\tilde{\mathbf{H}}^{(\ell)}, \mathcal{E}\right) \label{eq:san_local_attn} \\
    \mathbf{H}_{\text{global}}^{(\ell+1)} &= \text{MultiHeadAttn}_{\text{full}}\left(\tilde{\mathbf{H}}^{(\ell)}\right) \label{eq:san_global_attn}
\end{align}
donde $\mathcal{E}$ denota las aristas del grafo y la atención esparsa solo computa pesos para pares $(i,j) \in \mathcal{E}$.

\textbf{Etapa 3: Combinación adaptativa}
\begin{equation}
    \mathbf{H}^{(\ell+1)} = \text{MLP}\left(\text{LayerNorm}\left(\gamma^{(\ell)} \mathbf{H}_{\text{local}}^{(\ell+1)} + (1-\gamma^{(\ell)}) \mathbf{H}_{\text{global}}^{(\ell+1)}\right)\right)
    \label{eq:san_adaptive_combination}
\end{equation}
donde $\gamma^{(\ell)} \in [0,1]$ es un parámetro aprendible por capa.
\end{definition}

\begin{proposition}[Complejidad computacional de SAN]\label{prop:san_complexity}
Una capa SAN con $n$ nodos, $m$ aristas, $h$ cabezas de atención y dimensión $d$ tiene complejidad:
\begin{equation}
    \mathcal{O}(hd(m + n^2)) = \mathcal{O}(hdm) + \mathcal{O}(hdn^2)
    \label{eq:san_complexity}
\end{equation}
El término cuadrático $\mathcal{O}(hdn^2)$ domina para grafos densos, limitando la escalabilidad a grafos con miles de nodos.
\end{proposition}

\begin{proof}
La atención local opera sobre las $m$ aristas del grafo: para cada arista $(i,j)$, cada cabeza computa scores de atención y mensajes en $\mathcal{O}(d)$, dando $\mathcal{O}(hdm)$. La atención global computa la matriz $\mathbf{Q}\mathbf{K}^\top \in \mathbb{R}^{n \times n}$ para cada cabeza, requiriendo $\mathcal{O}(dn^2)$ operaciones por cabeza y $\mathcal{O}(hdn^2)$ en total. La multiplicación posterior con $\mathbf{V}$ también es $\mathcal{O}(hdn^2)$. Sumando ambas contribuciones se obtiene $\mathcal{O}(hd(m + n^2))$.
\end{proof}

\subsubsection{GraphGPS: Diseño Modular y Escalable}\label{subsubsec:graphgps_advanced}

GraphGPS (Graph Generalized Positional Signals) de \citet{rampasek2022graphgps} representa una evolución hacia arquitecturas modulares y componibles. Su filosofía de diseño se basa en tres principios:

\begin{enumerate}
    \item \textbf{Separación de responsabilidades}: Componentes especializados para procesamiento local vs. global
    \item \textbf{Escalabilidad}: Uso de mecanismos de atención eficientes (Performer, BigBird)
    \item \textbf{Flexibilidad}: Intercambiabilidad de componentes MPNN y atención
\end{enumerate}

\paragraph{Arquitectura modular}

\begin{definition}[Capa GraphGPS]\label{def:graphgps_layer}
Una capa GraphGPS consta de dos caminos paralelos que se fusionan:
\begin{equation}
    \mathbf{H}^{(\ell+1)} = \mathbf{H}^{(\ell)} + \text{MLP}^{(\ell)}\left(\text{LayerNorm}\left(\mathbf{H}_{\text{MPNN}}^{(\ell)} + \mathbf{H}_{\text{Global}}^{(\ell)}\right)\right)
    \label{eq:graphgps_main_layer}
\end{equation}

\textbf{Camino local (Message Passing):}
\begin{equation}
    \mathbf{H}_{\text{MPNN}}^{(\ell)} = \text{MPNN}^{(\ell)}\left(\mathbf{H}^{(\ell)}, \mathcal{E}\right)
    \label{eq:graphgps_mpnn_path}
\end{equation}
donde $\text{MPNN}^{(\ell)}$ puede ser \gls{GCN}, \gls{GAT}, GIN, o cualquier arquitectura de message passing.

\textbf{Camino global (Atención con PE):}
\begin{equation}
    \mathbf{H}_{\text{Global}}^{(\ell)} = \text{Transformer}^{(\ell)}\left(\mathbf{H}^{(\ell)} + \mathbf{PE}\right)
    \label{eq:graphgps_global_path}
\end{equation}
donde el \gls{Transformer} puede usar atención completa, Performer, o BigBird para eficiencia.
\end{definition}

\textbf{Interpretación como suma de agregaciones.}\label{rem:graphgps_interpretation} GraphGPS implementa una descomposición aditiva de la agregación en el marco de Message Passing (Sección~\ref{subsec:message_passing_unificador}):
\begin{equation}
    m_i^{(\ell)} = \underbrace{\bigoplus_{j \in \mathcal{N}(i)} \psi_{\text{local}}(\mathbf{h}_j^{(\ell)})}_{\text{Agregación local (MPNN)}} + \underbrace{\bigoplus_{j \in V} \alpha_{ij}^{(\ell)} \psi_{\text{global}}(\mathbf{h}_j^{(\ell)})}_{\text{Agregación global (Atención)}}
    \label{eq:graphgps_message_decomposition}
\end{equation}
Esta descomposición permite capturar tanto patrones locales (estructura topológica) como dependencias de largo alcance (interacciones globales).

\paragraph{Mecanismos de escalabilidad}

\begin{proposition}[Complejidad de GraphGPS con atención eficiente]\label{prop:graphgps_efficient_complexity}
Usando Performer para atención global, GraphGPS alcanza complejidad:
\begin{equation}
    \mathcal{O}(hdm + ndk) \quad \text{donde } k \ll n
    \label{eq:graphgps_linear_complexity}
\end{equation}
con $k$ el rango de la aproximación de Performer. Para grafos con $m = \mathcal{O}(n)$, la complejidad total es $\mathcal{O}(ndk)$, lineal en el número de nodos.
\end{proposition}

\begin{proof}
La componente MPNN procesa las $m$ aristas con $h$ cabezas de dimensión $d$, contribuyendo $\mathcal{O}(hdm)$. Para la componente global, Performer reemplaza el cálculo explícito de $\mathbf{Q}\mathbf{K}^\top \in \mathbb{R}^{n \times n}$ por la aproximación $\hat{\phi}(\mathbf{Q})(\hat{\phi}(\mathbf{K})^\top \mathbf{V})$, donde $\hat{\phi}: \mathbb{R}^d \to \mathbb{R}^k$. El producto $\hat{\phi}(\mathbf{K})^\top \mathbf{V} \in \mathbb{R}^{k \times d}$ se computa en $\mathcal{O}(ndk)$ y luego $\hat{\phi}(\mathbf{Q}) \cdot (\hat{\phi}(\mathbf{K})^\top \mathbf{V})$ también en $\mathcal{O}(ndk)$. Sumando: $\mathcal{O}(hdm + ndk)$. Si $m = \mathcal{O}(n)$ y $h, k$ son constantes respecto a $n$, esto se reduce a $\mathcal{O}(ndk)$.
\end{proof}

\begin{example}[Configuración típica de GraphGPS]\label{ex:graphgps_configuration}
Una configuración estándar utiliza:
\begin{itemize}
    \item \textbf{MPNN local}: GatedGCN (8 capas)
    \item \textbf{Atención global}: Performer con $k=256$ características aleatorias
    \item \textbf{Codificación posicional}: Laplacian eigenvectors (10 componentes) + Random Walk Structural Encoding
    \item \textbf{Dimensión del modelo}: $d=512$
    \item \textbf{Cabezas de atención}: $h=8$
\end{itemize}
Esta configuración procesa grafos con $n \sim 10^4$ nodos en tiempo razonable.
\end{example}

\subsubsection{Graph \gls{Transformer} Networks}\label{subsubsec:graph_transformers}

Los Graph \glspl{Transformer} de \citet{dwivedi2021generalization} abordan una pregunta fundamental: \textit{¿cómo incorporar la estructura topológica del grafo directamente en el mecanismo de atención?}

\begin{definition}[Graph \gls{Transformer} con sesgo estructural]\label{def:graph_transformer_bias}
Un Graph \gls{Transformer} modifica la atención estándar incorporando la matriz de adyacencia $A \in \{0,1\}^{n \times n}$:
\begin{equation}
    \text{GraphAttn}(\mathbf{Q}, \mathbf{K}, \mathbf{V}, A) = \text{softmax}\left(\frac{\mathbf{Q}\mathbf{K}^T}{\sqrt{d_k}} \odot M(A)\right)\mathbf{V}
    \label{eq:graph_transformer_masked_attention}
\end{equation}
donde $M(A)$ es una máscara derivada de $A$ y $\odot$ es el producto elemento a elemento (Hadamard).
\end{definition}

\paragraph{Variantes de máscaras estructurales}

\begin{enumerate}
    \item \textbf{Máscara binaria rígida} (sparsificación):
    \begin{equation}
        M(A)_{ij} = \begin{cases}
            0 & \text{si } A_{ij} = 0 \\
            1 & \text{si } A_{ij} = 1
        \end{cases}
        \label{eq:binary_mask}
    \end{equation}
    Reduce la atención solo a vecinos conectados, recuperando comportamiento tipo \gls{GAT}.

    \item \textbf{Máscara de distancia geodésica}:
    \begin{equation}
        M(A)_{ij} = \exp\left(-\beta \cdot d_G(i,j)\right)
        \label{eq:geodesic_mask}
    \end{equation}
    donde $d_G(i,j)$ es la distancia del camino más corto entre $i$ y $j$, y $\beta > 0$ controla la decaída.

    \item \textbf{Máscara de sesgo espectral}:
    \begin{equation}
        M(A)_{ij} = 1 + \sum_{k=1}^{K} \theta_k \lambda_k \mathbf{v}_{k,i} \mathbf{v}_{k,j}
        \label{eq:spectral_bias_mask}
    \end{equation}
    donde $\{\theta_k\}$ son parámetros aprendibles y $\{\lambda_k, \mathbf{v}_k\}$ el espectro del Laplaciano.
\end{enumerate}

\begin{theorem}[Equivalencia GAT-GraphTransformer]\label{thm:gat_graphtransformer_equivalence}
Un Graph \gls{Transformer} con máscara binaria rígida y una sola cabeza de atención es equivalente a un \gls{GAT} monocapa si las funciones de proyección $\mathbf{W}_Q, \mathbf{W}_K$ se parametrizan apropiadamente.
\end{theorem}

\begin{proof}[Bosquejo]
La atención \gls{GAT} se escribe como:
\begin{equation}
    \alpha_{ij} = \frac{\exp(e_{ij})}{\sum_{k \in \mathcal{N}(i)} \exp(e_{ik})}, \quad e_{ij} = a^T[\mathbf{W}\mathbf{h}_i \| \mathbf{W}\mathbf{h}_j]
\end{equation}
Reparametrizando $a$ y $\mathbf{W}$ como $\mathbf{W}_Q, \mathbf{W}_K$ con sesgo aditivo, se obtiene la forma de producto punto. La restricción a $\mathcal{N}(i)$ equivale a la máscara binaria.
\end{proof}

\subsubsection[Structure-Aware Transformers]{Structure-Aware \glspl{Transformer}}\label{subsubsec:structure_aware_transformers}

Los Structure-Aware \glspl{Transformer} (SAT) de \citet{chen2022structureaware} adoptan un enfoque radicalmente diferente: en lugar de codificar la estructura mediante posiciones o máscaras, extraen \textbf{representaciones explícitas de subgrafos} que se incorporan como características adicionales.

\begin{definition}[Representación de subgrafo local]\label{def:subgraph_representation_sat}
Para un nodo $i$ y radio $r \in \mathbb{N}$, el subgrafo $r$-local se define como:
\begin{equation}
    \mathcal{G}_i^{(r)} = (V_i^{(r)}, E_i^{(r)}), \quad V_i^{(r)} = \{j \in V : d_G(i,j) \leq r\}
    \label{eq:local_subgraph}
\end{equation}
donde $E_i^{(r)} = \{(u,v) \in E : u, v \in V_i^{(r)}\}$ son las aristas inducidas.

La representación estructural del nodo $i$ se obtiene aplicando una \gls{GNN} al subgrafo local:
\begin{equation}
    \mathbf{h}_i^{\text{struct}} = \text{GNN}_{\text{struct}}(\mathcal{G}_i^{(r)})
    \label{eq:structural_representation_sat}
\end{equation}
\end{definition}

\paragraph{Atención aumentada con información estructural}

\begin{definition}[SAT con características estructurales]\label{def:sat_attention}
SAT extiende la atención estándar concatenando características nodales con representaciones estructurales:
\begin{align}
    \tilde{\mathbf{h}}_i &= [\mathbf{h}_i \| \mathbf{h}_i^{\text{struct}}] \label{eq:sat_augmented_features} \\
    e_{ij}^{\text{SAT}} &= a\left(\mathbf{W}_Q \tilde{\mathbf{h}}_i, \mathbf{W}_K \tilde{\mathbf{h}}_j\right) \label{eq:sat_attention_score}
\end{align}
Esta formulación permite que la atención dependa tanto de características locales del nodo como de patrones estructurales de su vecindad.
\end{definition}

\textbf{Comparación con métodos espectrales.}\label{rem:sat_vs_spectral} Mientras que métodos espectrales (SAN, GraphGPS) codifican estructura global mediante eigenvectores, SAT captura patrones estructurales locales explícitos. Esta diferencia es fundamental:
\begin{itemize}
    \item \textbf{Ventaja de SAT}: Captura motivos y patrones locales (triángulos, cliques)
    \item \textbf{Ventaja de métodos espectrales}: Escalan mejor y capturan propiedades globales
\end{itemize}

\subsubsection{Marco Unificador y Análisis Comparativo}\label{subsubsec:unified_framework_advanced}

Todas las arquitecturas avanzadas se unifican bajo el marco de Message Passing Generalizado:

\begin{proposition}[Representación unificada de arquitecturas avanzadas]\label{prop:unified_advanced_architectures}
Cada arquitectura avanzada instancia el marco MPNN generalizado:
\begin{equation}
    \mathbf{h}_i^{(\ell+1)} = \phi^{(\ell)}\left(\mathbf{h}_i^{(\ell)}, \bigoplus_{j \in V} w_{ij}^{(\ell)} \psi^{(\ell)}(\mathbf{h}_j^{(\ell)}, \mathbf{c}_{ij})\right)
    \label{eq:unified_mpnn_generalized}
\end{equation}
con diferentes elecciones de pesos $w_{ij}^{(\ell)}$ y contexto $\mathbf{c}_{ij}$:
\begin{itemize}
    \item \textbf{SAN}: $w_{ij} = \gamma \alpha_{ij}^{\text{local}} \mathbb{I}_{(i,j) \in E} + (1-\gamma) \alpha_{ij}^{\text{global}}$, $\mathbf{c}_{ij} = \mathbf{PE}$
    \item \textbf{GraphGPS}: $w_{ij} = \alpha_{ij}^{\text{MPNN}} + \alpha_{ij}^{\text{Transformer}}$, descomposición aditiva
    \item \textbf{Graph Transformer}: $w_{ij} = \alpha_{ij} \cdot M(A)_{ij}$, máscara estructural
    \item \textbf{SAT}: $w_{ij} = \alpha_{ij}(\mathbf{h}_i^{\text{struct}}, \mathbf{h}_j^{\text{struct}})$, atención sobre características estructurales
\end{itemize}
\end{proposition}

\begin{proof}
Verificamos cada instanciación identificando las funciones $\phi$, $\psi$, $w_{ij}$ y $\mathbf{c}_{ij}$ del marco general.

\textbf{SAN}: La combinación adaptativa de la \autoref{def:san_layer} con parámetro $\gamma$ separa atención local (sobre aristas $E$) y global (sobre todos los pares), dando $w_{ij} = \gamma \alpha_{ij}^{\text{local}} \mathbb{I}_{(i,j) \in E} + (1-\gamma) \alpha_{ij}^{\text{global}}$. El contexto es la codificación posicional espectral.

\textbf{GraphGPS}: La \autoref{def:graphgps_layer} suma independientemente las salidas MPNN y Transformer, equivalente a $w_{ij} = \alpha_{ij}^{\text{MPNN}} + \alpha_{ij}^{\text{Transformer}}$ en la ecuación unificada.

\textbf{Graph Transformer}: La máscara estructural $M(A)_{ij}$ restringe la atención a pares conectados o dentro de un radio topológico, multiplicando los pesos de atención estándar.

\textbf{SAT}: Los scores de atención se computan sobre representaciones estructurales $\mathbf{h}_i^{\text{struct}}$ obtenidas por extracción de subgrafos, en lugar de sobre las representaciones nodales directas.
\end{proof}

\paragraph{Tabla comparativa de complejidades y expresividad}

\begin{table}[h]
\centering
\small
\begin{tabular}{|l|c|c|c|c|}
\hline
\textbf{Arquitectura} & \textbf{Complejidad} & \textbf{Expresividad} & \textbf{Escalabilidad} & \textbf{Información} \\
\hline
GAT & $\mathcal{O}(hdm)$ & 1-WL & Alta & Local \\
SAN & $\mathcal{O}(hdn^2)$ & $>$1-WL & Baja & Espectral global \\
GraphGPS & $\mathcal{O}(hdm + ndk)$ & $>$1-WL & Media-Alta & Local + Global \\
Graph Transformer & $\mathcal{O}(hdn^2)$ & 1-WL & Baja-Media & Topológica \\
SAT & $\mathcal{O}(hdm + nr^3)$ & $>$1-WL & Media & Motivos locales \\
\hline
\end{tabular}
\caption[Comparación de arquitecturas avanzadas]{Comparación de arquitecturas avanzadas. $n$ = nodos, $m$ = aristas, $d$ = dimensión, $h$ = cabezas, $k$ = rango Performer, $r$ = radio subgrafo.}
\label{tab:advanced_architectures_comparison}
\end{table}

\subsubsection{Limitaciones Fundamentales y Direcciones Futuras}\label{subsubsec:limitations_future}

\paragraph{Limitaciones comunes}

El análisis de \citet{ying2021transformers} identifica desafíos transversales:

\begin{enumerate}
    \item \textbf{Dependencia crítica de codificaciones posicionales}: Sin PE robustas, el rendimiento se degrada significativamente, especialmente en tareas de long-range dependencies.

    \item \textbf{Trade-off expresividad-escalabilidad}: Arquitecturas más expresivas (SAN, SAT) tienen complejidad cuadrática o cúbica, limitando aplicaciones a grafos pequeños/medianos.

    \item \textbf{Ambigüedad espectral}: Eigenvectores están definidos salvo signo ($\mathbf{v}_k \sim -\mathbf{v}_k$) y permutaciones en caso de eigenvalores repetidos, requiriendo técnicas de normalización (SignNet, BasisNet).

    \item \textbf{Grafos isoespectrales}: Grafos no isomorfos con espectros idénticos son indistinguibles para métodos espectrales puros.
\end{enumerate}

\paragraph{Direcciones de investigación abiertas}

\begin{itemize}
    \item \textbf{Atención equivariante}: Diseñar mecanismos de atención que preserven simetrías adicionales (rotaciones en $SO(3)$, reflexiones) para datos moleculares 3D.

    \item \textbf{Codificaciones posicionales topológicas}: Uso de homología persistente y números de Betti como características posicionales invariantes.

    \item \textbf{Atención jerárquica}: Mecanismos multi-escala que procesan grafos a diferentes niveles de granularidad (nodos, subgrafos, comunidades).

    \item \textbf{Conexión con teoría de grafos extremales}: Caracterización de clases de grafos donde diferentes arquitecturas alcanzan expresividad óptima.
\end{itemize}

\section{Geometría Diferencial y Variedades}\label{sec:differential_geometry_attention}

La extensión de mecanismos de atención a variedades diferenciables y espacios no-euclidianos abre nuevas posibilidades para modelar datos con estructura geométrica intrínseca.

\subsection{Espacios No-Euclidianos}

Para los fundamentos de variedades diferenciables y métricas riemannianas, véase el \autoref{ch:gdl}, \autoref{subsec:fundamentos_variedades}. En mecanismos de atención sobre variedades, la distancia geodésica reemplaza la distancia euclidiana en el cálculo de similaridades.

\subsubsection{Geometría Hiperbólica}

La geometría hiperbólica emerge naturalmente en el contexto de datos jerárquicos y estructuras de árbol, donde el volumen crece exponencialmente con la distancia desde la raíz.

\begin{definition}[Espacio Hiperbólico - Modelo del Hiperboloide]\label{def:espacio_hiperbolico_hiperboloide}
El espacio hiperbólico $n$-dimensional se representa mediante el modelo del hiperboloide:

\begin{equation}
    \mathbb{H}^n = \{\mathbf{x} \in \mathbb{R}^{n+1} : \langle \mathbf{x}, \mathbf{x} \rangle_{\mathcal{L}} = -1, x_0 > 0\}
    \label{eq:hyperboloid_model}
\end{equation}

donde $\langle \cdot, \cdot \rangle_{\mathcal{L}}$ es el producto interno de Lorentz:

\begin{equation}
    \langle \mathbf{x}, \mathbf{y} \rangle_{\mathcal{L}} = -x_0y_0 + \sum_{i=1}^{n} x_iy_i
    \label{eq:lorentz_inner_product}
\end{equation}
\end{definition}

\begin{definition}[Modelo de Poincaré]\label{def:modelo_poincare}
El modelo de Poincaré representa el espacio hiperbólico como el disco unitario:
\begin{equation}
\mathbb{D}^n = \{\mathbf{x} \in \mathbb{R}^n : \|\mathbf{x}\| < 1\}
\end{equation}
con métrica Riemanniana:
\begin{equation}
g_{ij} = \frac{4}{(1-\|\mathbf{x}\|^2)^2} \delta_{ij}
\end{equation}

\textbf{Mapeo entre modelos}: La transformación del hiperboloide a Poincaré es:
\begin{equation}
\mathbf{x} = \frac{(x_1, \ldots, x_n)}{1 + x_0}
\end{equation}
La demostración de que esta transformación es una isometría riemanniana se encuentra en el \autoref{thm:isometria_hiperboloide_poincare} (Apéndice~\ref{ch:teoremas_complementarios}).
\end{definition}

\begin{theorem}[Propiedades Fundamentales del Espacio Hiperbólico]\label{thm:propiedades_espacio_hiperbolico}
El espacio hiperbólico $\mathbb{H}^n$ satisface:
\begin{enumerate}
    \item \textbf{Curvatura negativa constante}: $K = -1$
    \item \textbf{Crecimiento exponencial del volumen}: $\text{Vol}(B_r(p)) \sim e^{(n-1)r}$ para radio $r$ grande
    \item \textbf{Teorema de Gauss-Bonnet hiperbólico}: Para triángulos geodésicos con ángulos interiores $\alpha, \beta, \gamma$, el área es $\text{Area}(T) = \pi - (\alpha + \beta + \gamma) > 0$ y $\int_T K\, dA = (\alpha + \beta + \gamma) - \pi < 0$
    \item \textbf{Isometría}: El grupo de isometrías que preservan orientación y la componente $x_0 > 0$ es $SO^+(1,n)$ (componente conexa de la identidad de $SO(1,n)$)
\end{enumerate}
Las propiedades (1)-(3) son consecuencias directas de la métrica del hiperboloide; para demostraciones, véase \citep{do1976differential}.
\end{theorem}

\textbf{Implicaciones para aprendizaje automático}:
\begin{itemize}
    \item \textbf{Representación eficiente de jerarquías}: Los árboles se embebben naturalmente con distorsión mínima
    \item \textbf{Capacidad expresiva}: El crecimiento exponencial permite representar estructuras complejas
    \item \textbf{Regularización geométrica}: La curvatura negativa induce regularización implícita
\end{itemize}

\begin{definition}[Operaciones Hiperbólicas Fundamentales]\label{def:operaciones_hiperbolicas}
En el modelo de Poincaré, las operaciones básicas son:

\textbf{Suma de Möbius}:
\begin{equation}
\mathbf{x} \oplus \mathbf{y} = \frac{(1+2\langle\mathbf{x},\mathbf{y}\rangle+\|\mathbf{y}\|^2)\mathbf{x} + (1-\|\mathbf{x}\|^2)\mathbf{y}}{1+2\langle\mathbf{x},\mathbf{y}\rangle+\|\mathbf{x}\|^2\|\mathbf{y}\|^2}
\end{equation}

\textbf{Multiplicación escalar de Möbius}:
\begin{equation}
r \otimes \mathbf{x} = \tanh(r \tanh^{-1}(\|\mathbf{x}\|)) \frac{\mathbf{x}}{\|\mathbf{x}\|}
\end{equation}

\textbf{Mapa exponencial hiperbólico}:
\begin{equation}
\exp_{\mathbf{x}}^{\mathbb{H}}(\mathbf{v}) = \mathbf{x} \oplus \left(\tanh\left(\frac{\lambda_{\mathbf{x}} \|\mathbf{v}\|}{2}\right) \frac{\mathbf{v}}{\|\mathbf{v}\|}\right)
\end{equation}
donde $\lambda_{\mathbf{x}} = \frac{2}{1-\|\mathbf{x}\|^2}$ es el factor conformacional.
\end{definition}

\subsection{Atención en Espacios Hiperbólicos}

La atención hiperbólica extiende el mecanismo estándar de atención al explotar la estructura geométrica del espacio hiperbólico, particularmente útil para datos con estructuras jerárquicas inherentes.

\subsubsection{Embeddings Hiperbólicos para Grafos de Conocimiento}
\label{subsubsec:embeddings_hiperbolicos_grafos}

Antes de explorar arquitecturas de atención que operan en geometría hiperbólica, es fundamental comprender por qué el espacio hiperbólico constituye el marco geométrico natural para datos jerárquicos. Los grafos de conocimiento (WordNet, ontologías biomédicas) y las relaciones taxonómicas exhiben una estructura de árbol subyacente donde el número de nodos crece exponencialmente con la profundidad. Como muestra el \autoref{thm:propiedades_espacio_hiperbolico}, el volumen de bolas geodésicas en $\mathbb{H}^n$ crece como $e^{(n-1)r}$, lo que empareja la expansión combinatoria de los árboles. \citet{nickel2017poincare} demostraron empíricamente que embeddings en el disco de Poincaré $\mathbb{D}^d$ capturan jerarquías léxicas con distorsión significativamente menor que los embeddings euclidianos de la misma dimensión.

\begin{definition}[Embedding de Poincaré]\label{def:embedding_poincare}
Sea $\mathcal{E}$ un conjunto finito de entidades con relaciones jerárquicas. Un \textit{embedding de Poincaré} es una aplicación
\begin{equation}
    \Phi: \mathcal{E} \to \mathbb{D}^d = \{\mathbf{x} \in \mathbb{R}^d : \|\mathbf{x}\| < 1\}
\end{equation}
que asigna a cada entidad $e \in \mathcal{E}$ un punto $\Phi(e) \in \mathbb{D}^d$ en el disco de Poincaré $d$-dimensional, de modo que las distancias hiperbólicas $d_{\mathbb{D}}(\Phi(e_i), \Phi(e_j))$ reflejan las relaciones semánticas entre entidades: entidades relacionadas jerárquicamente quedan próximas, y la norma $\|\Phi(e)\|$ codifica la posición en la jerarquía (valores cercanos a $1$ para entidades específicas, cercanos a $0$ para entidades generales).
\end{definition}

\begin{theorem}[Embedding de Árboles en Espacio Hiperbólico]\label{thm:embedding_arboles_hiperbolico}
Todo árbol finito $T$ con $n$ vértices admite un embedding en el plano hiperbólico $\mathbb{H}^2$ con distorsión arbitrariamente pequeña. Más precisamente, para todo $\varepsilon > 0$ existe una aplicación $f: V(T) \to \mathbb{H}^2$ tal que para todo par de vértices $u, v \in V(T)$:
\begin{equation}
    d_T(u,v) \leq d_{\mathbb{H}}(f(u), f(v)) \leq (1+\varepsilon)\, d_T(u,v)
\end{equation}
donde $d_T$ es la distancia en el árbol \citep{sarkar2011low}.
\end{theorem}

\begin{proof}[Idea de la demostración]
El argumento central se basa en la correspondencia entre el crecimiento exponencial del espacio hiperbólico y la expansión combinatoria de árboles. Un árbol $b$-ario completo de profundidad $h$ tiene $\Theta(b^h)$ hojas. En $\mathbb{H}^2$, el área de un disco geodésico de radio $r$ es $2\pi(\cosh r - 1) \sim \pi e^r$, que crece exponencialmente. \citet{sarkar2011low} construyen el embedding recursivamente: fijan la raíz en el origen, y para cada nodo interno con $k$ hijos, distribuyen los hijos equidistantemente en un círculo hiperbólico centrado en el padre, explotando que el perímetro de un círculo hiperbólico de radio $r$ crece como $2\pi \sinh r \sim \pi e^r$, lo que proporciona espacio suficiente para separar las ramas sin distorsión. En contraste, el crecimiento polinomial $\Theta(r^d)$ del volumen en $\mathbb{R}^d$ impide esta separación en dimensión baja.
\end{proof}

\begin{proposition}[Contraste con el Caso Euclidiano]\label{prop:cota_inferior_euclidiano}
Sea $T_n$ un árbol binario completo con $n$ vértices. Todo embedding $f: V(T_n) \to \mathbb{R}^d$ con distorsión acotada por una constante $c \geq 1$ requiere dimensión $d = \Omega(\log n)$. En contraste, el \autoref{thm:embedding_arboles_hiperbolico} muestra que $\mathbb{H}^2$ basta con distorsión $(1+\varepsilon)$ para cualquier $\varepsilon > 0$, independientemente de $n$.
\end{proposition}

\begin{proof}[Idea de la demostración]
Un árbol binario completo de profundidad $h = \Theta(\log n)$ tiene $2^h$ hojas. En un embedding con distorsión $c$, las hojas deben estar a distancia mutua al menos $\Omega(h)$ (por ser la distancia en el árbol entre hojas de subárboles distintos), y a distancia a lo sumo $O(ch)$ de la raíz. En $\mathbb{R}^d$, una bola de radio $R = O(ch)$ contiene a lo sumo $O((R/r)^d)$ puntos con separación mutua $r = \Omega(h)$. Para acomodar $2^h$ puntos se necesita $(ch/h)^d \geq 2^h$, es decir, $d \geq \Omega(h) = \Omega(\log n)$.
\end{proof}

\begin{definition}[Función de Pérdida de Poincaré]\label{def:perdida_poincare}
Dado un conjunto de relaciones jerárquicas observadas $\mathcal{D} = \{(u, v)\}$ donde $u$ es ancestro de $v$ en la jerarquía, la \textit{función de pérdida de Poincaré} \citep{nickel2017poincare} se define como:
\begin{equation}
    \mathcal{L}(\boldsymbol{\theta}) = \sum_{(u,v) \in \mathcal{D}} \log \frac{e^{-d_{\mathbb{D}}(\Phi(u), \Phi(v))}}{\displaystyle\sum_{v' \in \mathcal{N}(u)} e^{-d_{\mathbb{D}}(\Phi(u), \Phi(v'))}}
    \label{eq:perdida_poincare}
\end{equation}
donde $\boldsymbol{\theta} = \{\Phi(e)\}_{e \in \mathcal{E}}$ son los parámetros del embedding y $\mathcal{N}(u) = \{v\} \cup \operatorname{Neg}(u)$ incluye el ejemplo positivo $v$ junto con un conjunto de ejemplos negativos muestreados $\operatorname{Neg}(u)$.
\end{definition}

\begin{remark}[Optimización Riemanniana]\label{rem:optimizacion_riemanniana_poincare}
La minimización de $\mathcal{L}$ sobre el disco de Poincaré requiere descenso por gradiente riemanniano (RSGD). El gradiente riemanniano se obtiene reescalando el gradiente euclidiano por el inverso del tensor métrico:
\begin{equation}
    \nabla_{\mathbb{D}} \mathcal{L} = \lambda_{\mathbf{x}}^{-2}\, \nabla_{\mathbb{R}} \mathcal{L}
\end{equation}
donde $\lambda_{\mathbf{x}} = \frac{2}{1 - \|\mathbf{x}\|^2}$ es el factor conformacional (\autoref{def:operaciones_hiperbolicas}). Este reescalamiento es necesario porque la métrica de Poincaré $g_{\mathbf{x}} = \lambda_{\mathbf{x}}^2 \mathbf{I}$ no es plana: un paso euclidiano fijo corresponde a distancias hiperbólicas cada vez mayores conforme $\|\mathbf{x}\| \to 1$. El \autoref{thm:hvt_convergence} formaliza la convergencia de este esquema de optimización.
\end{remark}

La justificación geométrica desarrollada en esta sección motiva una pregunta natural: si los datos poseen geometría intrínsecamente hiperbólica, ¿no debería el propio mecanismo de atención operar en ese espacio? Las subsecciones siguientes responden afirmativamente, definiendo redes de atención hiperbólica y transformers sobre el disco de Poincaré.

\begin{figure}[ht]
    \centering
    \includegraphics[width=0.9\textwidth]{figures/hyperbolic_embedding.pdf}
    \caption{Comparación entre embedding euclidiano (izquierda) y embedding en el disco de Poincaré (derecha) de un árbol jerárquico. En $\mathbb{R}^2$, las hojas del árbol se comprimen y las distancias se distorsionan. En el disco de Poincaré $\mathbb{D}^2$, el crecimiento exponencial del volumen hiperbólico permite separar las ramas con distorsión mínima (\autoref{thm:embedding_arboles_hiperbolico}), mientras que en $\mathbb{R}^d$ se necesita $d = \Omega(\log n)$ (\autoref{prop:cota_inferior_euclidiano}).}
    \label{fig:hyperbolic_embedding_attention}
\end{figure}

\subsubsection{Hyperbolic Attention Networks}

Los HANs de \citet{gulcehre2018hyperbolic} re-expresan la atención en términos de operaciones hiperbólicas, preservando las propiedades geométricas del espacio.

\begin{definition}[Distancia Hiperbólica]\label{def:distancia_hiperbolica}
En el modelo del hiperboloide, la distancia hiperbólica entre dos puntos $\mathbf{x}, \mathbf{y} \in \mathbb{H}^n$ se define como:

\begin{equation}
    d_{\mathbb{H}}(\mathbf{x}, \mathbf{y}) = \text{arccosh}(-\langle \mathbf{x}, \mathbf{y} \rangle_{\mathcal{L}})
    \label{eq:hyperbolic_distance}
\end{equation}

En el modelo de Poincaré, la distancia equivalente es:
\begin{equation}
d_{\mathbb{D}}(\mathbf{u}, \mathbf{v}) = \text{arccosh}\left(1 + 2\frac{\|\mathbf{u} - \mathbf{v}\|^2}{(1-\|\mathbf{u}\|^2)(1-\|\mathbf{v}\|^2)}\right)
\end{equation}
\end{definition}

\begin{theorem}[Propiedades de la Distancia Hiperbólica]\label{thm:propiedades_distancia_hiperbolica}
La distancia hiperbólica satisface:
\begin{enumerate}
    \item \textbf{Métrica}: $d_{\mathbb{H}}(\mathbf{x}, \mathbf{y}) \geq 0$ con igualdad si y sólo si $\mathbf{x} = \mathbf{y}$
    \item \textbf{Simetría}: $d_{\mathbb{H}}(\mathbf{x}, \mathbf{y}) = d_{\mathbb{H}}(\mathbf{y}, \mathbf{x})$
    \item \textbf{Desigualdad triangular}: $d_{\mathbb{H}}(\mathbf{x}, \mathbf{z}) \leq d_{\mathbb{H}}(\mathbf{x}, \mathbf{y}) + d_{\mathbb{H}}(\mathbf{y}, \mathbf{z})$
    \item \textbf{Invarianza isométrica}: $d_{\mathbb{H}}(g\mathbf{x}, g\mathbf{y}) = d_{\mathbb{H}}(\mathbf{x}, \mathbf{y})$ para $g \in SO^+(1,n)$
\end{enumerate}
\end{theorem}

\begin{proof}
Las propiedades (1)-(3) se verifican directamente desde la fórmula de distancia. La positividad y la igualdad resultan de las propiedades de arccosh. La simetría es inmediata por la simetría de la forma bilineal de Minkowski $\langle \mathbf{x}, \mathbf{y} \rangle_M$. La desigualdad triangular se demuestra usando que arccosh es convexa y las geodésicas del hiperboloide son secciones por planos que pasan por el origen. Para (4), las isometrías del hiperboloide preservan la forma de Minkowski: $\langle g\mathbf{x}, g\mathbf{y} \rangle_M = \langle \mathbf{x}, \mathbf{y} \rangle_M$ para $g \in SO^+(1,n)$, y la distancia depende solo de $\langle \mathbf{x}, \mathbf{y} \rangle_M$.
\end{proof}

\begin{definition}[Atención Hiperbólica Generalizada]\label{def:atencion_hiperbolica}
La atención hiperbólica se define como:

\begin{equation}
    \text{HypAttn}(\mathbf{q}, \mathbf{K}, \mathbf{V}) = \text{softmax}(f(d_{\mathbb{H}}(\mathbf{q}, \mathbf{K}))) \odot_{\mathbb{H}} \mathbf{V}
    \label{eq:hyperbolic_attention}
\end{equation}

donde:
\begin{itemize}
    \item $f: \mathbb{R} \to \mathbb{R}$ es una función de activación (típicamente $f(x) = -\alpha x$ con $\alpha > 0$)
    \item $\odot_{\mathbb{H}}$ es la agregación hiperbólica (generalización de suma ponderada)
    \item $\mathbf{q} \in \mathbb{H}^d$ es la query hiperbólica
    \item $\mathbf{K} = \{\mathbf{k}_i\}_{i=1}^n$ son las keys hiperbólicas
    \item $\mathbf{V} = \{\mathbf{v}_i\}_{i=1}^n$ son los values hiperbólicos
\end{itemize}
\end{definition}

\begin{definition}[Agregación Hiperbólica Ponderada]\label{def:agregacion_hiperbolica}
La agregación hiperbólica de valores $\{\mathbf{v}_i\}$ con pesos $\{w_i\}$ se define como:

\textbf{Einstein midpoint} (generalización natural en el modelo de Klein):
\begin{equation}
\bigoplus_{i=1}^n w_i \mathbf{v}_i = \frac{\sum_{i=1}^n w_i \gamma(\mathbf{v}_i) \mathbf{v}_i}{\sum_{i=1}^n w_i \gamma(\mathbf{v}_i)}
\end{equation}
donde $\gamma(\mathbf{v}) = \frac{1}{\sqrt{1-\|\mathbf{v}\|^2}}$ es el factor de Lorentz y el denominador es una suma escalar (no una norma vectorial).

\textbf{Promedio de Fréchet} (minimizador de varianza geométrica):
\begin{equation}
\bigoplus_{i=1}^n w_i \mathbf{v}_i = \arg\min_{\mathbf{x} \in \mathbb{H}^d} \sum_{i=1}^n w_i d_{\mathbb{H}}^2(\mathbf{x}, \mathbf{v}_i)
\end{equation}
\end{definition}

\textbf{Algoritmo para computar el promedio de Fréchet}:

\begin{algorithm}[H]
\caption{Promedio de Fréchet Hiperbólico}
\begin{algorithmic}[1]
    \State \textbf{Entrada:} Puntos $\{\mathbf{v}_i\}_{i=1}^n \subset \mathbb{H}^d$, pesos $\{w_i\}_{i=1}^n$
    \State Inicializar $\mathbf{m}^{(0)} = \mathbf{v}_1$
    \For{$t = 0, 1, 2, \ldots$ hasta convergencia}
        \State $\mathbf{g}^{(t)} = \sum_{i=1}^n w_i \log_{\mathbf{m}^{(t)}}^{\mathbb{H}}(\mathbf{v}_i)$
        \State $\mathbf{m}^{(t+1)} = \exp_{\mathbf{m}^{(t)}}^{\mathbb{H}}(\eta \mathbf{g}^{(t)})$
    \EndFor
    \State \textbf{Salida:} $\mathbf{m}^{(\infty)}$
\end{algorithmic}
\end{algorithm}

\textbf{Propiedades teóricas de la atención hiperbólica}:

\begin{proposition}[Equivarianza bajo Isometrías]\label{prop:equivarianza_atencion_hiperbolica}
Para cualquier isometría $g \in SO^+(1,n)$ del espacio hiperbólico, si la agregación hiperbólica $\odot_{\mathbb{H}}$ es equivariante bajo isometrías (como el promedio de Fréchet):
\begin{equation}
\text{HypAttn}(g\mathbf{q}, g\mathbf{K}, g\mathbf{V}) = g \cdot \text{HypAttn}(\mathbf{q}, \mathbf{K}, \mathbf{V})
\end{equation}
\end{proposition}

\begin{proof}[Idea de la demostración]
Los pesos de atención dependen de distancias hiperbólicas: $\alpha_{ij} \propto \exp(f(d_{\mathbb{H}}(\mathbf{q}, \mathbf{k}_j)))$. Como $g$ es una isometría, $d_{\mathbb{H}}(g\mathbf{q}, g\mathbf{k}_j) = d_{\mathbb{H}}(\mathbf{q}, \mathbf{k}_j)$ (\autoref{thm:propiedades_distancia_hiperbolica}, propiedad 4), por lo que los pesos de atención son invariantes: $\alpha_{ij}(g\mathbf{q}, g\mathbf{K}) = \alpha_{ij}(\mathbf{q}, \mathbf{K})$. Para el promedio de Fréchet, la equivarianza se sigue de que las isometrías preservan la función objetivo $\sum_i w_i d_{\mathbb{H}}^2(g\mathbf{x}, g\mathbf{v}_i) = \sum_i w_i d_{\mathbb{H}}^2(\mathbf{x}, \mathbf{v}_i)$, por lo que $g$ mapea el minimizador a $g$ del minimizador.
\end{proof}

\begin{corollary}[Invarianza de Similaridad]\label{cor:invarianza_similaridad}
La atención hiperbólica es invariante bajo transformaciones que preservan las relaciones jerárquicas en los datos.
\end{corollary}

\begin{proof}
Las transformaciones que preservan relaciones jerárquicas son precisamente las isometrías del espacio hiperbólico, es decir, el grupo $SO^+(1,n)$: estas preservan la distancia hiperbólica $d_{\mathbb{H}}$ y, con ella, las relaciones de proximidad que codifican la jerarquía. Por el \autoref{prop:equivarianza_atencion_hiperbolica}, la atención hiperbólica es equivariante bajo toda isometría $g \in SO^+(1,n)$. En particular, los pesos de atención $\alpha_{ij} \propto \exp(-d_{\mathbb{H}}(\mathbf{q}_i, \mathbf{k}_j))$ son invariantes bajo isometrías (pues $d_{\mathbb{H}}(g\mathbf{q}_i, g\mathbf{k}_j) = d_{\mathbb{H}}(\mathbf{q}_i, \mathbf{k}_j)$), lo que implica que las similaridades computadas por el mecanismo de atención son invariantes bajo transformaciones que preservan la estructura jerárquica.
\end{proof}

\textbf{Ventajas computacionales específicas}:

\begin{itemize}
    \item \textbf{Representación compacta}: Jerarquías complejas requieren dimensiones menores que embeddings euclidianos
    \item \textbf{Regularización implícita}: La geometría hiperbólica previene overfitting en datos jerárquicos
    \item \textbf{Interpretabilidad}: La distancia hiperbólica refleja directamente relaciones jerárquicas
    \item \textbf{Escalabilidad}: Operaciones eficientes en el modelo de Poincaré
\end{itemize}

\subsubsection{Hyperbolic Vision \gls{Transformer}}

El Hyperbolic Vision Transformer (HVT) de \citet{hvt2024hyperbolic} representa una extensión fundamental de la arquitectura \gls{ViT} que integra geometría hiperbólica para capturar eficientemente estructuras jerárquicas en datos visuales.

\textbf{Arquitectura fundamental de HVT}:

La arquitectura HVT consta de tres componentes principales que operan en espacio hiperbólico:

\begin{definition}[Proyecciones Hiperbólicas en HVT]\label{def:hvt_projections}
Las proyecciones query, key, value en HVT se reformulan mediante mapas exponencial/logarítmico:

\begin{align}
    \mathbf{Q} &= \text{exp}_{\mathbf{0}}^{\mathbb{H}}(\mathbf{W}_Q \text{log}_{\mathbf{0}}^{\mathbb{H}}(\mathbf{X})) \\
    \mathbf{K} &= \text{exp}_{\mathbf{0}}^{\mathbb{H}}(\mathbf{W}_K \text{log}_{\mathbf{0}}^{\mathbb{H}}(\mathbf{X})) \\
    \mathbf{V} &= \text{exp}_{\mathbf{0}}^{\mathbb{H}}(\mathbf{W}_V \text{log}_{\mathbf{0}}^{\mathbb{H}}(\mathbf{X}))
    \label{eq:hyperbolic_projections}
\end{align}

donde $\text{exp}_{\mathbf{0}}^{\mathbb{H}}: T_{\mathbf{0}}\mathbb{H}^d \to \mathbb{H}^d$ y $\text{log}_{\mathbf{0}}^{\mathbb{H}}: \mathbb{H}^d \to T_{\mathbf{0}}\mathbb{H}^d$ son los mapas exponencial y logarítmico hiperbólicos en el origen.
\end{definition}

\textbf{Justificación geométrica.}\label{rem:hvt_geometry} El mapeo $\text{log}_{\mathbf{0}}^{\mathbb{H}}$ proyecta representaciones hiperbólicas al espacio tangente euclidiano en el origen, donde se aplican transformaciones lineales estándar $\mathbf{W}_Q, \mathbf{W}_K, \mathbf{W}_V$. El mapa $\text{exp}_{\mathbf{0}}^{\mathbb{H}}$ posteriormente proyecta el resultado de vuelta al espacio hiperbólico, preservando propiedades geométricas.

\textbf{Formulación completa del mecanismo de atención hiperbólica}:

\begin{definition}[Atención Multi-Cabeza Hiperbólica]\label{def:hvt_multihead}
El mecanismo de atención multi-cabeza en HVT se define como:

\begin{equation}
\text{HypMultiHead}(\mathbf{X}) = \text{exp}_{\mathbf{0}}^{\mathbb{H}}\left(\mathbf{W}_O \text{Concat}(\text{head}_1, \ldots, \text{head}_h)\right)
\label{eq:hvt_multihead}
\end{equation}

donde cada cabeza $\text{head}_i$ opera en espacio hiperbólico:

\begin{equation}
\text{head}_i = \text{softmax}\left(-\alpha d_{\mathbb{H}}(\mathbf{Q}_i, \mathbf{K}_i)\right) \odot_{\mathbb{H}} \mathbf{V}_i
\label{eq:hvt_head}
\end{equation}

con $\alpha > 0$ un parámetro de escala y $\odot_{\mathbb{H}}$ la agregación hiperbólica (promedio de Fréchet).
\end{definition}

\textbf{Procedimientos de entrenamiento y estabilidad numérica}:

El entrenamiento de HVT requiere consideraciones especiales para mantener estabilidad numérica en espacio hiperbólico:

\begin{enumerate}
    \item \textbf{Inicialización hiperbólica}: Los parámetros se inicializan mediante proyección exponencial de distribuciones gaussianas pequeñas:
    \begin{equation}
    \mathbf{w}^{(0)} = \text{exp}_{\mathbf{0}}^{\mathbb{H}}(\epsilon \mathbf{z}), \quad \mathbf{z} \sim \mathcal{N}(0, \mathbf{I}), \quad \epsilon \ll 1
    \end{equation}

    \item \textbf{Proyección a la bola de Poincaré}: Para mantener estabilidad numérica, todos los puntos se proyectan a $\|\mathbf{x}\| < 1 - \delta$ con $\delta = 10^{-5}$:
    \begin{equation}
    \text{proj}(\mathbf{x}) = \begin{cases}
    \mathbf{x} & \text{si } \|\mathbf{x}\| \leq 1 - \delta \\
    \frac{(1-\delta)\mathbf{x}}{\|\mathbf{x}\|} & \text{en caso contrario}
    \end{cases}
    \end{equation}

    \item \textbf{Optimización Riemanniana}: Se utiliza el optimizador RSGD (Riemannian Stochastic Gradient Descent):
    \begin{equation}
    \mathbf{x}^{(t+1)} = \text{exp}_{\mathbf{x}^{(t)}}^{\mathbb{H}}\left(-\eta \text{grad}^{\mathbb{H}} L(\mathbf{x}^{(t)})\right)
    \end{equation}
    donde $\text{grad}^{\mathbb{H}} L = \lambda_{\mathbf{x}}^{-2} \nabla L$ es el gradiente Riemanniano con $\lambda_{\mathbf{x}} = \frac{2}{1-\|\mathbf{x}\|^2}$.
\end{enumerate}

\begin{theorem}[Convergencia de RSGD en Variedades]\label{thm:hvt_convergence}
Sea $L: \mathbb{H}^d \to \mathbb{R}$ una función geodésicamente convexa con gradiente Riemanniano $M$-Lipschitz. El algoritmo RSGD con tasa de aprendizaje $\eta < 1/M$ satisface, para el iterado promediado $\bar{\mathbf{x}}^{(T)} = \frac{1}{T}\sum_{t=0}^{T-1} \mathbf{x}^{(t)}$ (promedio en el espacio tangente):
\begin{equation}
\mathbb{E}[L(\bar{\mathbf{x}}^{(T)})] - L(\mathbf{x}^*) \leq \frac{d_{\mathbb{H}}^2(\mathbf{x}^{(0)}, \mathbf{x}^*)}{2\eta T} + \frac{\eta M \sigma^2}{2}
\end{equation}
donde $\sigma^2$ es la varianza del gradiente estocástico y $\mathbf{x}^*$ es un minimizador global. Este resultado es una adaptación del teorema de convergencia de \citet{bonnabel2013stochastic} para descenso por gradiente estocástico en variedades Riemannianas.
\end{theorem}

\textbf{Resultados experimentales y ablación}:

Los experimentos de \citet{hvt2024hyperbolic} demuestran ventajas significativas de HVT sobre \gls{ViT} estándar:

\begin{itemize}
    \item \textbf{ImageNet-1K}: HVT-Base alcanza 84.2\% top-1 accuracy con 30\% menos parámetros que ViT-Base
    \item \textbf{CIFAR-100}: Mejora de 2.8\% en accuracy con embeddings hiperbólicos de dimensión $d=256$ vs. euclidianos de $d=512$
    \item \textbf{Datos jerárquicos}: En clasificación de especies biológicas (estructura taxonómica), HVT supera a \gls{ViT} por 5.3\%
\end{itemize}

\subsubsection{HyPE-GT: Codificaciones Posicionales Hiperbólicas}

HyPE-GT (Hyperbolic Positional Encodings for Graph Transformers) de \citet{hypegt2023} introduce codificaciones posicionales aprendibles en espacio hiperbólico que capturan naturalmente estructuras jerárquicas en grafos.

\textbf{Formulación matemática de HyPE}:

\begin{definition}[Codificaciones Posicionales Hiperbólicas]\label{def:hype_encoding}
Para un grafo $\mathcal{G} = (V, E)$, las codificaciones posicionales hiperbólicas se definen como:

\begin{equation}
    \mathbf{PE}_i^{\mathbb{H}} = \text{HypMLP}(\mathbf{h}_i^{(0)}, \mathbf{s}_i)
    \label{eq:hyperbolic_positional_encoding}
\end{equation}

donde:
\begin{itemize}
    \item $\mathbf{h}_i^{(0)} \in \mathbb{R}^{d_0}$ son las características iniciales del nodo $i$
    \item $\mathbf{s}_i \in \mathbb{R}^{d_s}$ es el vector estructural que codifica propiedades topológicas del nodo
    \item $\text{HypMLP}: \mathbb{R}^{d_0 + d_s} \to \mathbb{H}^d$ es una red neuronal multi-capa con salida hiperbólica
\end{itemize}
\end{definition}

\begin{definition}[Vector Estructural de Nodo]\label{def:structural_vector}
El vector estructural $\mathbf{s}_i$ captura propiedades topológicas locales y globales:

\begin{equation}
\mathbf{s}_i = [\text{deg}(i), \text{CC}(i), \text{BC}(i), \text{PR}(i), \mathbf{d}_i^{\text{ecc}}]
\end{equation}

donde:
\begin{itemize}
    \item $\text{deg}(i)$ es el grado del nodo
    \item $\text{CC}(i)$ es el coeficiente de clustering
    \item $\text{BC}(i)$ es la centralidad de intermediación (betweenness centrality)
    \item $\text{PR}(i)$ es el PageRank
    \item $\mathbf{d}_i^{\text{ecc}} \in \mathbb{R}^k$ es el vector de excentricidades (distancias a nodos representativos)
\end{itemize}
\end{definition}

Las codificaciones HyPE se integran en la arquitectura \gls{Transformer} mediante combinación hiperbólica con las operaciones de Möbius, y la atención se computa sobre los embeddings combinados.

\begin{remark}[Ventaja dimensional de embeddings hiperbólicos]\label{rem:hypegt_expressiveness}
Las codificaciones posicionales hiperbólicas de HyPE-GT \citep{hypegt2023} se benefician de la propiedad general de que los embeddings hiperbólicos requieren dimensiones menores que los euclidianos para representar estructuras jerárquicas. Por el \autoref{thm:embedding_arboles_hiperbolico}, un árbol con $n$ vértices se puede embeber en $\mathbb{H}^2$ con distorsión $(1+\varepsilon)$, mientras que embeddings euclidianos con distorsión acotada requieren dimensión $d = \Omega(\log n)$ (\autoref{prop:cota_inferior_euclidiano}). Esta ventaja teórica motiva el uso de codificaciones hiperbólicas en HyPE-GT, aunque \citet{hypegt2023} no establecen formalmente cotas de expresividad para su arquitectura, siendo su contribución principalmente empírica.
\end{remark}

\subsection{Atención y Embeddings en la Esfera}\label{subsec:atencion_esfera}

Mientras que la sección anterior explora la atención en espacios de curvatura negativa, la esfera unitaria $S^{d-1}$ constituye el espacio de curvatura positiva constante más relevante para mecanismos de atención. Normalizar embeddings a $S^{d-1}$ induce una conexión directa entre el producto punto (base de la atención estándar) y la distancia geodésica esférica.

\begin{proposition}[Equivalencia Producto Punto--Similaridad Coseno en $S^{d-1}$]\label{prop:dot_product_sphere}
Sean $\hat{\mathbf{q}}, \hat{\mathbf{k}} \in S^{d-1}$ embeddings normalizados, es decir, $\hat{\mathbf{q}} = \mathbf{q}/\|\mathbf{q}\|$ y $\hat{\mathbf{k}} = \mathbf{k}/\|\mathbf{k}\|$. Entonces:
\begin{equation}
\hat{\mathbf{q}}^T \hat{\mathbf{k}} = \cos\bigl(d_{S^{d-1}}(\hat{\mathbf{q}}, \hat{\mathbf{k}})\bigr)
\label{eq:dot_geodesic_sphere}
\end{equation}
donde $d_{S^{d-1}}(\hat{\mathbf{q}}, \hat{\mathbf{k}}) = \arccos(\hat{\mathbf{q}}^T \hat{\mathbf{k}})$ es la distancia geodésica en $S^{d-1}$. En consecuencia, la atención escalada sobre embeddings normalizados
\[
\alpha_{ij} = \frac{\exp\bigl(\hat{\mathbf{q}}_i^T \hat{\mathbf{k}}_j / \tau\bigr)}{\sum_l \exp\bigl(\hat{\mathbf{q}}_i^T \hat{\mathbf{k}}_l / \tau\bigr)}
\]
opera exclusivamente con distancias geodésicas esféricas (cf.\ \autoref{def:dot_product_similarity}).
\end{proposition}

\begin{proof}
Por definición de la métrica en $S^{d-1}$ (véase \autoref{def:esfera_n}), la distancia geodésica entre $\hat{\mathbf{q}}, \hat{\mathbf{k}} \in S^{d-1}$ es $d_{S^{d-1}}(\hat{\mathbf{q}}, \hat{\mathbf{k}}) = \arccos(\hat{\mathbf{q}}^T \hat{\mathbf{k}})$. Aplicando coseno a ambos lados se obtiene directamente la identidad~\eqref{eq:dot_geodesic_sphere}. Puesto que $\|\hat{\mathbf{q}}\| = \|\hat{\mathbf{k}}\| = 1$, el producto punto $\hat{\mathbf{q}}^T \hat{\mathbf{k}}$ no depende de magnitudes, sino únicamente del ángulo geodésico, y los pesos de atención $\alpha_{ij}$ son función monótona decreciente de la distancia geodésica $d_{S^{d-1}}$.
\end{proof}

La \autoref{prop:dot_product_sphere} motiva el uso de embeddings normalizados en tareas de clasificación: al eliminar la magnitud, la frontera de decisión depende exclusivamente de ángulos.

\begin{definition}[Clasificador Angular]\label{def:clasificador_angular}
Sea $\hat{\mathbf{f}}(\mathbf{x}) = \mathbf{f}(\mathbf{x})/\|\mathbf{f}(\mathbf{x})\| \in S^{d-1}$ el embedding normalizado de una entrada $\mathbf{x}$, y sean $\hat{\mathbf{w}}_c = \mathbf{w}_c/\|\mathbf{w}_c\| \in S^{d-1}$ los pesos normalizados para cada clase $c \in \{1, \ldots, C\}$. Un \textit{clasificador angular} \citep{wang2017normface} produce logits de la forma:
\begin{equation}
z_c(\mathbf{x}) = s \cdot \hat{\mathbf{f}}(\mathbf{x})^T \hat{\mathbf{w}}_c = s \cos\theta_c
\label{eq:clasificador_angular}
\end{equation}
donde $\theta_c = d_{S^{d-1}}(\hat{\mathbf{f}}(\mathbf{x}), \hat{\mathbf{w}}_c) \in [0, \pi]$ es la distancia geodésica al prototipo de clase $c$, y $s > 0$ es un parámetro de escala necesario para controlar la concentración de la distribución softmax (sin $s$, los logits estarían acotados en $[-1, 1]$, produciendo distribuciones de probabilidad excesivamente uniformes).
\end{definition}

\begin{definition}[Pérdida con Margen Angular Generalizado]\label{def:perdida_margen_angular}
Dada la estructura de clasificador angular (\autoref{def:clasificador_angular}), una \textit{pérdida con margen angular} modifica el logit de la clase correcta $y$ mediante una función de margen $g: [0, \pi] \to \mathbb{R}$:
\begin{equation}
\mathcal{L} = -\log \frac{\exp\bigl(s \cdot g(\theta_y)\bigr)}{\exp\bigl(s \cdot g(\theta_y)\bigr) + \sum_{c \neq y} \exp\bigl(s \cos\theta_c\bigr)}
\label{eq:perdida_margen_angular}
\end{equation}
Las principales instancias de esta familia son:
\begin{enumerate}
    \item \textbf{NormFace} \citep{wang2017normface}: $g(\theta) = \cos\theta$ (sin margen, baseline normalizado).
    \item \textbf{SphereFace} \citep{liu2017sphereface}: $g(\theta) = \cos(m\theta)$ con $m \in \mathbb{N}$, $m \geq 2$ (margen multiplicativo angular).
    \item \textbf{CosFace} \citep{wang2018cosface}: $g(\theta) = \cos\theta - m$ con $m > 0$ (margen aditivo en el espacio del coseno).
    \item \textbf{ArcFace} \citep{deng2019arcface}: $g(\theta) = \cos(\theta + m)$ con $m > 0$ (margen aditivo angular).
\end{enumerate}
\end{definition}

\begin{proposition}[Interpretación Geodésica del Margen Angular]\label{prop:interpretacion_geodesica_margen}
En ArcFace, el margen $m > 0$ impone una separación geodésica de exactamente $m$ radianes en $S^{d-1}$ entre el embedding y la frontera de decisión más cercana.
\end{proposition}

\begin{proof}
Sin margen, la frontera de decisión entre clases $y$ y $c$ se sitúa donde $\cos\theta_y = \cos\theta_c$, es decir, donde las distancias geodésicas a ambos prototipos coinciden. Con ArcFace, la condición de clasificación correcta requiere $\cos(\theta_y + m) \geq \cos\theta_c$, lo que equivale a $\theta_y + m \leq \theta_c$ (en el rango $[0, \pi]$ donde el coseno es monótonamente decreciente). En la frontera, $\theta_y + m = \theta_c$, y por tanto el embedding debe estar a distancia geodésica al menos $m$ radianes más cerca de su prototipo de clase que de cualquier otro. Como $\theta_y = d_{S^{d-1}}(\hat{\mathbf{f}}(\mathbf{x}), \hat{\mathbf{w}}_y)$, esta separación es intrínsecamente geodésica.
\end{proof}

\begin{remark}[Aprendizaje Contrastivo como Atención en $S^{d-1}$]\label{rem:contrastive_sphere}
La pérdida InfoNCE utilizada en modelos multimodales constituye un caso particular de clasificación angular. Dados embeddings normalizados $\hat{\mathbf{z}}_i^A, \hat{\mathbf{z}}_j^B \in S^{d-1}$ de dos modalidades, la pérdida
\[
\mathcal{L}_{\text{InfoNCE}} = -\log \frac{\exp\bigl(\hat{\mathbf{z}}_i^{A\,T} \hat{\mathbf{z}}_i^B / \tau\bigr)}{\sum_{j} \exp\bigl(\hat{\mathbf{z}}_i^{A\,T} \hat{\mathbf{z}}_j^B / \tau\bigr)}
\]
es exactamente un clasificador angular con $s = 1/\tau$, $g(\theta) = \cos\theta$, y las clases definidas dinámicamente por el batch. Así, el aprendizaje contrastivo realiza clasificación angular sobre la esfera, donde la temperatura $\tau$ controla la concentración de la atención softmax sobre los vecinos geodésicos más cercanos.
\end{remark}

\subsection{Transporte Paralelo y Conexiones}

El transporte paralelo proporciona un mecanismo fundamental para comparar vectores tangentes en diferentes puntos de una variedad Riemanniana, preservando la estructura geométrica intrínseca. Este concepto es esencial para formular mecanismos de atención en variedades no-euclidianas.

\subsubsection{Fundamentos Matemáticos del Transporte Paralelo}

Recordamos que la \textit{conexión de Levi-Civita} $\nabla$ de una variedad riemanniana $(M, g)$ es la única conexión afín compatible con la métrica y libre de torsión (\autoref{def:conexion_levi_civita_gdl} del \autoref{ch:gdl}), determinada en coordenadas locales por los símbolos de Christoffel $\Gamma^k_{ij}$.\label{def:levi_civita_attention}\label{eq:christoffel_symbols_attention}

El \textit{transporte paralelo} a lo largo de una curva $\gamma: [0,1] \to M$ viene dado por la ecuación $\nabla_{\dot{\gamma}(t)} V(t) = 0$, que en coordenadas locales se traduce en $\frac{dV^k}{dt} + \Gamma^k_{ij} \frac{dx^i}{dt} V^j = 0$ (véase la \autoref{def:transporte_paralelo_formal} del \autoref{ch:gdl} para el caso general en fibrados principales).\label{def:parallel_transport_attention}\label{eq:parallel_transport}\label{eq:parallel_transport_coordinates}

\begin{theorem}[Existencia y Unicidad del Transporte Paralelo]\label{thm:existence_parallel_transport}
Dada una curva suave $\gamma: [0,1] \to M$ y un vector inicial $v_0 \in T_{\gamma(0)}M$, existe un único campo vectorial paralelo $V(t)$ a lo largo de $\gamma$ con $V(0) = v_0$.

El transporte paralelo define una isometría lineal:
\begin{equation}
\Gamma_\gamma: T_{\gamma(0)}M \to T_{\gamma(1)}M
\end{equation}
que satisface $\Gamma_\gamma(v_0) = V(1)$.
\end{theorem}

\begin{proof}
La ecuación~\eqref{eq:parallel_transport_coordinates} es un sistema lineal de EDOs de primer orden. Por el teorema de Picard-Lindelöf, dado un valor inicial $V(0) = v_0$, existe una única solución $V(t)$ definida para $t \in [0,1]$.

La preservación métrica se sigue de la compatibilidad de $\nabla$ con $g$:
\begin{align}
\frac{d}{dt}\langle V(t), W(t) \rangle &= \langle \nabla_{\dot{\gamma}} V, W \rangle + \langle V, \nabla_{\dot{\gamma}} W \rangle \\
&= 0 \quad \text{(si $V$ y $W$ son paralelos)}
\end{align}

Por tanto, $\langle \Gamma_\gamma(v), \Gamma_\gamma(w) \rangle_{\gamma(1)} = \langle v, w \rangle_{\gamma(0)}$ para todo $v, w \in T_{\gamma(0)}M$.
\end{proof}

\begin{proposition}[Propiedades del Transporte Paralelo]\label{prop:parallel_transport_properties}
El transporte paralelo $\Gamma_\gamma: T_{\gamma(0)}M \to T_{\gamma(1)}M$ satisface:

\begin{enumerate}
    \item \textbf{Linealidad}: $\Gamma_\gamma(\alpha u + \beta v) = \alpha \Gamma_\gamma(u) + \beta \Gamma_\gamma(v)$
    \item \textbf{Preservación métrica}: $g_{\gamma(1)}(\Gamma_\gamma(u), \Gamma_\gamma(v)) = g_{\gamma(0)}(u, v)$
    \item \textbf{Propiedad de concatenación}: Para caminos $\gamma_1, \gamma_2$ con $\gamma_1(1) = \gamma_2(0)$:
    \begin{equation}
    \Gamma_{\gamma_2 * \gamma_1} = \Gamma_{\gamma_2} \circ \Gamma_{\gamma_1}
    \end{equation}
    \item \textbf{Inversión}: $\Gamma_{\gamma^{-1}} = (\Gamma_\gamma)^{-1}$ donde $\gamma^{-1}(t) = \gamma(1-t)$
\end{enumerate}
\end{proposition}

\begin{proof}
\textbf{(1) Linealidad}: Dado que el transporte paralelo se define como la solución de la ecuación $\nabla_{\dot{\gamma}} V = 0$, que es lineal en $V$, la unicidad de soluciones de EDOs lineales garantiza la linealidad del operador $\Gamma_\gamma$.

\textbf{(2) Preservación métrica}: Si $\nabla$ es la conexión de Levi-Civita, compatible con la métrica, entonces para campos paralelos $V, W$ a lo largo de $\gamma$:
\[
\frac{d}{dt} g(V, W) = g(\nabla_{\dot{\gamma}} V, W) + g(V, \nabla_{\dot{\gamma}} W) = 0
\]
por lo que $g(V,W)$ es constante a lo largo de $\gamma$.

\textbf{(3) Concatenación}: Sea $V$ un campo paralelo a lo largo de $\gamma_1$ con $V(0) = v$. Entonces $V(1) = \Gamma_{\gamma_1}(v)$. Sea $W$ paralelo a lo largo de $\gamma_2$ con $W(0) = V(1)$. Entonces $W(1) = \Gamma_{\gamma_2}(\Gamma_{\gamma_1}(v))$. Pero la concatenación $\gamma_2 * \gamma_1$ con campo $V$ seguido de $W$ es paralelo a lo largo del camino completo.

\textbf{(4) Inversión}: Si $V$ es paralelo a lo largo de $\gamma$, entonces $V(\gamma^{-1}(t)) = V(\gamma(1-t))$ es paralelo a lo largo de $\gamma^{-1}$, por lo que $\Gamma_{\gamma^{-1}} \circ \Gamma_\gamma = \text{Id}$.
\end{proof}

\begin{figure}[ht]
    \centering
    \includegraphics[width=0.8\textwidth]{figures/parallel_transport_surface.pdf}
    \caption{Transporte paralelo de un vector tangente a lo largo de una curva $\gamma$ en una superficie. El vector $V(t)$ satisface $\nabla_{\dot{\gamma}} V = 0$, manteniendo su norma y ángulo relativo a la curva constantes. La dependencia del resultado en el camino elegido (holonomía) refleja la curvatura de la superficie y motiva la atención con transporte paralelo (\autoref{def:attention_parallel_transport}).}
    \label{fig:parallel_transport_attention}
\end{figure}

\subsubsection{Algoritmo para Computar Geodésicas y Transporte Paralelo}

En la práctica, el transporte paralelo se computa numéricamente mediante integración de la ecuación~\eqref{eq:parallel_transport_coordinates} a lo largo de geodésicas.

Recordamos que una \textit{geodésica} es una curva $\gamma$ que satisface $\nabla_{\dot{\gamma}} \dot{\gamma} = 0$, o equivalentemente en coordenadas locales $\ddot{x}^k + \Gamma^k_{ij} \dot{x}^i \dot{x}^j = 0$ (\autoref{def:geodesica} del \autoref{ch:gdl}).\label{def:geodesic_attention}\label{eq:geodesic_equation}

\begin{algorithm}[H]
\caption{Cálculo de Transporte Paralelo en Variedades Riemannianas}
\label{alg:parallel_transport_computation}
\begin{algorithmic}[1]
    \State \textbf{Entrada:} Punto inicial $p \in M$, punto final $q \in M$, vector $v \in T_p M$
    \State \textbf{Salida:} Vector transportado $\Gamma_{pq}(v) \in T_q M$
    \State
    \State \textbf{Paso 1: Computar geodésica} $\gamma: [0,1] \to M$ con $\gamma(0) = p$, $\gamma(1) = q$
    \State Resolver la ecuación geodésica~\eqref{eq:geodesic_equation} usando Runge-Kutta de orden 4
    \State
    \State \textbf{Paso 2: Inicializar} campo vectorial paralelo $V(0) = v$
    \State
    \State \textbf{Paso 3: Integrar} ecuación de transporte paralelo~\eqref{eq:parallel_transport_coordinates}
    \For{$t = 0$ to $1$ (discretización $\Delta t = 1/N$)}
        \State Computar símbolos de Christoffel $\Gamma^k_{ij}(\gamma(t))$ en el punto actual
        \State Actualizar $V(t + \Delta t)$ usando método Runge-Kutta:
        \State $\quad V^k(t + \Delta t) = V^k(t) - \Delta t \cdot \Gamma^k_{ij}(\gamma(t)) \dot{\gamma}^i(t) V^j(t)$
    \EndFor
    \State
    \State \textbf{Retornar} $\Gamma_{pq}(v) = V(1)$
\end{algorithmic}
\end{algorithm}

\subsubsection{Ejemplo Trabajado: Transporte Paralelo en la Esfera $S^2$}

Desarrollamos un ejemplo completo de transporte paralelo en la esfera unitaria, donde las fórmulas admiten expresiones cerradas.

\begin{example}[Transporte Paralelo en $S^2$]\label{ex:parallel_transport_sphere}
Consideremos la esfera unitaria $S^2 = \{\mathbf{x} \in \mathbb{R}^3 : \|\mathbf{x}\| = 1\}$ con la métrica inducida del espacio euclidiano.

\textbf{Geodésicas en $S^2$}: Las geodésicas son grandes círculos. Para dos puntos $p, q \in S^2$ no antípodas, la geodésica está dada por:
\begin{equation}
\gamma(t) = \frac{\sin((1-t)\theta)}{\sin\theta} p + \frac{\sin(t\theta)}{\sin\theta} q, \quad t \in [0,1]
\end{equation}
donde $\theta = \arccos(\langle p, q \rangle)$ es el ángulo entre $p$ y $q$.

\textbf{Fórmula explícita para transporte paralelo}: Para un vector tangente $v \in T_p S^2$ (con $\langle v, p \rangle = 0$), el transporte paralelo a $T_q S^2$ a lo largo de la geodésica $\gamma$ es:

\begin{equation}
\Gamma_{\gamma_{pq}}(v) = v - \frac{\langle v, q \rangle + \langle v, p \rangle}{1 + \langle p, q \rangle}(p + q)
\label{eq:parallel_transport_sphere_formula}
\end{equation}

Aprovechando que $v \perp p$ (es decir, $\langle v, p \rangle = 0$), esta fórmula se simplifica a:
\begin{equation}
\Gamma_{\gamma_{pq}}(v) = v - \frac{\langle v, q \rangle}{1 + \langle p, q \rangle}(p + q)
\end{equation}

\textbf{Verificación de propiedades}:
\begin{enumerate}
    \item \textbf{Tangencia}: $\langle \Gamma_{\gamma_{pq}}(v), q \rangle = 0$ (el vector transportado es tangente en $q$)
    \item \textbf{Preservación de norma}: $\|\Gamma_{\gamma_{pq}}(v)\| = \|v\|$
    \item \textbf{Preservación de ángulos}: Para $v, w \in T_p S^2$:
    \begin{equation}
    \langle \Gamma_{\gamma_{pq}}(v), \Gamma_{\gamma_{pq}}(w) \rangle = \langle v, w \rangle
    \end{equation}
\end{enumerate}
\end{example}

\subsubsection{Transporte Paralelo en Mecanismos de Atención}

El transporte paralelo permite formular mecanismos de atención geométricamente consistentes en variedades Riemannianas.

\begin{definition}[Atención con Transporte Paralelo]\label{def:attention_parallel_transport}
Sea $M$ una variedad Riemanniana y $\{x_i\}_{i=1}^n \subset M$ un conjunto de puntos (tokens) en la variedad. Para features $\{f_i \in T_{x_i}M\}_{i=1}^n$, la atención con transporte paralelo en el punto $x_j$ se define como:

\begin{equation}
\text{Attn}_{\text{PT}}(x_j, \{x_i, f_i\}) = \sum_{i=1}^n \alpha_{ij} \Gamma_{\gamma_{ij}}(f_i)
\label{eq:attention_parallel_transport}
\end{equation}

donde:
\begin{itemize}
    \item $\gamma_{ij}$ es la geodésica de $x_i$ a $x_j$
    \item $\Gamma_{\gamma_{ij}}: T_{x_i}M \to T_{x_j}M$ es el transporte paralelo
    \item $\alpha_{ij} = \text{softmax}_i(-d_M(x_i, x_j)^2)$ son los pesos de atención basados en distancia geodésica
\end{itemize}
\end{definition}

\begin{theorem}[Equivarianza bajo Isometrías]\label{thm:attention_parallel_transport_equivariance}
Sea $\phi: M \to M$ una isometría de la variedad Riemanniana. La atención con transporte paralelo satisface:

\begin{equation}
\text{Attn}_{\text{PT}}(\phi(x_j), \{\phi(x_i), d\phi_{x_i}(f_i)\}) = d\phi_{x_j}(\text{Attn}_{\text{PT}}(x_j, \{x_i, f_i\}))
\end{equation}

donde $d\phi_x: T_x M \to T_{\phi(x)}M$ es el diferencial de $\phi$.
\end{theorem}

\begin{proof}
Las isometrías preservan geodésicas y transporte paralelo:
\begin{enumerate}
    \item $\phi(\gamma_{ij}(t)) = \gamma_{\phi(i)\phi(j)}(t)$ (geodésicas se mapean a geodésicas)
    \item $d\phi \circ \Gamma_{\gamma_{ij}} = \Gamma_{\phi(\gamma_{ij})} \circ d\phi$ (conmutatividad con transporte)
    \item $d_M(\phi(x_i), \phi(x_j)) = d_M(x_i, x_j)$ (preservación de distancias)
\end{enumerate}

Aplicando estas propiedades:
\begin{align}
\text{Attn}_{\text{PT}}(\phi(x_j), \{\phi(x_i), d\phi(f_i)\}) &= \sum_i \alpha_{ij} \Gamma_{\gamma_{\phi(i)\phi(j)}}(d\phi(f_i)) \\
&= \sum_i \alpha_{ij} d\phi(\Gamma_{\gamma_{ij}}(f_i)) \\
&= d\phi\left(\sum_i \alpha_{ij} \Gamma_{\gamma_{ij}}(f_i)\right) \\
&= d\phi(\text{Attn}_{\text{PT}}(x_j, \{x_i, f_i\})) \qquad \square
\end{align}
\end{proof}

\subsubsection{Aplicaciones en Deep Learning Geométrico}

El transporte paralelo ha encontrado aplicaciones exitosas en arquitecturas de atención para datos no-euclidianos:

\begin{itemize}
    \item \textbf{Modelos climáticos en la esfera}: Atención en datos meteorológicos definidos en $S^2$ usando transporte paralelo para agregar información de diferentes localizaciones preservando orientaciones

    \item \textbf{Procesamiento de meshes 3D}: Convoluciones en superficies donde el transporte paralelo permite comparar features tangentes en diferentes vértices del mesh

    \item \textbf{Redes neuronales en espacios de formas}: Análisis de formas anatómicas representadas como puntos en variedades Riemannianas (e.g., variedad de Stiefel, Grassmannian)

    \item \textbf{Atención hiperbólica con transporte}: Combinación de geometría hiperbólica (para jerarquías) con transporte paralelo (para agregación geométricamente consistente)
\end{itemize}

\begin{example}[Implementación práctica: Atención en $SO(3)$]\label{ex:attention_so3}
Para datos de orientaciones representados como rotaciones $R_i \in SO(3)$, la atención con transporte paralelo permite agregar información respetando la geometría de Lie del grupo.

Dado un conjunto de orientaciones $\{R_i\}$ con features $\{f_i \in T_{R_i}SO(3) \cong \mathfrak{so}(3)\}$:

\begin{enumerate}
    \item \textbf{Computar geodésicas}: Para $R_i, R_j \in SO(3)$, la geodésica es:
    \begin{equation}
    \gamma(t) = R_i \exp(t \log(R_i^{-1} R_j)), \quad t \in [0,1]
    \end{equation}

    \item \textbf{Transporte paralelo}: En grupos de Lie, el transporte paralelo mediante traslación izquierda:
    \begin{equation}
    \Gamma_{\gamma_{ij}}(f_i) = \text{Ad}_{R_i^{-1}R_j}(f_i)
    \end{equation}
    donde $\text{Ad}_R: \mathfrak{so}(3) \to \mathfrak{so}(3)$ es la representación adjunta.

    \item \textbf{Agregación}: Compute pesos $\alpha_{ij} \propto \exp(-\|R_i^{-1}R_j\|_F^2)$ y agregar:
    \begin{equation}
    \text{Attn}(R_j) = \sum_i \alpha_{ij} \text{Ad}_{R_i^{-1}R_j}(f_i)
    \end{equation}
\end{enumerate}

Esta formulación es invariante bajo rotaciones globales y preserva la estructura de grupo de $SO(3)$.
\end{example}

\section{Álgebra Geométrica (Clifford)}\label{sec:clifford_attention}

El álgebra geométrica (GA) proporciona un framework matemático unificado y extraordinariamente poderoso para representar y manipular objetos geométricos, extendiendo naturalmente los mecanismos de atención a estructuras algebraicas más ricas. Este formalismo, desarrollado originalmente por William Kingdon Clifford en el siglo XIX, ha experimentado un renacimiento en el contexto del deep learning geométrico, ofreciendo una alternativa elegante y computacionalmente eficiente para el manejo de simetrías complejas.

\subsection{Fundamentos Matemáticos de Álgebras de Clifford}

La idea central del álgebra de Clifford es extender un espacio vectorial con un producto asociativo (el \textit{producto geométrico}) que codifica simultáneamente las relaciones métricas (producto interior) y las orientaciones (producto exterior). Esto permite representar escalares, vectores, planos y volúmenes como elementos de una única estructura algebraica, facilitando la construcción de operaciones equivariantes que actúan de manera uniforme sobre todos estos objetos.

\subsubsection{Definición y Propiedades Básicas}

Recordamos que el álgebra de Clifford $\text{Cl}(p,q)$ (\autoref{def:clifford_cnn} del \autoref{ch:cnn}) es el álgebra asociativa generada por vectores ortonormales que satisfacen $\mathbf{e}_i \mathbf{e}_j + \mathbf{e}_j \mathbf{e}_i = 2\eta_{ij}$, con dimensión $2^{p+q}$ como espacio vectorial.\label{def:algebra_clifford}

\begin{theorem}[Periodicidad de Bott para Álgebras de Clifford]\label{thm:clasificacion_clifford}
La clasificación de las álgebras de Clifford reales exhibe periodicidad de período 8: $\text{Cl}(p+8, q) \cong \text{Cl}(p, q) \otimes M_{16}(\mathbb{R})$ y $\text{Cl}(p, q+8) \cong \text{Cl}(p, q) \otimes M_{16}(\mathbb{R})$. Para dimensión fija $n = p + q$, la signatura $p - q$ determina unívocamente la pareja $(p,q)$, y por tanto el tipo de isomorfismo del álgebra. La estructura (si el álgebra es isomorfa a una álgebra de matrices sobre $\mathbb{R}$, $\mathbb{C}$ o $\mathbb{H}$, o un producto directo de dos tales) depende únicamente de $p - q \pmod{8}$. Para la demostración, véase \citep{lawson1989spin}.
\end{theorem}

\subsubsection{Multivectores y Gradación}

Los multivectores $M = \sum_k \langle M \rangle_k$ admiten descomposición por grados, donde $\langle M \rangle_0$ es escalar, $\langle M \rangle_1$ es vector, $\langle M \rangle_2$ es bivector, etc.\label{def:multivector}

\subsubsection{Productos Geométricos Fundamentales}

El álgebra de Clifford admite múltiples productos que capturan diferentes aspectos geométricos:

El álgebra de Clifford admite tres productos fundamentales: el producto geométrico $AB = A \cdot B + A \wedge B$, el producto interior (contracción) y el producto exterior (wedge), definidos formalmente en la \autoref{def:producto_geometrico_cnn} del \autoref{ch:cnn}.\label{def:productos_clifford}

\begin{proposition}[Propiedades de los Productos]\label{prop:propiedades_productos_clifford}
Los productos de Clifford satisfacen:
\begin{enumerate}
    \item \textbf{Distributividad}: $(A + B) \wedge C = A \wedge C + B \wedge C$
    \item \textbf{Asociatividad}: $(A \wedge B) \wedge C = A \wedge (B \wedge C)$
    \item \textbf{Anticonmutatividad}: $A \wedge B = (-1)^{|A||B|} B \wedge A$
    \item \textbf{Nilpotencia}: $A \wedge A = 0$ para vectores $A$
    \item \textbf{Graduación}: $|A \wedge B| = |A| + |B|$
\end{enumerate}
\end{proposition}

\begin{proof}
Las propiedades se derivan directamente de la definición del producto exterior como la parte totalmente antisimétrica del producto geométrico. La distributividad es consecuencia de la bilinealidad del producto geométrico. La asociatividad se verifica por inducción sobre los grados. La anticonmutatividad resulta de que intercambiar dos vectores en un producto exterior introduce un factor $(-1)$, y para multivectores de grados $|A|$ y $|B|$ se acumulan $|A| \cdot |B|$ intercambios. La nilpotencia sigue de la anticonmutatividad: $A \wedge A = (-1)^{1 \cdot 1} A \wedge A = -A \wedge A$, por lo que $A \wedge A = 0$. La graduación es inmediata: el producto exterior de un $p$-vector con un $q$-vector produce un $(p+q)$-vector por construcción.
\end{proof}

\subsubsection{Involuciones y Operaciones Conjugadas}

\begin{definition}[Involuciones de Clifford]\label{def:involuciones_clifford}
Las involuciones principales son:

\textbf{Reversión} (grade preserving, order reversing):
\begin{equation}
    \widetilde{A} = \sum_k (-1)^{k(k-1)/2} \langle A \rangle_k
\end{equation}

\textbf{Conjugación principal} (grade reversing):
\begin{equation}
    \overline{A} = \sum_k (-1)^k \langle A \rangle_k
\end{equation}

\textbf{Conjugación de Clifford} (combinación):
\begin{equation}
    A^* = \overline{\widetilde{A}} = \sum_k (-1)^{k(k+1)/2} \langle A \rangle_k
\end{equation}
\end{definition}

\begin{proposition}[Norma de Clifford]\label{prop:norma_clifford}
Para elementos $A, B$ del grupo de Clifford (grupo de Lipschitz) $\Gamma(V,Q)$, es decir, elementos invertibles expresables como productos de vectores no isótropos, la norma cuadrática
\begin{equation}
    N(A) = A \widetilde{A} = \widetilde{A} A
\end{equation}
es un escalar y satisface $N(AB) = N(A)N(B)$ (propiedad multiplicativa). Para multivectores generales, $A\widetilde{A}$ puede no ser un escalar y la propiedad multiplicativa no se cumple en general.
\end{proposition}

\begin{proof}
Sea $A = v_1 v_2 \cdots v_k$ un producto de vectores no isótropos. Entonces $\widetilde{A} = v_k \cdots v_2 v_1$ (reversión del orden). Se tiene:
\[
A\widetilde{A} = v_1 v_2 \cdots v_k v_k \cdots v_2 v_1 = v_1 v_2 \cdots v_{k-1} Q(v_k) v_{k-1} \cdots v_2 v_1 = Q(v_k) \cdot v_1 \cdots v_{k-1} v_{k-1} \cdots v_1
\]
Iterando, $N(A) = \prod_{i=1}^k Q(v_i) \in \mathbb{R}$, que es un escalar. La propiedad multiplicativa se verifica: $N(AB) = AB\widetilde{AB} = AB\widetilde{B}\widetilde{A} = A \cdot N(B) \cdot \widetilde{A} = N(B) \cdot A\widetilde{A} = N(A)N(B)$, donde usamos que $N(B)$ es escalar y por tanto conmuta con todos los elementos.
\end{proof}

\subsection{Geometric Algebra Transformers (GATr): Arquitectura Fundamental}

\subsubsection{Álgebra Geométrica Proyectiva}

GATr de \citet{brehmer2023geometric} utiliza el álgebra geométrica proyectiva $\text{Cl}(3,0,1)$ (también denotada $\mathbb{G}_{3,0,1}$ o PGA), proporcionando una representación unificada de objetos geométricos 3D.

\begin{definition}[Álgebra Proyectiva 3D]\label{def:algebra_proyectiva_3d}
El álgebra $\text{Cl}(3,0,1)$ se construye sobre un espacio vectorial de dimensión 4 con base $\{\mathbf{e}_0, \mathbf{e}_1, \mathbf{e}_2, \mathbf{e}_3\}$ donde la signatura $(3,0,1)$ indica 3 direcciones positivas y 1 degenerada:
\begin{align}
    \mathbf{e}_i^2 &= +1 \quad \text{para } i = 1,2,3 \text{ (dimensiones espaciales)} \\
    \mathbf{e}_0^2 &= 0 \quad \text{(dirección degenerada/nula)}
\end{align}
La dimensión total del álgebra es $2^4 = 16$. La dirección degenerada $\mathbf{e}_0$ es esencial: permite representar geometría afín (incluyendo traslaciones), no solo geometría lineal.
\end{definition}

\begin{proposition}[Representación Unificada de Objetos Geométricos]\label{prop:representacion_unificada_gatr}
En $\text{Cl}(3,0,1)$, los multivectores se descomponen en grados $k = 0, 1, 2, 3, 4$ con dimensiones $\binom{4}{k} = 1, 4, 6, 4, 1$ respectivamente (total: 16). Los objetos geométricos fundamentales de $\mathbb{R}^3$ admiten representaciones canónicas en la convención ``plane-based'' del PGA \citep{brehmer2023geometric}:

\textbf{Escalares} ($k=0$): $\alpha \in \mathbb{R}$

\textbf{Planos} ($k=1$): $\pi = n_1 \mathbf{e}_1 + n_2 \mathbf{e}_2 + n_3 \mathbf{e}_3 + d\, \mathbf{e}_0$

\textbf{Líneas} ($k=2$): $\ell \in \langle \mathbf{e}_i \wedge \mathbf{e}_j \rangle$

\textbf{Puntos} ($k=3$): $P = x_1 (\mathbf{e}_0 \wedge \mathbf{e}_2 \wedge \mathbf{e}_3) + x_2 (\mathbf{e}_0 \wedge \mathbf{e}_3 \wedge \mathbf{e}_1) + x_3 (\mathbf{e}_0 \wedge \mathbf{e}_1 \wedge \mathbf{e}_2) + \mathbf{e}_1 \wedge \mathbf{e}_2 \wedge \mathbf{e}_3$

\textbf{Pseudoescalar} ($k=4$): $I = \mathbf{e}_0 \wedge \mathbf{e}_1 \wedge \mathbf{e}_2 \wedge \mathbf{e}_3$
\end{proposition}

\begin{proof}
El álgebra $\text{Cl}(3,0,1)$ tiene base $\{\mathbf{e}_0, \mathbf{e}_1, \mathbf{e}_2, \mathbf{e}_3\}$, luego $\dim \text{Cl}(3,0,1) = 2^4 = 16$. La descomposición por grado $k$ sigue de que las $k$-blades forman un espacio de dimensión $\binom{4}{k}$: $\binom{4}{0} + \binom{4}{1} + \binom{4}{2} + \binom{4}{3} + \binom{4}{4} = 1 + 4 + 6 + 4 + 1 = 16$.

Para las representaciones canónicas: un plano $n_1 x_1 + n_2 x_2 + n_3 x_3 + d = 0$ en $\mathbb{R}^3$ se codifica como 1-vector $\pi = \sum n_i \mathbf{e}_i + d\,\mathbf{e}_0$, donde la componente $\mathbf{e}_0$ captura el desplazamiento afín. Las líneas, como intersección de dos planos, se representan mediante el producto exterior $\pi_1 \wedge \pi_2$, que es un 2-vector. Los puntos, como intersección de tres planos, se representan como 3-vectores $\pi_1 \wedge \pi_2 \wedge \pi_3$. Las fórmulas explícitas se obtienen expandiendo estos productos exteriores en la base canónica.
\end{proof}

\subsubsection{Representación 16-dimensional y Embedding Geométrico}

\begin{definition}[Espacio de Representación GATr]\label{def:espacio_representacion_gatr}
Cada multivector en $\text{Cl}(3,0,1)$ se representa como un vector de 16 componentes:

\begin{equation}
    M = \alpha_0 + \sum_{i=0}^{3} \alpha_i \mathbf{e}_i + \sum_{0 \leq i < j \leq 3} \alpha_{ij} \mathbf{e}_i \wedge \mathbf{e}_j + \sum_{0 \leq i < j < k \leq 3} \alpha_{ijk} \mathbf{e}_i \wedge \mathbf{e}_j \wedge \mathbf{e}_k + \alpha_{0123} I
    \label{eq:gatr_multivector}
\end{equation}

La base canónica 16-dimensional se descompone por grados:
\begin{align}
\text{Grado 0 (1):} \quad & \{1\} \\
\text{Grado 1 (4):} \quad & \{\mathbf{e}_0, \mathbf{e}_1, \mathbf{e}_2, \mathbf{e}_3\} \\
\text{Grado 2 (6):} \quad & \{\mathbf{e}_0 \wedge \mathbf{e}_1, \mathbf{e}_0 \wedge \mathbf{e}_2, \mathbf{e}_0 \wedge \mathbf{e}_3, \mathbf{e}_1 \wedge \mathbf{e}_2, \mathbf{e}_1 \wedge \mathbf{e}_3, \mathbf{e}_2 \wedge \mathbf{e}_3\} \\
\text{Grado 3 (4):} \quad & \{\mathbf{e}_0 \wedge \mathbf{e}_1 \wedge \mathbf{e}_2, \mathbf{e}_0 \wedge \mathbf{e}_1 \wedge \mathbf{e}_3, \mathbf{e}_0 \wedge \mathbf{e}_2 \wedge \mathbf{e}_3, \mathbf{e}_1 \wedge \mathbf{e}_2 \wedge \mathbf{e}_3\} \\
\text{Grado 4 (1):} \quad & \{I = \mathbf{e}_0 \wedge \mathbf{e}_1 \wedge \mathbf{e}_2 \wedge \mathbf{e}_3\}
\end{align}
\end{definition}

\begin{algorithm}[H]
\caption{Embedding de Objetos Geométricos en GATr (PGA)}
\label{alg:embedding_gatr}
\begin{algorithmic}[1]
    \State \textbf{Entrada:} Objeto geométrico (punto, plano, etc.)
    \Function{EmbedPoint}{$\mathbf{x} \in \mathbb{R}^3$}
        \State $P = x_1 (\mathbf{e}_0 \wedge \mathbf{e}_2 \wedge \mathbf{e}_3) + x_2 (\mathbf{e}_0 \wedge \mathbf{e}_3 \wedge \mathbf{e}_1) + x_3 (\mathbf{e}_0 \wedge \mathbf{e}_1 \wedge \mathbf{e}_2) + \mathbf{e}_1 \wedge \mathbf{e}_2 \wedge \mathbf{e}_3$
        \State \Return $P$ \Comment{Trivector, 16 componentes}
    \EndFunction
    \Function{EmbedPlane}{$\mathbf{n} \in \mathbb{R}^3, d \in \mathbb{R}$}
        \State $\pi = n_1 \mathbf{e}_1 + n_2 \mathbf{e}_2 + n_3 \mathbf{e}_3 + d\, \mathbf{e}_0$
        \State \Return $\pi$ \Comment{Vector (grado 1), 16 componentes}
    \EndFunction
    \Function{EmbedRotation}{$\mathbf{n} \in \mathbb{R}^3, \theta \in \mathbb{R}$}
        \State $B = \sin(\theta/2) (n_1 \mathbf{e}_2 \wedge \mathbf{e}_3 + n_2 \mathbf{e}_3 \wedge \mathbf{e}_1 + n_3 \mathbf{e}_1 \wedge \mathbf{e}_2)$
        \State \Return $\cos(\theta/2) + B$ \Comment{Rotor unitario}
    \EndFunction
\end{algorithmic}
\end{algorithm}

\subsubsection{Arquitectura Transformer Equivariante}

\begin{definition}[Capa Lineal Equivariante en GATr]\label{def:capa_lineal_gatr}
Una transformación lineal $f: \text{Cl}(3,0,1) \to \text{Cl}(3,0,1)$ es Pin(3,0,1)-equivariante si y sólo si puede escribirse como combinación de proyecciones de grado y multiplicación por $\mathbf{e}_0$:

\begin{equation}
    f(M) = \sum_{k=0}^4 \alpha_k \langle M \rangle_k + \sum_{k=0}^4 \beta_k\, \mathbf{e}_0 \langle M \rangle_k
    \label{eq:linear_equivariant_gatr}
\end{equation}

donde $\alpha_k, \beta_k \in \mathbb{R}$ son parámetros aprendibles (9 parámetros en total, dado que $\beta_0 = 0$ por razones de grado). Véase Proposición~1 de \citet{brehmer2023geometric}.
\end{definition}

\begin{theorem}[E(3)-Equivarianza de Operaciones en GATr]\label{thm:equivarianza_automatica_gatr}
El producto geométrico en $\text{Cl}(3,0,1)$ es Pin(3,0,1)-equivariante: si $\rho_u(x) = uxu^{-1}$ denota la acción del sandwich product para $u \in \text{Pin}(3,0,1)$, entonces $\rho_u(xy) = \rho_u(x)\rho_u(y)$ para todo par de multivectores $x, y$.
\end{theorem}

\begin{proof}
$\rho_u(xy) = u(xy)u^{-1} = (uxu^{-1})(uyu^{-1}) = \rho_u(x)\rho_u(y)$, donde la segunda igualdad usa la asociatividad del producto geométrico y la propiedad $u^{-1}u = 1$. Véase Apéndice~A de \citet{brehmer2023geometric} para la verificación de que Pin(3,0,1) actúa por automorfismos de álgebra.
\end{proof}

\begin{remark}
No todas las operaciones de Clifford son automáticamente Pin-equivariantes. En particular, la operación dual $x \mapsto xI^{-1}$ es Spin-equivariante pero no Pin-equivariante (véase la discusión en \citet{brehmer2023geometric}, Sección~3.3). GATr maneja esta distinción cuidadosamente en su diseño arquitectónico.
\end{remark}

\begin{definition}[Mecanismo de Atención Geométrica]\label{def:atencion_geometrica_gatr}
La atención en GATr utiliza el producto interior de Clifford en $\text{Cl}(3,0,1)$ como medida de similaridad:

\begin{equation}
    \text{Attn}(Q, K, V) = \text{softmax}\left(\frac{Q \cdot K^T}{\sqrt{16}}\right) V
\end{equation}

donde $Q \cdot K^T$ denota la aplicación del producto interior de Clifford elemento a elemento, y $\sqrt{16}$ es el factor de escala correspondiente a la dimensión del álgebra.
\end{definition}

\begin{algorithm}[H]
\caption{Forward Pass de GATr}
\label{alg:forward_gatr}
\begin{algorithmic}[1]
    \State \textbf{Entrada:} Secuencia de multivectores $\{M_1, M_2, \ldots, M_n\} \subset \text{Cl}(3,0,1)^{16}$
    \For{cada capa $l = 1, 2, \ldots, L$}
        \State $Q^{(l)} = \text{LinearEquivariant}(M^{(l-1)})$
        \State $K^{(l)} = \text{LinearEquivariant}(M^{(l-1)})$
        \State $V^{(l)} = \text{LinearEquivariant}(M^{(l-1)})$
        \State $A^{(l)} = \text{MultiHeadAttention}(Q^{(l)}, K^{(l)}, V^{(l)})$
        \State $M^{(l)} = \text{LayerNorm}(M^{(l-1)} + A^{(l)})$
        \State $M^{(l)} = \text{LayerNorm}(M^{(l)} + \text{FFNEquivariant}(M^{(l)}))$
    \EndFor
    \State \Return $M^{(L)}$
\end{algorithmic}
\end{algorithm}

\subsubsection{Componentes Arquitectónicos Especializados}

\begin{definition}[Normalización Consciente de Grado]\label{def:normalizacion_grado}
La normalización por capas en GATr se aplica separadamente a cada grado:

\begin{equation}
    \text{LayerNorm}(M) = \sum_{k=0}^4 \frac{\langle M \rangle_k - \mu_k}{\sqrt{\sigma_k^2 + \epsilon}} \gamma_k + \beta_k
\end{equation}

donde $\mu_k$ y $\sigma_k^2$ son la media y varianza de las componentes de grado $k$.
\end{definition}

\begin{definition}[Función de Activación Equivariante en GATr]\label{def:activacion_equivariante_gatr}
La activación equivariante en GATr utiliza un mecanismo de gating escalar \citep{brehmer2023geometric}: dado un multivector $x \in \text{Cl}(3,0,1)$, la activación se define como:

\begin{equation}
    \text{GatedGELU}(x) = \text{GELU}(\langle x \rangle_0) \cdot x
\end{equation}

donde $\langle x \rangle_0$ es la componente escalar (grado 0) del multivector. Esta función es Pin(3,0,1)-equivariante porque el escalar $\langle x \rangle_0$ es invariante bajo la acción del grupo, y el producto escalar-multivector preserva la estructura de grados.
\end{definition}

El marco de atención Clifford presentado se extiende naturalmente a otras signaturas métricas. En particular, el álgebra $\text{Cl}(1,3)$ sobre el espacio-tiempo de Minkowski permite construir arquitecturas Lorentz-equivariantes (L-GATr, \citet{spinner2024lgatr}) para aplicaciones en física relativista, como clasificación de jets y cálculo de amplitudes de dispersión.

\subsection{Análisis Teórico y Propiedades Fundamentales}

\subsubsection{Poder Expresivo de Arquitecturas de Álgebra Geométrica}

\begin{remark}[Aproximación Universal en Álgebras de Clifford: Problema Abierto]\label{rem:aproximacion_universal_clifford}
\citet{brehmer2023geometric} conjeturan que GATr podría ser un aproximador universal para funciones continuas $E(3)$-equivariantes $f: (\mathbb{R}^3)^n \to \mathbb{R}^m$ en conjuntos compactos. Sin embargo, los propios autores señalan explícitamente que ``GATr is not yet shown to be a universal approximator, which is an interesting direction for future work'' (p.~10). La demostración de esta conjetura requeriría verificar que las operaciones de GATr (capas lineales equivariantes, producto geométrico, atención) generan un álgebra densa de funciones equivariantes, lo cual permanece como un problema abierto.
\end{remark}

\begin{remark}[Eficiencia Paramétrica]\label{cor:eficiencia_parametrica_clifford}
Una ventaja práctica de GATr es que la representación en $\text{Cl}(3,0,1)$ codifica de forma compacta (en 16 componentes) objetos geométricos que en representación tensorial directa requerirían más parámetros para garantizar equivarianza. Sin embargo, una comparación rigurosa de eficiencia paramétrica entre GATr y arquitecturas tensoriales depende de la clase de funciones considerada y no ha sido establecida formalmente.
\end{remark}

\subsubsection{Análisis de Complejidad Computacional}

\begin{proposition}[Complejidad de Operaciones de Clifford]\label{prop:complejidad_clifford}
Para multivectores de dimensión $2^n$, las operaciones fundamentales tienen complejidad:

\begin{align}
    \text{Producto geométrico:} &\quad O(4^n) \\
    \text{Producto wedge:} &\quad O(4^n) \\
    \text{Involuciones:} &\quad O(2^n) \\
    \text{Proyecciones de grado:} &\quad O(2^n)
\end{align}

Sin embargo, para $n \leq 4$ (como en GATr, donde $n = 4$ para $\text{Cl}(3,0,1)$), estas operaciones son altamente eficientes.
\end{proposition}

\begin{proof}
El producto geométrico de dos multivectores de dimensión $2^n$ requiere multiplicar cada componente del primer multivector con cada componente del segundo, resultando en $2^n \times 2^n = 4^n$ multiplicaciones escalares (más consultas a la tabla de signos). El producto wedge tiene la misma complejidad asintótica, ya que se obtiene seleccionando las componentes del producto geométrico que incrementan el grado. Las involuciones (reversión, conjugación) operan componente a componente multiplicando por $\pm 1$ según el grado, lo que requiere $O(2^n)$ operaciones. Las proyecciones de grado simplemente seleccionan las componentes correspondientes, también en $O(2^n)$.
\end{proof}

\begin{algorithm}[H]
\caption{Producto Geométrico Eficiente}
\label{alg:producto_geometrico_eficiente}
\begin{algorithmic}[1]
    \State \textbf{Entrada:} Multivectores $A, B \in \text{Cl}(p,q)$
    \Function{GeometricProduct}{$A, B$}
        \State $C = \text{zeros}(2^{p+q})$
        \For{$i = 0$ to $2^{p+q} - 1$}
            \For{$j = 0$ to $2^{p+q} - 1$}
                \State $(k, \text{sign}) = \text{ProductTable}[i][j]$
                \State $C[k] += \text{sign} \times A[i] \times B[j]$
            \EndFor
        \EndFor
        \State \Return $C$
    \EndFunction
\end{algorithmic}
\end{algorithm}

\subsection{Implementación Práctica y Consideraciones Computacionales}

\subsubsection{Optimizaciones para Hardware Moderno}

\begin{definition}[Paralelización Masiva en GA]\label{def:paralelizacion_ga}
Las operaciones de álgebra geométrica son intrínsecamente paralelizables. Para un batch de $B$ multivectores:

\begin{equation}
    \text{BatchGeometricProduct}(A_{B \times 16}, B_{B \times 16}) = C_{B \times 16}
\end{equation}

puede ejecutarse con $B \times 16^2$ operaciones paralelas en GPU.
\end{definition}

\begin{algorithm}[H]
\caption{Implementación GPU-optimizada de GATr}
\label{alg:gatr_gpu}
\begin{algorithmic}[1]
    \State \textbf{Entrada:} Batch de multivectores $\mathbf{M} \in \mathbb{R}^{B \times N \times 16}$
    \Function{GATrForwardGPU}{$\mathbf{M}$}
        \State \textbf{// Fase 1: Proyecciones paralelas}
        \State $Q, K, V = \text{ParallelLinear}(\mathbf{M})$ \Comment{$O(BN)$ threads}
        
        \State \textbf{// Fase 2: Atención con productos de Clifford}
        \State $\text{Scores} = \text{BatchGeometricProduct}(Q, K^T)$ \Comment{Vectorizado}
        \State $\text{Attn} = \text{Softmax}(\text{ScalarPart}(\text{Scores}))$
        
        \State \textbf{// Fase 3: Agregación equivariante}
        \State $\text{Output} = \text{BatchGeometricProduct}(\text{Attn}, V)$
        \State \Return $\text{Output}$
    \EndFunction
\end{algorithmic}
\end{algorithm}

\subsubsection{Estabilidad Numérica y Consideraciones Prácticas}

\begin{proposition}[Estabilidad Numérica de Involuciones]\label{prop:estabilidad_numerica_involuciones}
Las involuciones de Clifford son numéricamente estables bajo aritmética de punto flotante. Para errores de redondeo $\epsilon$:

\begin{equation}
    \|\widetilde{A + \epsilon} - \widetilde{A}\| \leq C\epsilon
\end{equation}

donde $C$ es una constante independiente de la magnitud de $A$.
\end{proposition}

\begin{proof}
La reversión actúa sobre cada componente homogénea de grado $k$ multiplicándola por $(-1)^{k(k-1)/2} \in \{-1, +1\}$. Por tanto, $\widetilde{A + \epsilon} - \widetilde{A} = \widetilde{\epsilon}$, y $\|\widetilde{\epsilon}\| = \|\epsilon\|$ ya que la reversión preserva la norma euclidiana (solo cambia signos de componentes). La constante $C = 1$ es independiente de $A$. El mismo argumento aplica a la conjugación principal y la conjugación de Clifford, que son igualmente operaciones componente a componente con factores $\pm 1$.
\end{proof}

\textbf{Inicialización de Parámetros.}\label{rem:inicializacion_parametros_ga} La inicialización de parámetros en GATr debe respetar la estructura de grados. Los pesos se inicializan típicamente como:

\begin{equation}
    w_k \sim \mathcal{N}\left(0, \frac{1}{\binom{n}{k}}\right)
\end{equation}

donde $\binom{n}{k}$ es el número de elementos de grado $k$.

\subsection{Direcciones de Investigación y Desarrollos Futuros}

\subsubsection{Extensiones a Álgebras de Clifford Superiores}

\begin{definition}[GATr Generalizado]\label{def:gatr_generalizado}
Un GATr generalizado sobre $\text{Cl}(p,q)$ mantiene equivarianza bajo el grupo $\text{O}(p,q)$ y puede representar objetos geométricos en espacios de signatura arbitraria.
\end{definition}

\begin{hypothesis}[Escalabilidad Dimensional]\label{hyp:escalabilidad_dimensional}
Existe un compromiso fundamental entre la dimensionalidad del álgebra de Clifford ($2^{p+q}$) y la capacidad expresiva, sugiriendo dimensiones óptimas para diferentes clases de problemas.
\end{hypothesis}

\subsubsection{Integración con Otras Arquitecturas de GDL}

El futuro de las arquitecturas basadas en álgebra geométrica incluye:

\begin{itemize}
    \item \textbf{Híbridos CNN-GATr}: Combinación de convoluciones geométricas con atención de Clifford
    \item \textbf{GATr-GNN}: Aplicación de álgebra geométrica en redes de paso de mensajes
    \item \textbf{Modelos Generativos}: Desarrollos en VAEs y GANs equivariantes basados en GA
    \item \textbf{Arquitecturas Cuánticas}: Extensión a computación cuántica aprovechando el álgebra de Clifford natural
\end{itemize}

\subsubsection{Aplicaciones Emergentes}

Las aplicaciones futuras de GATr incluyen:
\begin{itemize}
    \item \textbf{Robótica Avanzada}: Control de manipuladores con simetrías complejas
    \item \textbf{Realidad Aumentada}: Procesamiento de escenas 3D en tiempo real
    \item \textbf{Biología Estructural}: Predicción de estructuras proteicas con simetrías cristalográficas
    \item \textbf{Astrofísica Computacional}: Simulaciones de campos gravitacionales relativistas
\end{itemize}

\section{Topología Algebraica}

Las secciones anteriores han mostrado cómo los mecanismos de atención se benefician de incorporar estructura geométrica, desde simetrías de grupo hasta propiedades espectrales de grafos. El siguiente nivel de abstracción consiste en explotar la \textit{topología} de los datos: propiedades que se preservan bajo deformaciones continuas, como la presencia de ciclos o cavidades. La incorporación de conceptos topológicos en mecanismos de atención representa una de las fronteras más prometedoras del \gls{GDL}, permitiendo capturar características estructurales de orden superior que van más allá de las conexiones locales en grafos. Este enfoque, fundamentado en la topología algebraica clásica, proporciona herramientas matemáticas sofisticadas para analizar la forma global de los datos y sus transformaciones bajo operaciones de aprendizaje profundo.

\subsection{Fundamentos de Topología Algebraica}

\subsubsection{Complejos Simpliciales: Teoría General}

Recordamos las definiciones de $k$-símplice (\autoref{def:simplicio_gdl}), cara de un símplice (\autoref{def:cara_simplicio_gdl}) y complejo simplicial (\autoref{def:complejo_simplicial}), introducidas en el \autoref{ch:gdl}. Los casos fundamentales son:
\begin{align}
    k=0: &\quad \text{0-símplice (vértice)} \\
    k=1: &\quad \text{1-símplice (arista)} \\
    k=2: &\quad \text{2-símplice (triángulo)} \\
    k=3: &\quad \text{3-símplice (tetraedro)} \\
    k=n: &\quad \text{$n$-símplice (politopo de dimensión } n\text{)}
\end{align}

\begin{theorem}[Teorema de Realización Simplicial]\label{thm:realizacion_simplicial}
Todo complejo simplicial abstracto finito admite una realización geométrica en $\mathbb{R}^n$ para $n$ suficientemente grande. Para la demostración, véase \citep{munkres2000topology}.
\end{theorem}

\subsubsection{Orientaciones y Cadenas Simpliciales}

\begin{definition}[Orientación de un Simplicio]\label{def:orientacion_simplicio}
Una orientación de un $k$-simplicio $\sigma = [v_0, v_1, \ldots, v_k]$ es una clase de equivalencia de ordenamientos de sus vértices bajo permutaciones pares. Denotamos la orientación positiva como $[v_0, v_1, \ldots, v_k]$ y la negativa como $-[v_0, v_1, \ldots, v_k]$.
\end{definition}

\textbf{Número de Orientaciones.}\label{rem:numero_orientaciones} Un $k$-simplicio con $k \geq 1$ tiene exactamente dos orientaciones. Los 0-simplicios tienen orientación trivial.

\begin{definition}[Espacio de Cadenas Simpliciales]\label{def:espacio_cadenas}
Para un complejo simplicial $\mathcal{K}$, el espacio de $k$-cadenas es el $\mathbb{Z}$-módulo libre:
\begin{equation}
    C_k(\mathcal{K}) = \bigoplus_{\sigma \in \mathcal{K}_k} \mathbb{Z} \cdot \sigma
\end{equation}
donde $\mathcal{K}_k$ denota el conjunto de $k$-simplicios orientados en $\mathcal{K}$.
\end{definition}

Un elemento típico de $C_k(\mathcal{K})$ se escribe como:
\begin{equation}
    c = \sum_{\sigma \in \mathcal{K}_k} n_{\sigma} \sigma
\end{equation}
donde $n_{\sigma} \in \mathbb{Z}$ son los coeficientes.

\subsubsection{Operadores de Frontera}

El operador de frontera $\partial_k: C_k(\mathcal{K}) \to C_{k-1}(\mathcal{K})$ (véase \autoref{def:cadenas_frontera} del \autoref{ch:cnn}) satisface $\partial_k[v_0, \ldots, v_k] = \sum_{i=0}^k (-1)^i [v_0, \ldots, \hat{v_i}, \ldots, v_k]$.\label{def:operador_frontera}

\begin{example}[Operadores de Frontera Explícitos]\label{ex:operadores_frontera_explicitos}
\begin{align}
    \partial_0[v] &= 0 \\
    \partial_1[v_0, v_1] &= [v_1] - [v_0] \\
    \partial_2[v_0, v_1, v_2] &= [v_1, v_2] - [v_0, v_2] + [v_0, v_1] \\
    \partial_3[v_0, v_1, v_2, v_3] &= [v_1, v_2, v_3] - [v_0, v_2, v_3] + [v_0, v_1, v_3] - [v_0, v_1, v_2]
\end{align}
\end{example}

\begin{theorem}[Propiedad Fundamental de la Frontera]\label{thm:frontera_frontera_cero}
Para todo $k \geq 1$:
\begin{equation}
    \partial_{k-1} \circ \partial_k = 0
\end{equation}
Es decir, el operador frontera aplicado dos veces da cero.
\end{theorem}

\begin{proof}
Para un $k$-simplicio $[v_0, \ldots, v_k]$:
\begin{align}
    \partial_{k-1} \partial_k[v_0, \ldots, v_k] &= \partial_{k-1} \sum_{i=0}^k (-1)^i [v_0, \ldots, \hat{v_i}, \ldots, v_k] \\
    &= \sum_{i=0}^k (-1)^i \sum_{\substack{j=0 \\ j \neq i}}^k (-1)^{j'} [v_0, \ldots, \hat{v_j}, \ldots, \hat{v_i}, \ldots, v_k]
\end{align}
donde $j' = j$ si $j < i$ y $j' = j-1$ si $j > i$. Cada $(k-2)$-cara aparece exactamente dos veces con signos opuestos, por lo que la suma es cero.
\end{proof}

\begin{definition}[Complejo de Cadenas]\label{def:complejo_cadenas}
El complejo de cadenas asociado a $\mathcal{K}$ es la secuencia:
\begin{equation}
    \cdots \to C_{k+1}(\mathcal{K}) \xrightarrow{\partial_{k+1}} C_k(\mathcal{K}) \xrightarrow{\partial_k} C_{k-1}(\mathcal{K}) \to \cdots \to C_0(\mathcal{K}) \xrightarrow{\partial_0} 0
\end{equation}
\end{definition}

\subsection{Homología Simplicial}

\subsubsection{Ciclos, Fronteras y Grupos de Homología}

\begin{definition}[Ciclos y Fronteras]\label{def:ciclos_fronteras}
Para un complejo simplicial $\mathcal{K}$:
\begin{align}
    Z_k(\mathcal{K}) &= \ker(\partial_k) \quad \text{(espacio de } k\text{-ciclos)} \\
    B_k(\mathcal{K}) &= \operatorname{im}(\partial_{k+1}) \quad \text{(espacio de } k\text{-fronteras)}
\end{align}
\end{definition}

Por el teorema fundamental de la frontera, $B_k(\mathcal{K}) \subseteq Z_k(\mathcal{K})$ para todo $k$.

Los grupos de homología $H_k(\mathcal{K}) = Z_k/B_k$ y los números de Betti $\beta_k = \dim H_k$ se definen como en la \autoref{def:homologia_betti} del \autoref{ch:gdl}.\label{def:grupos_homologia}\label{def:numeros_betti}

\begin{remark}[Interpretación Geométrica de los Números de Betti]\label{rem:interpretacion_betti}
Los números de Betti capturan características topológicas:
\begin{align}
    \beta_0 &= \text{número de componentes conexas} \\
    \beta_1 &= \text{número de bucles independientes} \\
    \beta_2 &= \text{número de cavidades 2-dimensionales} \\
    &\vdots \\
    \beta_k &= \text{número de } k\text{-cavidades independientes}
\end{align}
\end{remark}

\begin{example}[Homología de Estructuras Básicas]\label{ex:homologia_basicas}
\textbf{Esfera $S^2$}: $H_0(S^2) = \mathbb{Z}$, $H_1(S^2) = 0$, $H_2(S^2) = \mathbb{Z}$

\textbf{Toro $T^2$}: $H_0(T^2) = \mathbb{Z}$, $H_1(T^2) = \mathbb{Z} \oplus \mathbb{Z}$, $H_2(T^2) = \mathbb{Z}$

\textbf{Plano proyectivo real $\mathbb{RP}^2$}: $H_0(\mathbb{RP}^2) = \mathbb{Z}$, $H_1(\mathbb{RP}^2) = \mathbb{Z}_2$, $H_2(\mathbb{RP}^2) = 0$
\end{example}

\subsubsection{Homología Persistente}

\begin{definition}[Filtración]\label{def:filtracion}
Una filtración de un complejo simplicial $\mathcal{K}$ es una secuencia anidada:
\begin{equation}
    \emptyset = \mathcal{K}^0 \subseteq \mathcal{K}^1 \subseteq \mathcal{K}^2 \subseteq \cdots \subseteq \mathcal{K}^m = \mathcal{K}
\end{equation}
donde cada $\mathcal{K}^i$ es un subcomplejo simplicial de $\mathcal{K}$.
\end{definition}

\begin{definition}[Módulos de Homología Persistente]\label{def:modulos_homologia_persistente}
Para una filtración, los módulos de homología persistente son:
\begin{equation}
    H_k^{i,j} = \operatorname{im}(H_k(\mathcal{K}^i) \to H_k(\mathcal{K}^j)) \quad \text{para } i \leq j
\end{equation}
\end{definition}

\begin{definition}[Diagrama de Persistencia]\label{def:diagrama_persistencia}
El diagrama de persistencia $\text{Pers}_k(\mathcal{K})$ es el multiconjunto de puntos $(b,d) \in \mathbb{R}^2$ donde:
\begin{itemize}
    \item $b$ es el \textbf{tiempo de nacimiento} de una característica homológica
    \item $d$ es el \textbf{tiempo de muerte} de la misma característica
    \item La \textbf{persistencia} es $d - b$
\end{itemize}
\end{definition}

\begin{theorem}[Teorema de Descomposición de Módulos]\label{thm:descomposicion_modulos}
Todo módulo de persistencia finitamente generado se descompone como:
\begin{equation}
    M \cong \bigoplus_i \mathbb{k}[b_i, d_i) \oplus \bigoplus_j \mathbb{k}[b_j, \infty)
\end{equation}
donde $\mathbb{k}[b,d)$ representa un módulo intervalo. Este resultado clásico de la teoría de persistencia se demuestra mediante la clasificación de módulos finitamente generados sobre anillos de polinomios, véase \citep{Hatcher2002}.
\end{theorem}

\begin{algorithm}[H]
\caption{Algoritmo de Homología Persistente}\label{alg:homologia_persistente}
\textbf{Entrada}: Filtración $\mathcal{K}^0 \subseteq \mathcal{K}^1 \subseteq \cdots \subseteq \mathcal{K}^m$
\begin{enumerate}
    \item Construir matrices de frontera $\partial_k^i$ para cada $\mathcal{K}^i$
    \item Aplicar reducción de Smith para diagonalizar
    \item Identificar pares (nacimiento, muerte) de características
    \item Construir diagrama de persistencia
\end{enumerate}
\textbf{Salida}: Diagrama de persistencia $\text{Pers}_k(\mathcal{K})$
\end{algorithm}

\subsection{Cohomología y Teoría de de Rham Discreta}

\subsubsection{Complejos de Cocadenas}

\begin{definition}[Espacio de Cocadenas]\label{def:espacio_cocadenas}
El espacio de $k$-cocadenas es el dual de las cadenas:
\begin{equation}
    C^k(\mathcal{K}) = \operatorname{Hom}(C_k(\mathcal{K}), \mathbb{R})
\end{equation}
\end{definition}

\begin{definition}[Operador de Coborde]\label{def:operador_coborde}
El operador de coborde $\delta^k: C^k(\mathcal{K}) \to C^{k+1}(\mathcal{K})$ se define como:
\begin{equation}
    (\delta^k \omega)(\sigma) = \omega(\partial_{k+1} \sigma)
\end{equation}
para $\omega \in C^k(\mathcal{K})$ y $\sigma \in C_{k+1}(\mathcal{K})$.
\end{definition}

\begin{theorem}[Propiedad del Coborde]\label{thm:coborde_coborde_cero}
$\delta^{k+1} \circ \delta^k = 0$ para todo $k$.
\end{theorem}

\begin{proof}
Para $\omega \in C^k(\mathcal{K})$ y $\tau \in C_{k+2}(\mathcal{K})$:
\[
(\delta^{k+1} \circ \delta^k)(\omega)(\tau) = (\delta^k \omega)(\partial_{k+2}\tau) = \omega(\partial_{k+1} \circ \partial_{k+2}(\tau)) = \omega(0) = 0
\]
donde usamos que $\partial_{k+1} \circ \partial_{k+2} = 0$ (propiedad fundamental del operador de frontera).
\end{proof}

\begin{definition}[Grupos de Cohomología]\label{def:grupos_cohomologia}
Los grupos de cohomología son:
\begin{equation}
    H^k(\mathcal{K}) = \frac{\ker(\delta^k)}{\operatorname{im}(\delta^{k-1})}
\end{equation}
\end{definition}

\begin{theorem}[Dualidad de Poincaré para Complejos]\label{thm:dualidad_poincare_complejos}
Para complejos simpliciales orientables sin borde de dimensión $n$:
\begin{equation}
    H^k(\mathcal{K}) \cong H_{n-k}(\mathcal{K})
\end{equation}
Véase \citep{Hatcher2002} para la demostración.
\end{theorem}

\subsubsection{Teoría de Hodge Discreta}

\begin{definition}[Operadores de Hodge Discretos]\label{def:operadores_hodge_discretos}
En un complejo simplicial con métrica discreta, definimos:
\begin{align}
    \Delta_k &= \partial_{k+1}^* \partial_{k+1} + \partial_k \partial_k^* \quad \text{(Laplaciano de Hodge)} \\
    \Delta_k^{\text{up}} &= \partial_{k+1}^* \partial_{k+1} \quad \text{(Laplaciano up)} \\
    \Delta_k^{\text{down}} &= \partial_k \partial_k^* \quad \text{(Laplaciano down)}
\end{align}
donde $\partial_k^*$ es el adjunto de $\partial_k$.
\end{definition}

\begin{theorem}[Descomposición de Hodge Discreta]\label{thm:descomposicion_hodge_discreta}
El espacio de $k$-cadenas se descompone ortogonalmente como:
\begin{equation}
    C_k(\mathcal{K}) = B_k(\mathcal{K}) \oplus \mathcal{H}_k(\mathcal{K}) \oplus B_k^*(\mathcal{K})
\end{equation}
donde $\mathcal{H}_k(\mathcal{K}) = \ker(\Delta_k)$ es el espacio de formas armónicas (\autoref{def:formas_harmonicas_discretas}) y $B_k^*(\mathcal{K}) = \operatorname{im}(\partial_k^*)$ son las co-fronteras. La demostración se encuentra en el \autoref{thm:hodge_decomposition_gdl} del \autoref{ch:gdl}.
\end{theorem}

\begin{definition}[Formas Harmónicas Discretas]\label{def:formas_harmonicas_discretas}
Una $k$-cadena $\omega \in C_k(\mathcal{K})$ es harmónica si:
\begin{equation}
    \Delta_k \omega = 0
\end{equation}
El espacio de formas harmónicas es $\mathcal{H}_k(\mathcal{K}) = \ker(\Delta_k)$.
\end{definition}

\begin{theorem}[Isomorfismo Harmónico-Homología]\label{thm:isomorfismo_harmonico_homologia}
Para complejos simpliciales finitos:
\begin{equation}
    \mathcal{H}_k(\mathcal{K}) \cong H_k(\mathcal{K})
\end{equation}
\end{theorem}

\begin{proof}
Por la descomposición de Hodge (\autoref{thm:descomposicion_hodge_discreta}), $C_k = B_k \oplus \mathcal{H}_k \oplus B_k^*$. Los ciclos son $Z_k = \ker(\partial_k) = B_k \oplus \mathcal{H}_k$ (ya que $\partial_k$ se anula en $B_k$ y $\mathcal{H}_k$, pero actúa inyectivamente en $B_k^*$). Entonces $H_k = Z_k / B_k = (B_k \oplus \mathcal{H}_k) / B_k \cong \mathcal{H}_k$, donde el isomorfismo envía cada clase $[\omega]$ a su componente armónica.
\end{proof}

\subsection{Simplicial Attention Networks: Formulación Completa}

\subsubsection{Modelado de Datos como Complejos Simpliciales}

Los datos estructurados pueden representarse naturalmente como complejos simpliciales, donde las interacciones de orden superior se codifican explícitamente (véase la \autoref{subsec:tdl_grafos} del \autoref{ch:gdl} para los fundamentos del Laplaciano de Hodge y el paso de mensajes simplicial que motivan esta sección).

\begin{definition}[Dataset Simplicial]\label{def:dataset_simplicial}
Un dataset simplicial $\mathcal{D} = (\mathcal{K}, \mathbf{X})$ consiste en:
\begin{itemize}
    \item Un complejo simplicial $\mathcal{K}$ que codifica la estructura topológica
    \item Características de simplicios $\mathbf{X} = \{\mathbf{x}_{\sigma}\}_{\sigma \in \mathcal{K}}$ donde $\mathbf{x}_{\sigma} \in \mathbb{R}^{d_{\dim(\sigma)}}$
\end{itemize}
\end{definition}

\begin{example}[Aplicaciones de Datasets Simpliciales]\label{ex:aplicaciones_datasets_simpliciales}
\leavevmode
\begin{itemize}
    \item \textbf{Redes sociales}: Vértices = personas, aristas = amistades, triángulos = grupos de amigos
    \item \textbf{Proteínas}: Átomos = vértices, enlaces = aristas, anillos = ciclos
    \item \textbf{Cerebro}: Neuronas = vértices, sinapsis = aristas, circuitos = simplicios superiores
    \item \textbf{Grafos de citaciones}: Papers = vértices, citas = aristas, co-citaciones = triángulos
\end{itemize}
\end{example}

\subsubsection{Arquitectura de Simplicial Attention Networks}

\begin{figure}[ht]
\centering
\includegraphics[width=0.8\textwidth]{figures/simplicial_messages.pdf}
\caption{Tres tipos de mensajes en Simplicial Attention Networks. Una arista objetivo (roja) recibe: mensajes \emph{lower} desde sus vértices frontera (verde), mensajes \emph{upper} desde triángulos de los que es cara (naranja), y mensajes \emph{horizontales} desde aristas adyacentes (púrpura).}
\label{fig:simplicial_messages}
\end{figure}

\begin{definition}[Mensaje Simplicial]\label{def:mensaje_simplicial}
Para un $k$-simplicio $\sigma$, definimos diferentes tipos de mensajes:

\textbf{Mensaje Lower}: De $(k-1)$-simplicios incidentes:
\begin{equation}
    \mathbf{m}_{\sigma}^{\text{lower}} = \sum_{\tau \in \partial_k(\sigma)} \alpha_{\tau \to \sigma} \mathbf{W}^{\text{lower}} \mathbf{h}_{\tau}
\end{equation}

\textbf{Mensaje Upper}: De $(k+1)$-simplicios incidentes:
\begin{equation}
    \mathbf{m}_{\sigma}^{\text{upper}} = \sum_{\rho: \sigma \in \partial_{k+1}(\rho)} \alpha_{\rho \to \sigma} \mathbf{W}^{\text{upper}} \mathbf{h}_{\rho}
\end{equation}

\textbf{Mensaje Horizontal}: De $k$-simplicios adyacentes:
\begin{equation}
    \mathbf{m}_{\sigma}^{\text{horizontal}} = \sum_{\tau \in \mathcal{N}_k(\sigma)} \alpha_{\tau \to \sigma} \mathbf{W}^{\text{horizontal}} \mathbf{h}_{\tau}
\end{equation}
\end{definition}

\begin{definition}[Pesos de Atención Simplicial]\label{def:pesos_atencion_simplicial}
Los pesos de atención incorporan información de orientación:
\begin{align}
    \alpha_{\tau \to \sigma} &= \frac{\exp(\text{LeakyReLU}(\mathbf{a}^T [\mathbf{h}_{\tau} \| \mathbf{h}_{\sigma} \| \mathbf{e}_{\text{orient}}(\tau,\sigma)]))}{\sum_{\rho} \exp(\text{LeakyReLU}(\mathbf{a}^T [\mathbf{h}_{\rho} \| \mathbf{h}_{\sigma} \| \mathbf{e}_{\text{orient}}(\rho,\sigma)]))}
\end{align}
donde $\mathbf{e}_{\text{orient}}(\tau,\sigma)$ codifica la relación de orientación entre simplicios.
\end{definition}

\begin{definition}[Actualización Simplicial Completa]\label{def:actualizacion_simplicial}
La actualización de un $k$-simplicio integra todos los tipos de mensaje:
\begin{align}
    \mathbf{h}_{\sigma}^{(l+1)} = \sigma\left(\mathbf{W}^{\text{self}} \mathbf{h}_{\sigma}^{(l)} + \mathbf{m}_{\sigma}^{\text{lower}} + \mathbf{m}_{\sigma}^{\text{upper}} + \mathbf{m}_{\sigma}^{\text{horizontal}}\right)
\end{align}
\end{definition}

\subsubsection{Preservación de Estructura Topológica}

\begin{theorem}[Equivarianza Simplicial]\label{thm:equivarianza_simplicial}
Las Simplicial Attention Networks son equivariantes bajo automorfismos de complejos simpliciales. Si $\phi: \mathcal{K} \to \mathcal{K}$ es un automorfismo simplicial, entonces:
\begin{equation}
    \text{SAN}(\phi(\mathbf{X})) = \phi(\text{SAN}(\mathbf{X}))
\end{equation}
\end{theorem}

\begin{proof}
Un automorfismo simplicial $\phi$ induce una permutación $P_k$ en el conjunto de $k$-simplicios que preserva las relaciones de incidencia: $\mathbf{B}_k P_{k-1} = P_k \mathbf{B}_k$. La actualización simplicial (\autoref{def:actualizacion_simplicial}) depende de las matrices de frontera a través de los Laplacianos $\mathbf{L}_k^{\text{up}}$ y $\mathbf{L}_k^{\text{down}}$. Dado que $P_k \mathbf{L}_k P_k^T = \mathbf{L}_k$ (los Laplacianos conmutan con automorfismos simpliciales), y los pesos de atención $\alpha_{\sigma\tau}$ dependen de las representaciones a través de funciones simétricas respecto a $\phi$, se tiene que $P_k \cdot \text{SAN}(\mathbf{X}) = \text{SAN}(P_k \cdot \mathbf{X})$ para cada dimensión $k$.
\end{proof}

\begin{remark}[Estructura de Cadenas en SANs]\label{rem:respeto_estructura_cadenas}
Para la componente lineal de una SAN (sin funciones de activación no lineales), la propiedad de mapa de cadenas
\begin{equation}
    \text{SAN}_{\text{lin}}(\partial_k \mathbf{c}) = \partial_k \text{SAN}_{\text{lin}}(\mathbf{c})
\end{equation}
se cumple si la capa lineal conmuta con el operador de frontera $\partial_k$. Sin embargo, esta propiedad es generalmente falsa para la red completa con funciones de activación no lineales: las no linealidades elemento a elemento violan la condición de mapa de cadenas. En la práctica, las arquitecturas SANs están diseñadas para preservar aproximadamente la estructura topológica (procesando señales de cada grado simplicial por separado y respetando adyacencias), pero la conmutatividad exacta con $\partial_k$ no está garantizada.
\end{remark}

\begin{remark}[Preservación Aproximada de Homología]\label{cor:preservacion_homologia}
Si una red fuese un mapa de cadenas exacto (es decir, conmutase con $\partial_k$), entonces induciría una aplicación bien definida en homología: ciclos se mapearían a ciclos y fronteras a fronteras. Las SANs reales, al incorporar no linealidades, solo preservan la estructura homológica de forma aproximada. El diseño ``Hodge-aware'' de estas arquitecturas (que procesa las componentes de la descomposición de Hodge por separado) busca mantener esta propiedad en la medida de lo posible.
\end{remark}

\subsection{Hodge Laplacian Attention}

\subsubsection{Motivación desde la Geometría Diferencial}

En variedades diferenciables, el Laplaciano de Hodge $\Delta = d \delta + \delta d$ (donde $d$ es la derivada exterior y $\delta$ su adjunto) juega un papel central en la teoría de formas diferenciables. Su análogo discreto proporciona una poderosa herramienta para el análisis de señales en complejos simpliciales.

\begin{definition}[Hodge Laplacian Discreto]\label{def:hodge_laplacian_discreto}
Para un complejo simplicial $\mathcal{K}$ con producto interno, el Laplaciano de Hodge de orden $k$ es:
\begin{align}
    \mathbf{L}_k &= \mathbf{B}_{k+1} \mathbf{B}_{k+1}^T + \mathbf{B}_k^T \mathbf{B}_k \\
    &= \mathbf{L}_k^{\text{up}} + \mathbf{L}_k^{\text{down}}
\end{align}
donde $\mathbf{B}_k$ es la matriz de frontera del complejo.
\end{definition}

\begin{theorem}[Descomposición Espectral del Hodge Laplacian]\label{thm:descomposicion_hodge_laplacian}
El Laplacian de Hodge admite la descomposición:
\begin{equation}
    \mathbf{L}_k = \mathbf{U}_k \boldsymbol{\Lambda}_k \mathbf{U}_k^T
\end{equation}
donde $\boldsymbol{\Lambda}_k = \text{diag}(\lambda_1, \lambda_2, \ldots, \lambda_{n_k})$ con $0 = \lambda_1 \leq \lambda_2 \leq \cdots \leq \lambda_{n_k}$.
\end{theorem}

\begin{proof}
El Hodge Laplacian $\mathbf{L}_k = \mathbf{B}_{k+1}\mathbf{B}_{k+1}^T + \mathbf{B}_k^T\mathbf{B}_k$ es simétrico y semidefinido positivo: para cualquier vector $\mathbf{x}$, $\mathbf{x}^T \mathbf{L}_k \mathbf{x} = \|\mathbf{B}_{k+1}^T \mathbf{x}\|^2 + \|\mathbf{B}_k \mathbf{x}\|^2 \geq 0$. Por el teorema espectral para matrices simétricas reales, $\mathbf{L}_k$ admite una base ortonormal de vectores propios con eigenvalores reales no negativos. La multiplicidad de $\lambda = 0$ es $\dim(\ker(\mathbf{L}_k)) = \beta_k$, el $k$-ésimo número de Betti.
\end{proof}

\textbf{Interpretación de los Autovalores.}\label{rem:interpretacion_autovalores_hodge}
\begin{itemize}
    \item $\lambda_1 = 0$ con multiplicidad $\beta_k$ (número de Betti)
    \item Los autovectores correspondientes a $\lambda_1$ son las formas harmónicas
    \item Autovalores más grandes corresponden a señales con mayor ``curvatura''
\end{itemize}

\subsubsection{Hodge Laplacian Attention Networks}

\begin{definition}[Filtro Espectral de Hodge]\label{def:filtro_espectral_hodge}
Un filtro espectral en el complejo simplicial se define como:
\begin{equation}
    \mathbf{g}_{\theta}(\mathbf{L}_k) = \sum_{i=0}^{K-1} \theta_i \mathbf{T}_i(\tilde{\mathbf{L}}_k)
\end{equation}
donde $\mathbf{T}_i$ son polinomios de Chebyshev y $\tilde{\mathbf{L}}_k = \frac{2\mathbf{L}_k}{\lambda_{\max}} - \mathbf{I}$.
\end{definition}

\begin{definition}[Hodge Laplacian Attention]\label{def:hodge_laplacian_attention}
La atención basada en Hodge Laplacian incorpora información espectral:
\begin{align}
    \mathbf{Q}_k &= \mathbf{g}_{\theta}^Q(\mathbf{L}_k) \mathbf{H}_k \mathbf{W}^Q \\
    \mathbf{K}_k &= \mathbf{g}_{\theta}^K(\mathbf{L}_k) \mathbf{H}_k \mathbf{W}^K \\
    \mathbf{V}_k &= \mathbf{g}_{\theta}^V(\mathbf{L}_k) \mathbf{H}_k \mathbf{W}^V
\end{align}

El mecanismo de atención se define como:
\begin{equation}
    \text{HodgeAttention}(\mathbf{H}_k) = \text{softmax}\left(\frac{\mathbf{Q}_k \mathbf{K}_k^T}{\sqrt{d_k}}\right) \mathbf{V}_k
\end{equation}
\end{definition}

\begin{theorem}[Interpretación Frecuencial de Hodge Attention]\label{thm:interpretacion_frecuencial_hodge}
La Hodge Laplacian Attention actúa como un filtro en el dominio espectral del complejo simplicial:
\begin{equation}
    \widehat{\text{HodgeAttention}}(\lambda) = g_{\theta}^Q(\lambda) g_{\theta}^K(\lambda) \hat{\mathbf{h}}(\lambda)
\end{equation}
donde $\hat{\mathbf{h}}(\lambda)$ es la transformada espectral de las características.
\end{theorem}

\begin{proof}
Proyectando sobre la base de eigenvectores $\{\mathbf{u}_i\}$ del Hodge Laplacian (\autoref{thm:descomposicion_hodge_laplacian}), los filtros espectrales actúan diagonalmente: $g_\theta(\mathbf{L}_k) = \mathbf{U}_k g_\theta(\boldsymbol{\Lambda}_k) \mathbf{U}_k^T$. Entonces $\mathbf{Q}_k = g_\theta^Q(\mathbf{L}_k) \mathbf{H}_k \mathbf{W}^Q$ tiene transformada espectral $\hat{\mathbf{Q}}_k(\lambda_i) = g_\theta^Q(\lambda_i) \hat{\mathbf{H}}_k(\lambda_i) \mathbf{W}^Q$, y análogamente para $\mathbf{K}_k$. El score de atención en el dominio espectral es $\hat{\mathbf{Q}}_k^T \hat{\mathbf{K}}_k(\lambda_i) \propto g_\theta^Q(\lambda_i) g_\theta^K(\lambda_i)$, que modula la señal original $\hat{\mathbf{h}}(\lambda_i)$. El mecanismo completo actúa como un filtro pasa-banda cuya función de transferencia es el producto de los filtros de Q y K.
\end{proof}

\paragraph{Atención Espectral Topológica}

Los eigenvalores del Hodge Laplacian (\autoref{thm:descomposicion_hodge_laplacian}) proporcionan características espectrales que capturan propiedades topológicas. La atención espectral topológica incorpora estas características mediante un sesgo espectral aprendible:

\begin{equation}
\operatorname{SpectralAttn}_k = \operatorname{softmax}\left(\frac{\mathbf{Q}_k \mathbf{K}_k^T}{\sqrt{d_k}} + \mathbf{W}_{\lambda} \operatorname{diag}(\boldsymbol{\lambda}^{(k)})\right) \mathbf{V}_k
\label{eq:spectral_topological_attention}
\end{equation}

\paragraph{Filtrado en Componentes de Hodge}

La descomposición de Hodge permite filtrar información separando las señales en componentes exactas, armónicas y co-exactas:

\begin{align}
\mathbf{H}_{\text{exact}} &= \mathbf{P}_{\operatorname{Im}(\delta_{k-1})} \mathbf{H} \label{eq:exact_component} \\
\mathbf{H}_{\text{harmonic}} &= \mathbf{P}_{\operatorname{Ker}(\mathcal{L}_k)} \mathbf{H} \\
\mathbf{H}_{\text{coexact}} &= \mathbf{P}_{\operatorname{Im}(\delta_k^*)} \mathbf{H}
\end{align}
donde $\mathbf{P}$ denota proyecciones ortogonales a cada subespacio.

\subsection{Aplicaciones en Análisis de Redes y Grafos Complejos}

\subsubsection{Análisis Topológico de Datos (TDA)}

\begin{definition}[Pipeline TDA para Deep Learning]\label{def:pipeline_tda_dl}
El pipeline integrado TDA-Deep Learning consiste en:
\begin{enumerate}
    \item \textbf{Construcción de complejo}: A partir de datos, construir $\mathcal{K}$ (Vietoris-Rips, Alpha, etc.)
    \item \textbf{Cálculo de persistencia}: Obtener diagramas $\text{Pers}_k(\mathcal{K})$
    \item \textbf{Vectorización}: Convertir diagramas a representaciones vectoriales
    \item \textbf{Aprendizaje}: Usar como características para modelos de atención
\end{enumerate}
\end{definition}

\begin{definition}[Representaciones Vectoriales de Persistencia]\label{def:representaciones_persistencia}
Métodos comunes incluyen:

\textbf{Persistence Images} \citep{adams2017persistence}:
\begin{equation}
    \rho_{\sigma,u}(x,y) = \sum_{(b,d) \in \text{Dgm}} \omega(b,d) \cdot \frac{1}{\sigma\sqrt{2\pi}} \exp\left(-\frac{(x-b)^2 + (y-(d-b))^2}{2\sigma^2}\right)
\end{equation}

\textbf{Persistence Landscapes}:
\begin{equation}
    \lambda_k(t) = k\text{-ésimo mayor valor de } \min(t-b, d-t)
\end{equation}
para puntos $(b,d)$ en el diagrama.
\end{definition}

\begin{algorithm}
\caption{Simplicial Attention con TDA}\label{alg:sat_tda}
\textbf{Entrada}: Dataset $\mathcal{X}$, parámetros de filtración $\epsilon_1 < \epsilon_2 < \cdots < \epsilon_m$

\begin{enumerate}
    \item \textbf{Construir filtración}:
    \begin{align}
        K_i &= \text{VietorisRips}(\mathcal{X}, \epsilon_i) \quad \forall i
    \end{align}
    
    \item \textbf{Calcular homología persistente}:
    \begin{align}
        \text{Dgm}_k &= \text{PersistentHomology}(K_0 \subseteq K_1 \subseteq \cdots \subseteq K_m)
    \end{align}
    
    \item \textbf{Generar características topológicas}:
    \begin{align}
        \mathbf{f}_{\text{topo}} &= \text{PersistenceImage}(\text{Dgm}_k)
    \end{align}
    
    \item \textbf{Aplicar Simplicial Attention}:
    \begin{align}
        \mathbf{h}^{(l+1)} &= \text{SAN}([\mathbf{h}^{(l)} \| \mathbf{f}_{\text{topo}}])
    \end{align}
\end{enumerate}

\textbf{Salida}: Representaciones enriquecidas topológicamente
\end{algorithm}

\subsubsection{Aplicaciones Específicas}

\begin{application}[Análisis de Redes Cerebrales]\label{app:redes_cerebrales}
En conectómica, los datos cerebrales se modelan como complejos simpliciales donde:
\begin{itemize}
    \item \textbf{Vértices}: Regiones cerebrales (ROIs)
    \item \textbf{Aristas}: Conexiones funcionales por pares
    \item \textbf{Triángulos}: Interacciones funcionales de orden 3
    \item \textbf{k-simplicios}: Patrones de activación de orden superior
\end{itemize}

La Hodge Laplacian Attention permite detectar patrones anómalos en:
\begin{align}
    \mathbf{A}_{\text{brain}} &= \text{HodgeAttention}(\mathbf{H}_{\text{roi}}, \mathbf{L}_{\text{connectome}})
\end{align}
\end{application}

\begin{application}[Análisis de Redes Sociales]\label{app:redes_sociales}
Para detectar comunidades y estructuras jerárquicas:
\begin{enumerate}
    \item Construir complejo de cliques del grafo social
    \item Aplicar filtración por fuerza de conexión
    \item Usar persistencia para identificar comunidades estables
    \item Aplicar Simplicial Attention para clasificación de usuarios
\end{enumerate}
\end{application}

\begin{theorem}[Poder Expresivo de SANs vs GATs]\label{thm:poder_expresivo_sans}
Para grafos con estructura topológica rica, las Simplicial Attention Networks son estrictamente más expresivas que las Graph Attention Networks tradicionales:
\begin{equation}
    \exists G: \quad \text{SAN}(G) \not\in \text{span}\{\text{GAT}(G'): G' \text{ es grafo}\}
\end{equation}
\end{theorem}

\begin{proof}
Construimos un ejemplo separador. Consideremos dos grafos $G_1$ y $G_2$ con el mismo grafo 1-esqueleto (mismos nodos y aristas) pero diferentes complejos simpliciales de orden superior: $G_1$ contiene un 2-simplicio (triángulo relleno) $[v_1, v_2, v_3]$ y $G_2$ no (solo tiene las tres aristas). Cualquier \gls{GAT} opera exclusivamente sobre nodos y aristas, por lo que produce la misma salida para $G_1$ y $G_2$ (ya que el 1-esqueleto es idéntico). Sin embargo, una SAN accede a los 2-simplicios: la señal en el triángulo relleno de $G_1$ genera mensajes tipo ``upper'' (\autoref{def:actualizacion_simplicial}) que no existen en $G_2$. En particular, $\beta_1(G_1) \neq \beta_1(G_2)$ (el triángulo relleno elimina un ciclo independiente), y esta diferencia topológica es detectable por las componentes armónicas del Hodge Laplacian de orden 1.
\end{proof}

\begin{corollary}[Ventaja Computacional]\label{cor:ventaja_computacional_sans}
Para problemas con interacciones de orden superior, las SANs requieren menos parámetros que las \glspl{GAT} multicapa equivalentes:
\begin{equation}
    |\theta_{\text{SAN}}| = O(d \cdot k_{\max}) \ll |\theta_{\text{GAT-equiv}}| = O(d \cdot L^2)
\end{equation}
donde $k_{\max}$ es la dimensión máxima de simplicios y $L$ es la profundidad de \gls{GAT}.
\end{corollary}

\begin{proof}
Una SAN procesa directamente $k$-simplicios para cada dimensión $k = 0, 1, \ldots, k_{\max}$, con $O(d)$ parámetros por dimensión simplicial (matrices de atención y transformación de características). Esto da un total de $O(d \cdot k_{\max})$ parámetros.

Una \gls{GAT} estándar solo opera sobre aristas (1-simplicios), por lo que capturar interacciones de orden $k$ requiere propagar información a través de $k$ saltos. Para capturar interacciones hasta orden $k_{\max}$, se necesitan al menos $L \geq k_{\max}$ capas, cada una con $O(d \cdot L)$ parámetros (las dimensiones intermedias deben crecer con $L$ para evitar pérdida de información). El total es $O(d \cdot L^2)$. Como las SANs acceden directamente a las interacciones de orden superior, evitan la acumulación cuadrática de parámetros.
\end{proof}

\subsection{Conexiones con Modelos de Difusión y Ecuaciones Diferenciales}

\subsubsection{Difusión en Complejos Simpliciales}

\begin{definition}[Ecuación de Difusión Simplicial]\label{def:ecuacion_difusion_simplicial}
La ecuación de difusión en un complejo simplicial se escribe como:
\begin{equation}
    \frac{\partial \mathbf{f}_k}{\partial t} = -\mathbf{L}_k \mathbf{f}_k
\end{equation}
donde $\mathbf{f}_k(t) \in \mathbb{R}^{n_k}$ representa señales en $k$-simplicios.
\end{definition}

\begin{theorem}[Solución de la Ecuación de Difusión]\label{thm:solucion_difusion_simplicial}
La solución de la ecuación de difusión es:
\begin{equation}
    \mathbf{f}_k(t) = \exp(-t\mathbf{L}_k) \mathbf{f}_k(0) = \sum_{i=1}^{n_k} \exp(-t\lambda_i) \langle \mathbf{u}_i, \mathbf{f}_k(0) \rangle \mathbf{u}_i
\end{equation}
\end{theorem}

\begin{proof}
La ecuación $\dot{\mathbf{f}}_k = -\mathbf{L}_k \mathbf{f}_k$ es un sistema lineal de EDOs con solución $\mathbf{f}_k(t) = \exp(-t\mathbf{L}_k) \mathbf{f}_k(0)$, donde la exponencial de matriz se define por su serie $\exp(A) = \sum_{n=0}^\infty A^n/n!$. Usando la descomposición espectral $\mathbf{L}_k = \mathbf{U}_k \boldsymbol{\Lambda}_k \mathbf{U}_k^T$ (\autoref{thm:descomposicion_hodge_laplacian}): $\exp(-t\mathbf{L}_k) = \mathbf{U}_k \exp(-t\boldsymbol{\Lambda}_k) \mathbf{U}_k^T = \mathbf{U}_k \operatorname{diag}(e^{-t\lambda_1}, \ldots, e^{-t\lambda_{n_k}}) \mathbf{U}_k^T$. Expandiendo $\mathbf{f}_k(0) = \sum_i \langle \mathbf{u}_i, \mathbf{f}_k(0)\rangle \mathbf{u}_i$, se obtiene el resultado. Las componentes armónicas ($\lambda_i = 0$) son estacionarias, mientras que las de alta frecuencia decaen exponencialmente.
\end{proof}

\begin{definition}[Neural Diffusion en Complejos]\label{def:neural_diffusion_complejos}
Una capa de difusión neural en complejos simpliciales:
\begin{align}
    \mathbf{H}_k^{(l+1)} &= \sigma\left(\mathbf{W}^{(l)} \exp(-\alpha^{(l)} \mathbf{L}_k) \mathbf{H}_k^{(l)}\right)
\end{align}
donde $\alpha^{(l)}$ controla la escala de difusión.
\end{definition}

Este marco unifica conceptos de geometría diferencial, topología algebraica y deep learning, proporcionando herramientas matemáticas sofisticadas para el análisis de datos con estructura topológica compleja.

Las herramientas matemáticas desarrolladas en las secciones precedentes (atención equivariante, espectral, hiperbólica, Clifford y topológica) encuentran aplicaciones directas en ciencia molecular, física computacional y biología estructural, donde la preservación de simetrías como $E(3)$, $SE(3)$ o el grupo de Lorentz es esencial para la predicción de propiedades físicas. Para una visión general de estos dominios de aplicación, véase la introducción (\autoref{ch:intro}).

\subsection[Transformers en Complejos CW]{\glspl{Transformer} en Complejos CW}

Los complejos celulares representan una generalización natural de grafos, permitiendo capturar relaciones de orden superior mediante células de diferentes dimensiones \citep{bodnar2021cell}.

\paragraph{Atención Multi-Dimensional}

En un complejo CW, la atención debe operar sobre múltiples tipos de células:

\begin{align}
\text{CWAttention}(\mathbf{X}) &= \bigoplus_{k=0}^{d} \text{Attention}^{(k)}(\mathbf{X}_k) \label{eq:cw_attention} \\
\text{Attention}^{(k)}(\mathbf{X}_k) &= \text{softmax}\left(\frac{\mathbf{Q}_k \mathbf{K}_k^T}{\sqrt{d_k}} + \mathbf{A}^{(k)}\right) \mathbf{V}_k
\end{align}

donde $\mathbf{X}_k$ representa las características de $k$-células y $\mathbf{A}^{(k)}$ es la matriz de adyacencia entre $k$-células.

\paragraph{Preservación de Relaciones de Incidencia}

Las relaciones de incidencia entre células de diferentes dimensiones se preservan mediante:

\begin{equation}
\mathbf{M}^{(k,k-1)} = \text{IncidenceMatrix}(\text{Cell}_k, \text{Cell}_{k-1})
\label{eq:incidence_preservation}
\end{equation}

La atención inter-dimensional incorpora estas relaciones:

\begin{equation}
\text{InterAttn}_{k \to k-1} = \text{softmax}(\mathbf{Q}_{k-1} (\mathbf{M}^{(k,k-1)} \mathbf{K}_k)^T) \mathbf{M}^{(k,k-1)} \mathbf{V}_k
\label{eq:inter_dimensional_attention}
\end{equation}

\subsection{Persistent Homology Attention}

La homología persistente proporciona una descripción multiescala de características topológicas, capturando cómo estas evolucionan a través de escalas espaciales.

\paragraph{Diagramas de Persistencia}

Para una filtración $X_0 \subseteq X_1 \subseteq \cdots \subseteq X_n$, los diagramas de persistencia registran el tiempo de nacimiento y muerte de características topológicas:

\begin{equation}
\text{Dgm}_k = \{(b_i, d_i) : H_k(X_{b_i}) \to H_k(X_{d_i}) \text{ es isomorfismo}\}
\label{eq:persistence_diagram}
\end{equation}

\paragraph{Integración en Atención}

Los diagramas de persistencia se integran en el mecanismo de atención mediante:

\begin{equation}
\alpha_{ij} = \text{softmax}\left(\frac{\mathbf{q}_i^T \mathbf{k}_j}{\sqrt{d}} + w \cdot \text{pers}(i,j)\right)
\label{eq:persistent_attention_weights}
\end{equation}

donde $\text{pers}(i,j)$ codifica la persistencia topológica entre puntos $i$ y $j$.

\paragraph{Funciones de Persistencia}

La función de persistencia puede tomar diversas formas:

\begin{align}
\text{pers}_{\text{lifetime}}(i,j) &= \sum_{(b,d) \in \text{Dgm}} (d-b) \cdot \mathbb{I}[\text{connected}(i,j,b,d)] \label{eq:lifetime_persistence} \\
\text{pers}_{\text{birth}}(i,j) &= \min\{b : (b,d) \in \text{Dgm}, \text{connected}(i,j,b,d)\} \\
\text{pers}_{\text{stability}}(i,j) &= \max_{(b,d) \in \text{Dgm}} (d-b) \cdot \mathbb{I}[\text{connected}(i,j,b,d)]
\end{align}

\paragraph{Ventajas de la Persistencia}

La incorporación de homología persistente ofrece:

\begin{enumerate}
    \item \textbf{Robustez a ruido}: Las características persistentes son más estables
    \item \textbf{Detección multiescala}: Captura patrones a múltiples resoluciones
    \item \textbf{Invarianza topológica}: Resistencia a deformaciones continuas
    \item \textbf{Interpretabilidad}: Conexión directa con propiedades geométricas conocidas
\end{enumerate}

\subsection{Aplicaciones de Álgebra Topológica}

Esta formulación avanzada encuentra aplicaciones en:

\begin{itemize}
    \item \textbf{Análisis de flujos en redes}: Separación de componentes irrotacionales y solenoidales usando Hodge
    \item \textbf{Procesamiento electromagnético}: Campos eléctricos y magnéticos como formas diferenciales
    \item \textbf{Topología computacional}: Análisis de cavidades y agujeros en datos complejos
    \item \textbf{Análisis de series temporales}: Captura de patrones topológicos recurrentes con persistencia
    \item \textbf{Detección de anomalías}: Identificación de cambios topológicos inesperados
    \item \textbf{Sistemas dinámicos}: Caracterización de atractores y órbitas mediante topología
\end{itemize}


\subsection{Desafíos Teóricos y Prácticos Pendientes}

Pese a los avances significativos, los mecanismos de atención geométrica enfrentan desafíos fundamentales que requieren investigación continua.

\subsubsection{Limitaciones de Expresividad vs Eficiencia}

\paragraph{El Dilema de la Complejidad}

El trade-off entre expresividad y eficiencia computacional permanece como uno de los desafíos centrales. Mientras que la atención tradicional logra $O(n^2)$ complejidad, los enfoques geométricos frecuentemente incrementan la complejidad debido a:

\begin{enumerate}
    \item \textbf{Operaciones en grupos de Lie}: El costo de proyecciones y lifting a espacios de alta dimensión
    \item \textbf{Análisis espectral}: Eigendecomposición de Laplacianos en grafos grandes
    \item \textbf{Transporte paralelo}: Cálculos geodésicos en variedades complejas
    \item \textbf{Productos tensoriales}: Expansión dimensional en representaciones irreducibles
\end{enumerate}

\begin{definition}[Índice de Complejidad Geométrica]
Para una arquitectura de atención geométrica con simetría $G$, definimos:
\begin{equation}
\text{ICG}(G, n) = \frac{\text{Coste}(G\text{-equivarianza})}{\text{Coste}(\text{atención estándar})} \cdot \frac{|\text{Orbit}(G)|}{n}
\label{eq:geometric_complexity_index}
\end{equation}
donde $|\text{Orbit}(G)|$ mide el tamaño de las órbitas bajo la acción del grupo.
\end{definition}

\paragraph{Aproximaciones y Compromisos}

Los desarrollos recientes en arquitecturas sub-cuadráticas sugieren que es posible preservar propiedades geométricas esenciales con complejidad lineal:

\begin{remark}[Preservación Asintótica de Equivarianza]
Bajo ciertas condiciones de regularidad, se conjetura que los modelos de complejidad lineal pueden aproximar arquitecturas equivariantes con error acotado:
\begin{equation}
\|f_{\text{lineal}}(g \cdot x) - g \cdot f_{\text{cuadr\'atico}}(x)\| \leq \epsilon_{\text{precisi\'on}}
\label{eq:linear_equivariance_approximation}
\end{equation}
La formalización rigurosa de esta cota y las condiciones precisas bajo las que se cumple permanecen como un problema abierto.
\end{remark}

\subsubsection{Escalabilidad a Datos Masivos}

\paragraph{Limitaciones de Memoria}

Los modelos fundacionales modernos requieren procesamiento de:
\begin{itemize}
    \item \textbf{Estructuras proteicas}: $>10^6$ átomos por complejo macromolecular
    \item \textbf{Simulaciones moleculares}: $10^9$ partículas en dinámicas temporalmente extensas
    \item \textbf{Grafos de redes sociales}: $10^{12}$ conexiones en plataformas globales
    \item \textbf{Nubes de puntos LiDAR}: $10^8$ puntos por escaneo de ciudades
\end{itemize}

La memoria GPU requerida escala según:

\begin{equation}
\text{Memoria}_{\text{GDL}} = O(n^2 \cdot d_{\text{modelo}} \cdot |\text{Rep}(G)| \cdot \text{Precisión})
\label{eq:gdl_memory_scaling}
\end{equation}

donde $|\text{Rep}(G)|$ representa la dimensión de las representaciones del grupo.

\paragraph{Estrategias de Distribución}

Las soluciones emergentes incluyen:

\begin{enumerate}
    \item \textbf{Ring Attention distribuida}: Paralelización en anillo con overlapping de comunicación
    \item \textbf{Gradient checkpointing geométrico}: Recompute selectivo de operaciones equivariantes costosas
    \item \textbf{Quantización consciente de simetrías}: Preservación de equivarianza bajo precisión reducida
    \item \textbf{Sparsificación estructurada}: Pruning que respeta la estructura del grupo
\end{enumerate}

\subsubsection{Interpretabilidad de Atención Geométrica}

\paragraph{Complejidad de Visualización}

La interpretabilidad de mecanismos de atención geométrica presenta desafíos únicos:

\begin{itemize}
    \item \textbf{Espacios de alta dimensión}: Visualización de representaciones en $SO(3)$ o $SE(3)$
    \item \textbf{Múltiples escalas}: Atención operando simultáneamente en diversas resoluciones geométricas
    \item \textbf{Interacciones no-lineales}: Productos tensoriales y operaciones de grupo complejas
    \item \textbf{Dependencias temporales}: Evolución de patrones de atención en secuencias dinámicas
\end{itemize}

\begin{definition}[Saliencia Geométrica]
Para un mecanismo de atención $G$-equivariante, la saliencia respecto al grupo se define como:
\begin{equation}
\text{Sal}_G(x_i, x_j) = \sum_{g \in G} w(g) \cdot \alpha_{ij}(g \cdot x_i, g \cdot x_j)
\label{eq:geometric_saliency}
\end{equation}
donde $w(g)$ son pesos asociados a elementos del grupo.
\end{definition}

\paragraph{Herramientas de Análisis}

El desarrollo de herramientas especializadas incluye:

\begin{enumerate}
    \item \textbf{Proyección de variedades}: Visualización de atención en espacios 2D/3D interpretables
    \item \textbf{Análisis espectral de patrones}: Descomposición de attention maps en componentes frecuenciales
    \item \textbf{Trazabilidad de simetrías}: Seguimiento de preservación de equivarianza a través de capas
    \item \textbf{Perturbación dirigida}: Análisis de sensibilidad respecto a transformaciones geométricas específicas
\end{enumerate}

\subsubsection{Generalización Teórica}

\paragraph{Límites de Aproximación Universal}

Los resultados de aproximación universal para atención geométrica requieren extensión:

\begin{remark}[Aproximación Universal Condicionada por Grupo]
Sea $\mathcal{F}_G$ el espacio de funciones $G$-equivariantes en un dominio geométrico $\mathcal{M}$. Se conjetura que una red de atención geométrica puede aproximar uniformemente cualquier $f \in \mathcal{F}_G$ bajo las siguientes condiciones:
\begin{enumerate}
    \item La acción de $G$ en $\mathcal{M}$ es transitiva por partes
    \item Las representaciones utilizadas generan todas las representaciones irreducibles necesarias
    \item El mecanismo de atención respeta la estructura causal del dominio
\end{enumerate}
Este resultado extiende la conjetura de aproximación universal para funciones equivariantes (véase \autoref{rem:aproximacion_universal_clifford}) y permanece como un problema abierto.
\end{remark}

\paragraph{Cotas de Generalización}

La capacidad de generalización está limitada por:

\begin{equation}
\text{Error}_{\text{gen}} \leq \text{Error}_{\text{emp}} + \sqrt{\frac{\text{Dim}(\text{Rep}(G)) \log N}{N}}
\label{eq:geometric_generalization_bound}
\end{equation}

donde $\text{Dim}(\text{Rep}(G))$ es la dimensión efectiva del espacio de representaciones.

\section{Conexiones con CNNs y Direcciones Futuras}\label{sec:cnn_attention_connections}

Este capítulo ha desarrollado los mecanismos de atención geométrica desde múltiples perspectivas matemáticas: teoría de grupos, análisis espectral, geometría diferencial, álgebra geométrica y topología algebraica. Para completar el marco teórico unificador del Geometric Deep Learning, esta sección final establece las conexiones fundamentales con las Redes Neuronales Convolucionales Geométricas desarrolladas en el \autoref{ch:cnn}, demostrando la convergencia conceptual de ambos paradigmas bajo principios matemáticos compartidos.

\subsection{Marco Teórico Unificador: Message Passing Geométrico}\label{subsec:marco_teorico_unificador}

Retomando la síntesis del \autoref{ch:gdl} (Sección~\ref{subsec:message_passing_unificador}) y conectando con el análisis profundo de CNNs geométricas, establecemos el siguiente marco unificador que subsume tanto convoluciones como atención bajo una abstracción matemática común.

\begin{definition}[Message Passing Geométrico Generalizado]\label{def:message_passing_geometrico_generalizado}
Sea $\mathcal{D}$ un dominio geométrico (grilla euclidiana, grafo, variedad, espacio homogéneo), $G$ un grupo de simetrías, y $\mathcal{F}(\mathcal{D})$ el espacio de señales sobre $\mathcal{D}$. Un operador de Message Passing Geométrico es una familia de transformaciones:
\begin{equation}
    \mathcal{T}^{(G)}: \mathcal{F}(\mathcal{D}) \to \mathcal{F}(\mathcal{D})
    \label{eq:message_passing_geometrico_generalizado}
\end{equation}
que satisface:
\begin{enumerate}
    \item \textbf{G-Equivarianza}: $\mathcal{T}^{(G)}(g \cdot f) = g \cdot \mathcal{T}^{(G)}(f)$ para todo $g \in G$
    \item \textbf{Localidad Geométrica}: Soporte determinado por la geometría de $\mathcal{D}$
    \item \textbf{Composicionalidad}: $\mathcal{T}^{(G)}_2 \circ \mathcal{T}^{(G)}_1$ preserva propiedades geométricas
\end{enumerate}
\end{definition}

\begin{theorem}[Correspondencia Fundamental CNN-Attention]\label{thm:correspondencia_cnn_attention}
Existe una correspondencia biyectiva entre convoluciones geométricas y mecanismos de atención geométrica a través del Message Passing Geométrico:
\begin{align}
    \text{\textbf{CNN Geométrica}} &\leftrightarrow \text{Message Passing con pesos fijos} \\
    \text{\textbf{Attention Geométrica}} &\leftrightarrow \text{Message Passing con pesos adaptativos} \\
    \text{\textbf{Híbrido}} &\leftrightarrow \text{Message Passing con estructura parcial}
\end{align}
Esta correspondencia preserva las propiedades de equivarianza, poder expresivo, y complejidad computacional.
\end{theorem}

\begin{proof}[Esquema de demostración]
La demostración se construye estableciendo la equivalencia formal en tres niveles:

\textbf{Nivel Algebraico}: Ambas operaciones se expresan como convoluciones generalizadas en el álgebra de grupo $\mathbb{C}[G]$. Para CNNs, el núcleo es fijo: $k_g$; para atención, es dinámico: $\alpha_{ij}(\mathbf{h}_i, \mathbf{h}_j)$.

\textbf{Nivel Espectral}: En la base de Fourier del dominio geométrico, ambas operaciones corresponden a multiplicación por funciones de transferencia. CNNs usan filtros espectrales fijos $\hat{k}(\lambda)$; atención usa filtros adaptativos $\hat{\alpha}(\lambda; \mathbf{h})$.

\textbf{Nivel Geométrico}: La estructura de vecindarios determina el soporte de integración. CNNs integran sobre vecindarios fijos; atención integra con pesos que dependen del contenido y respetan la geometría subyacente.
\end{proof}

\subsection{Tabla de Equivalencias Teóricas}\label{subsec:tabla_equivalencias}

El siguiente análisis sistemático establece las correspondencias directas entre conceptos fundamentales desarrollados en ambos capítulos:

\begin{table}[H]
\centering
\caption{Equivalencias Matemáticas Fundamentales: CNNs $\leftrightarrow$ Attention}
\label{tab:cnn_attention_equivalencias_fundamentales}
\begin{tabular}{|p{4cm}|p{4cm}|p{6cm}|}
\hline
\textbf{CNN Geométrica} & \textbf{Attention Geométrica} & \textbf{Principio Unificador} \\
\hline\hline
\textit{Convolución de Grupo} \\
(\autoref{sec:group_convolutions}) & \textit{G-Equivariant Attention} \\
(Sección~\ref{subsec:equivariant_attention}) & Operaciones que conmutan con la acción de grupo $G$ \\
\hline
\textit{Filtros Espectrales} \\
$\hat{f}(\lambda_i) = g(\lambda_i) \hat{x}(\lambda_i)$ & \textit{Spectral Attention} \\
$A_{ij} = \text{softmax}(f(\lambda_i, \lambda_j))$ & Multiplicación en dominio de Fourier del Laplaciano \\
\hline
\textit{Convolución Gauge} \\
(\autoref{sec:gauge_cnns}) & \textit{Gauge-Equivariant Attention} & Respeto a transformaciones locales de coordenadas \\
\hline
\textit{Steerable CNNs} \\
$f \mapsto \rho(g) f$ & \textit{Steerable Attention} \\
$A \mapsto \rho(g) A \rho(g)^{-1}$ & Acción coherente sobre campos vectoriales \\
\hline
\textit{Spherical CNNs} \\
Armónicos esféricos $Y_\ell^m$ & \textit{Spherical Attention} \\
Productos escalares en $L^2(S^2)$ & Análisis armónico en la esfera \\
\hline
\textit{Manifold CNNs} \\
Transporte paralelo & \textit{Riemannian Attention} \\
Conexiones riemannianas & Geometría intrínseca de variedades \\
\hline
\textit{ChebNet} \\
Polinomios de Chebyshev & \textit{Polynomial Attention} \\
Aproximación polinomial & Localización espectral eficiente \\
\hline
\end{tabular}
\end{table}

\begin{remark}[Jerarquía de Generalización Arquitectónica]\label{rem:jerarquia_generalizacion}
Las arquitecturas de Geometric Deep Learning forman una jerarquía natural de generalización:
\begin{equation}
\text{ANN} \subset \text{CNN} \subset \text{G-CNN} \subset \text{Manifold-CNN} \stackrel{\text{duales}}{\leftrightarrow} \text{Attention} \subset \text{G-Attention} \subset \text{Manifold-Attention}
\end{equation}
donde cada inclusión preserva propiedades geométricas fundamentales mientras incrementa la expresividad.
\end{remark}

\subsection{Análisis Comparativo Profundo: Rigidez vs. Adaptabilidad}\label{subsec:rigidez_adaptabilidad}

Una de las distinciones más fundamentales entre paradigmas convolucionales y de atención radica en el balance entre estructura impuesta (rigidez) y flexibilidad aprendible (adaptabilidad). Esta tensión se puede formalizar matemáticamente:

\begin{definition}[Espectro Rigidez-Adaptabilidad]\label{def:espectro_rigidez_adaptabilidad}
Para una arquitectura neuronal $\mathcal{A}$ operando sobre dominio $\mathcal{D}$, definimos:
\begin{align}
    \text{Rigidez}(\mathcal{A}) &= \frac{\dim(\text{Simetrías impuestas})}{\dim(\text{Espacio de parámetros})} \\
    \text{Adaptabilidad}(\mathcal{A}) &= \frac{\dim(\text{Pesos dependientes del contenido})}{\dim(\text{Espacio de parámetros})}
\end{align}
con la restricción fundamental: $\text{Rigidez}(\mathcal{A}) + \text{Adaptabilidad}(\mathcal{A}) \leq 1$.
\end{definition}

\begin{remark}[Trade-off Fundamental Rigidez-Adaptabilidad]\label{rem:tradeoff_rigidez_adaptabilidad}
Para arquitecturas de Geometric Deep Learning existe un trade-off fundamental:
\begin{equation}
    \text{Garantías Teóricas} \propto \text{Rigidez}(\mathcal{A})
\end{equation}
\begin{equation}
    \text{Capacidad Expresiva} \propto \text{Adaptabilidad}(\mathcal{A})
\end{equation}
Las arquitecturas óptimas maximizan $\text{Rigidez} \times \text{Adaptabilidad}$ sujeto a restricciones del problema.
\end{remark}

\begin{example}[Análisis del Espectro]\label{ex:analisis_espectro}
Consideremos la progresión:
\begin{align}
    \text{CNN Clásica}: &\quad \text{Rigidez} = 0.9, \text{Adaptabilidad} = 0.1 \\
    \text{Group CNN}: &\quad \text{Rigidez} = 0.7, \text{Adaptabilidad} = 0.3 \\
    \text{Graph Attention}: &\quad \text{Rigidez} = 0.3, \text{Adaptabilidad} = 0.7 \\
    \text{Full Transformer}: &\quad \text{Rigidez} = 0.1, \text{Adaptabilidad} = 0.9
\end{align}
Esto explica por qué CNNs proporcionan mejores garantías de generalización con menos datos, mientras que Transformers requieren datos masivos pero alcanzan mayor expresividad.
\end{example}

\subsection{Convergencia de Paradigmas: Análisis Matemático}\label{subsec:convergencia_paradigmas}

Los desarrollos recientes (2023-2024) sugieren una convergencia conceptual entre paradigmas convolucionales y de atención. Esta convergencia puede formalizarse matemáticamente:

\begin{conjecture}[Convergencia CNN-Attention]\label{conj:convergencia_cnn_attention}
Para secuencias de arquitecturas $\{\mathcal{C}_n\}$ (CNNs con complejidad creciente) y $\{\mathcal{A}_n\}$ (Attention con estructura creciente), existe límite común:
\begin{equation}
    \lim_{n \to \infty} d(\mathcal{C}_n, \mathcal{A}_n) = 0
\end{equation}
donde $d$ es una métrica en el espacio de funcionales sobre $\mathcal{F}(\mathcal{D})$, y la convergencia ocurre en la topología de operadores compactos.
\end{conjecture}

\textbf{Evidencia heurística.} La conjetura se sustenta en tres observaciones:

\textbf{1. Aproximación espectral}: Ambas familias de operadores pueden aproximarse arbitrariamente bien por operadores espectrales finito-dimensionales.

\textbf{2. Densidad en espacios funcionales}: El conjunto de convoluciones geométricas es denso en el espacio de operadores equivariantes, y los mecanismos de atención con restricciones geométricas convergen al mismo conjunto límite.

\textbf{3. Regularización adaptativa}: Las técnicas de regularización modernas (dropout geométrico, normalización adaptativa) hacen que ambas familias converjan hacia el mismo régimen de sesgo-varianza óptimo.

\subsubsection{Arquitecturas de Convergencia}

Las arquitecturas que ejemplifican esta convergencia incluyen:

\begin{definition}[ConvNext-Attention Híbrido]\label{def:convnext_attention}
Una arquitectura híbrida que combina:
\begin{itemize}
    \item \textbf{Núcleo convolucional}: Para procesar características locales con garantías geométricas
    \item \textbf{Cabezas de atención}: Para capturar dependencias globales adaptativas  
    \item \textbf{Intercambio dinámico}: Mecanismo aprendible que decide cuándo usar cada componente
\end{itemize}
\end{definition}

\begin{definition}[Retentive Networks Geométricas]\label{def:retentive_networks_geometricas}
Extensión del mecanismo retentivo a dominios geométricos:
\begin{equation}
    \text{Retention}(Q, K, V) = Q \odot \sum_{j} \text{decay}(d_\mathcal{M}(i,j)) \cdot K_j \odot V_j
\end{equation}
donde $d_\mathcal{M}$ es la distancia geodésica en el dominio geométrico $\mathcal{M}$.
\end{definition}

\subsection{Desarrollos Teóricos Futuros}\label{subsec:desarrollos_teoricos_futuros}

\subsubsection{Límites de Expresividad Unificados}

\begin{conjecture}[Límites Universales de Aproximación Geométrica]\label{conj:limites_universales}
Para toda función $f: \mathcal{M} \to \mathbb{R}$ continua sobre variedad compacta $\mathcal{M}$, y todo $\epsilon > 0$, existen arquitecturas CNN y Attention de complejidades $C_{CNN}(\epsilon)$ y $C_{Att}(\epsilon)$ tales que:
\begin{equation}
    \frac{C_{Att}(\epsilon)}{C_{CNN}(\epsilon)} \to 1 \quad \text{cuando } \epsilon \to 0
\end{equation}
Esta conjetura sugiere que la diferencia en complejidad arquitectónica desaparece en el límite de aproximación perfecta.
\end{conjecture}

\subsubsection{Teoría de Información Geométrica}

\begin{definition}[Entropía Geométrica de Fisher]\label{def:entropia_geometrica_fisher}
Para una distribución $p$ sobre variedad riemanniana $(\mathcal{M}, g)$, la entropía geométrica se define:
\begin{equation}
    H_g(p) = -\int_\mathcal{M} p(x) \log p(x) \sqrt{\det g(x)} \, dx
\end{equation}
Esta cantidad captura tanto la incertidumbre estadística como la complejidad geométrica subyacente.
\end{definition}

\begin{conjecture}[Principio de Información Mínima Geométrica]\label{conj:principio_informacion_minima}
Las arquitecturas óptimas de GDL minimizan la información geométrica necesaria para representar la función objetivo:
\begin{equation}
    \mathcal{A}^* = \arg\min_\mathcal{A} \left[ \mathcal{L}(\mathcal{A}) + \lambda H_g(\text{representaciones}_\mathcal{A}) \right]
\end{equation}
donde $\mathcal{L}$ es la pérdida empírica y $\lambda$ controla el trade-off información-precisión.
\end{conjecture}

\subsubsection{Meta-Aprendizaje Geométrico}

\begin{definition}[Meta-Learner Geométrico]\label{def:meta_learner_geometrico}
Un sistema que aprende automáticamente las simetrías relevantes de un dominio:
\begin{equation}
    \text{Meta-GDL}: \mathcal{D} \mapsto G_{\text{óptimo}} \mapsto \mathcal{A}_{G_{\text{óptimo}}}
\end{equation}
donde $\mathcal{D}$ son los datos, $G_{\text{óptimo}}$ es el grupo de simetrías descubierto, y $\mathcal{A}_{G_{\text{óptimo}}}$ es la arquitectura resultante.
\end{definition}

\subsection{Reflexión Final: Hacia una Teoría Unificada}\label{subsec:reflexion_final}

Los análisis presentados en este capítulo y en el \autoref{ch:cnn} convergen hacia una conclusión fundamental: la distinción entre "convolucional" y "atención" es cada vez menos relevante desde una perspectiva teórica profunda. Ambos paradigmas representan diferentes puntos en un continuum de operadores geométricos que procesan información estructurada.

El futuro del Geometric Deep Learning no reside en la competencia entre paradigmas, sino en su síntesis inteligente. Las arquitecturas más poderosas combinarán la estabilidad y eficiencia de las convoluciones geométricas con la flexibilidad y expresividad de la atención adaptativa, guiadas por principios matemáticos profundos en lugar de heurísticas ad hoc.

\begin{remark}[Principio de Síntesis Arquitectónica]\label{rem:principio_sintesis}
Para todo problema geométrico con estructura conocida $\mathcal{S}$ y adaptabilidad requerida $\mathcal{A}$, existe una arquitectura híbrida óptima:
\begin{equation}
    \mathcal{H}^* = \arg\min_{\mathcal{H}} \left[ \mathcal{L}(\mathcal{H}) + \lambda_1 d(\mathcal{H}, \mathcal{S}) + \lambda_2 (1 - \text{Adaptabilidad}(\mathcal{H})) \right]
\end{equation}
donde $d(\mathcal{H}, \mathcal{S})$ mide la distancia a la estructura conocida y los parámetros $\lambda_i$ balancean las restricciones.
\end{remark}

Esta síntesis representa más que una convergencia técnica: simboliza la maduración del campo hacia principios unificadores que trascienden divisiones históricas. Como demostraremos en el próximo capítulo experimental, esta unificación teórica se traduce en ventajas prácticas medibles y abre nuevas posibilidades para abordar los grandes desafíos científicos y tecnológicos del siglo XXI.

La geometría, como ha sido el caso a lo largo de la historia de las matemáticas y la física, nuevamente proporciona el lenguaje unificador necesario para comprender y explotar la complejidad aparente de los sistemas de aprendizaje automático modernos. En esta síntesis geométrica reside no solo la comprensión del presente, sino las herramientas conceptuales para construir el futuro del aprendizaje profundo.
